树、图和图论结构深度题卡 按题型拆成章内大节，但每个大节都先从 LeetCode 案例切入，再进入分流、案例矩阵、案例推导和实验复盘。这样读者看到的是从具体题推到规律，而不是先读抽象方法再猜它能用在哪。

\topic{二叉树与树形状态}
这一节先看 LeetCode 94 二叉树中序遍历、LeetCode 98 验证二叉搜索树、LeetCode 102 二叉树层序遍历 这几道 LeetCode 题，再整理同一题型的分流和推导。阅读顺序不是先背原理，而是先看案例的暴力瓶颈，再把重复出现的观察抽象成方法。

\subsection{原理和适用条件}

以 LeetCode 94 二叉树中序遍历、LeetCode 98 验证二叉搜索树、LeetCode 102 二叉树层序遍历 为主线：LeetCode 94 二叉树中序遍历：模型是“中序访问顺序是左子树、根、右子树，BST 中会得到有序序列。”，优化抓手是“用显式栈模拟一路向左，再回退访问根并转向右子树。”；LeetCode 98 验证二叉搜索树：模型是“BST 约束是全局上下界，不是只比较父子节点。”，优化抓手是“中序严格递增或显式传递上下界都可以；中序版本要保存前一个访问值。”；LeetCode 102 二叉树层序遍历：模型是“队列在每轮开始时保存当前层所有节点。”，优化抓手是“记录 level size 固定层边界，避免下一层节点混入当前层。”。这些案例共同推出的规律是：先想每个节点重复扫描子树会得到什么信息，再把这些信息压成节点向父节点返回的状态。树题要区分自顶向下约束、后序汇总和按层传播。 方法名只是复盘后的标签，真正要掌握的是从题面约束、暴力失败、状态设计到边界反例的推导链。

\begin{principlebox}{图示：从 LeetCode 案例到 二叉树与树形状态推导链}
\begin{tabular}{@{}p{0.18\linewidth}p{0.74\linewidth}@{}}
暴力基线 & 枚举候选、完整模拟或逐个检查，先保证覆盖全部可能答案。\\
瓶颈定位 & 慢点在于同一棵子树被重复求高度、最大值、是否包含目标等信息。若每个节点能一次性向父节点汇报足够信息，就不应重复下钻。\\
结构选择 & 树形 DP 把每个节点看成子问题根；BFS 把层次边界显式放进队列；BST 题利用中序有序或上下界约束。工程实现优先考虑显式栈避免深树调用栈风险。\\
不变量 & 返回值或栈元素必须有语义：当前子树高度、当前路径和、当前节点合法上下界、父指针、层号。不要只说遍历。\\
代码落地 & 示例中可以用指针节点，但比较节点身份时比较指针。极端深树用 \code{std::vector<TreeNode*>} 显式栈或队列。最大路径类题要用足够小的初值，不能默认零。\\
\end{tabular}
\end{principlebox}

\subsection{LeetCode 案例分流}

\begin{principlebox}{从 LeetCode 案例分流：二叉树与树形状态}
\textbf{遍历顺序。}
代表案例：LeetCode 94 中序遍历、LeetCode 102 层序遍历、LeetCode 103 锯齿层序、LeetCode 199 右视图。先看这些题的题面条件：题目要求按前序、中序、后序或层序收集节点。由此推出的处理方式是：用显式栈或队列保存节点和必要上下文。风险点：极端深树递归风险和层边界要单独考虑。

\textbf{BST 有序性。}
代表案例：LeetCode 98 验证 BST、LeetCode 230 第 K 小、LeetCode 235 BST 最近公共祖先、LeetCode 450 删除 BST 节点。先看这些题的题面条件：题目给定二叉搜索树或需要维护有序结构。由此推出的处理方式是：利用中序递增、上下界或前驱后继排除候选。风险点：普通二叉树不能套 BST 的值域剪枝。

\textbf{后序状态汇总。}
代表案例：LeetCode 110 平衡二叉树、LeetCode 124 最大路径和、LeetCode 337 打家劫舍 III、LeetCode 543 二叉树直径。先看这些题的题面条件：父节点答案依赖左右子树返回的信息。由此推出的处理方式是：子树返回高度、贡献、状态对或合法性，父节点合并。风险点：返回给父亲的值和全局答案不是同一个量。

\textbf{构造和变形。}
代表案例：LeetCode 105 前中构造、LeetCode 106 中后构造、LeetCode 108 有序数组转 BST、LeetCode 114 展开为链表。先看这些题的题面条件：题目要求根据遍历序列构造树，或原地改变树形。由此推出的处理方式是：用哈希定位根在中序中的位置，或用栈按目标顺序重连。风险点：节点值唯一性和区间边界是主要条件。

\textbf{Trie 和前缀树。}
代表案例：LeetCode 208 实现 Trie、LeetCode 211 添加与搜索单词、LeetCode 212 单词搜索 II。先看这些题的题面条件：字符串集合需要前缀共享或通配搜索。由此推出的处理方式是：每个字符沿边前进，节点保存子边和单词结束标记。风险点：前缀存在不等于单词存在，通配符会扩大搜索分支。

\end{principlebox}

\subsection{案例矩阵}

本节案例覆盖 LeetCode 94 二叉树中序遍历、LeetCode 98 验证二叉搜索树、LeetCode 102 二叉树层序遍历、LeetCode 103 二叉树锯齿形层序遍历、LeetCode 104 二叉树最大深度、LeetCode 105 前序中序构造二叉树、LeetCode 106 中序后序构造二叉树、LeetCode 108 有序数组转二叉搜索树、LeetCode 110 平衡二叉树、LeetCode 111 二叉树最小深度、LeetCode 112 路径总和、LeetCode 113 路径总和 II、LeetCode 114 二叉树展开为链表、LeetCode 124 二叉树最大路径和、LeetCode 199 二叉树右视图、LeetCode 208 实现 Trie、LeetCode 211 添加与搜索单词、LeetCode 226 翻转二叉树、LeetCode 230 二叉搜索树中第 K 小、LeetCode 235 二叉搜索树最近公共祖先、LeetCode 236 二叉树最近公共祖先、LeetCode 337 打家劫舍 III、LeetCode 450 删除二叉搜索树中的节点、LeetCode 543 二叉树直径、LeetCode 834 树中距离之和、LeetCode 1483 树节点的第 K 个祖先。阅读时不要按题号背答案，而要先判断题面落在哪个分流场景：查询目标是什么，输入约束给了什么单调性或等价关系，暴力慢在哪里，优化结构保存了什么状态。

\begin{principlebox}{案例矩阵}
\textbf{LeetCode 94 二叉树中序遍历。}
中序访问顺序是左子树、根、右子树，BST 中会得到有序序列。单点风险：空树、单节点、链式左树、链式右树。

\textbf{LeetCode 98 验证二叉搜索树。}
BST 约束是全局上下界，不是只比较父子节点。单点风险：重复值、int 最小最大值、局部合法但全局非法。

\textbf{LeetCode 102 二叉树层序遍历。}
队列在每轮开始时保存当前层所有节点。单点风险：空树、只有根、最后一层不满、左右顺序。

\textbf{LeetCode 103 二叉树锯齿形层序遍历。}
BFS 层边界不变，变化的是当前层写入结果的方向。单点风险：单层、偶数层、奇数层、空树。

\textbf{LeetCode 104 二叉树最大深度。}
深度可以看成层数，也可以作为栈中携带的状态。单点风险：空树返回零、根深度是一、链式树。

\textbf{LeetCode 105 前序中序构造二叉树。}
前序给根，中序把左右子树切开。单点风险：节点值唯一、空子树、左右子树长度计算、输入不合法。

\textbf{LeetCode 106 中序后序构造二叉树。}
后序末尾是根，中序位置确定左右子树规模。单点风险：空树、单节点、全左链、全右链。

\textbf{LeetCode 108 有序数组转二叉搜索树。}
有序数组中点作为根能让左右规模接近。单点风险：空区间、偶数长度选左中还是右中、高度平衡定义。

\textbf{LeetCode 110 平衡二叉树。}
每个节点需要知道子树高度以及是否已经失衡。单点风险：空树、单节点、左右高度差正好一、链式树。

\textbf{LeetCode 111 二叉树最小深度。}
最小深度是到最近叶子的节点数，不是左右深度简单取最小。单点风险：只有左子树、只有右子树、根为空、根是叶子。

\textbf{LeetCode 112 路径总和。}
路径必须从根到叶子，状态是剩余目标值。单点风险：负数节点、目标为零、非叶子提前达到目标、空树。

\textbf{LeetCode 113 路径总和 II。}
相比判定题，状态还要保存当前路径。单点风险：多个答案、负数、空树、路径数组复制时机。

\textbf{LeetCode 114 二叉树展开为链表。}
目标是按前序顺序把树原地改成右指针链。单点风险：空树、单节点、原有右子树不能丢、左指针置空。

\textbf{LeetCode 124 二叉树最大路径和。}
路径可以从任意节点到任意节点，但不能分叉向父节点贡献两边。单点风险：全负数、单节点、左右增益为负时取零、答案初始值。

\textbf{LeetCode 199 二叉树右视图。}
每层从右侧能看到的就是该层最后访问或最先访问的右侧节点。单点风险：空树、左链、右链、左右交错。

\textbf{LeetCode 208 实现 Trie。}
前缀树把字符串公共前缀共享成路径。单点风险：空字符串约定、单词结束标记、前缀存在不等于单词存在。

\textbf{LeetCode 211 添加与搜索单词。}
点号通配符让查询可能分叉，但前缀树仍能剪掉不匹配前缀。单点风险：连续通配符、空结果、单词结束标记、字符集大小。

\textbf{LeetCode 226 翻转二叉树。}
每个节点交换左右孩子即可，遍历顺序不影响最终结构。单点风险：空树、单节点、只有一侧子树、重复值。

\textbf{LeetCode 230 二叉搜索树中第 K 小。}
BST 中序遍历按升序访问节点。单点风险：k 为一、k 为节点数、树不平衡、是否有重复值。

\textbf{LeetCode 235 二叉搜索树最近公共祖先。}
BST 的值域关系能决定两个目标在当前节点的哪一侧。单点风险：一个节点是另一个祖先、目标顺序、重复值约定。

\textbf{LeetCode 236 二叉树最近公共祖先。}
普通树没有有序性，只能依靠祖先关系。单点风险：目标不存在、p 等于 q、一个是另一个祖先、比较节点地址。

\textbf{LeetCode 337 打家劫舍 III。}
树上相邻节点不能同时选，每个节点需要偷和不偷两个状态。单点风险：空树、负值是否可能、链式树、递归深度。

\textbf{LeetCode 450 删除二叉搜索树中的节点。}
删除节点后仍要保持 BST 中序有序。单点风险：删除叶子、一个孩子、两个孩子、删除根节点。

\textbf{LeetCode 543 二叉树直径。}
直径经过某节点时等于左高加右高，但向父亲只能贡献单边高度。单点风险：空树、单节点、链式树、边数还是节点数定义。

\textbf{LeetCode 834 树中距离之和。}
每个节点作为根的距离和可以由相邻换根关系转移。单点风险：单节点、链式树、星形树、n 很大递归栈、边是无向需建双向邻接。

\textbf{LeetCode 1483 树节点的第 K 个祖先。}
第 k 个祖先可以按 k 的二进制拆成若干 \ensuremath{2^{j}} 级跳跃。单点风险：根节点祖先不存在、k 为零、树很深、查询很多、表大小。

\end{principlebox}

\subsection{共性落点}

\begin{principlebox}{本组共性落点}
\textbf{慢版本。} 暴力做法对每个节点反复扫描子树，或者只检查父子局部关系。树题真正考的是状态如何从父到子、从子到父或按层传播。
\textbf{瓶颈。} 慢点在于同一棵子树被重复求高度、最大值、是否包含目标等信息。若每个节点能一次性向父节点汇报足够信息，就不应重复下钻。
\textbf{状态字段。} 节点指针、传入约束或返回值、当前答案、层号或路径状态；空节点的返回值必须和状态定义匹配。
\textbf{不变量。} 返回值或栈元素必须有语义：当前子树高度、当前路径和、当前节点合法上下界、父指针、层号。不要只说遍历。
\textbf{实现检查。} 递归版本先处理空节点；迭代版本的栈元素要带完整状态；全负值、重复值和极端深树要单独测。
\textbf{C++ 落地。} 示例中可以用指针节点，但比较节点身份时比较指针。极端深树用 \code{std::vector<TreeNode*>} 显式栈或队列。最大路径类题要用足够小的初值，不能默认零。
\textbf{证明抓手。} 证明从子树归纳或层序不变量开始：孩子状态正确时父节点合并正确；若用队列，当前轮队列必须正好保存当前层节点。
\end{principlebox}

\subsection{案例推导}

\noindent\textbf{LeetCode 94 二叉树中序遍历。}

\textbf{题目说明。} LeetCode 94 二叉树中序遍历 的题面入口要先还原成具体任务：中序访问顺序是左子树、根、右子树，BST 中会得到有序序列。

\textbf{样例入口。} 拿根 2、左 1、右 3 手算，中序结果必须是 1, 2, 3。显式栈做法先一路压到最左，弹出 1 后访问，再回到 2。

\textbf{暴力基线。} 慢版本在每个节点重新扫描子树或路径，先求出局部答案。重复扫描暴露了真正状态：每个节点只需要从孩子或父亲那里拿到少量信息。放到本题就是：先按“中序访问顺序是左子树、根、右子树，BST 中会得到有序序列”完整试慢版本；样例里已经能看到“规律是栈保存尚未访问但左侧已经处理完的祖先，BST 题里这个顺序会产生递增序列”，所以优化前必须先能说清慢版本枚举了哪些候选。

\textbf{暴力过程。} 拿根 2、左 1、右 3 手算，中序结果必须是 1, 2, 3。显式栈做法先一路压到最左，弹出 1 后访问，再回到 2。

\textbf{重复工作。} 题面模型：中序访问顺序是左子树、根、右子树，BST 中会得到有序序列。暴力版的慢点要落到本题：慢版本会在树上重复确认“中序访问顺序是左子树、根、右子树，BST 中会得到有序序列”。样例规律是“规律是栈保存尚未访问但左侧已经处理完的祖先，BST 题里这个顺序会产生递增序列”，说明父子之间真正需要传递的是高度、上下界、命中状态、层号或两类选择值，而不是反复扫描整棵子树。后面的优化只消掉这类重复工作，不能改变答案集合。

\textbf{优化条件。} 本题的关键观察是：用显式栈模拟一路向左，再回退访问根并转向右子树。这句话必须能翻译成可维护状态；如果题目条件改变后观察不再成立，就不能继续套原来的写法。

\textbf{相邻状态。} 规律是栈保存尚未访问但左侧已经处理完的祖先，BST 题里这个顺序会产生递增序列。把这个特例继续拆开看：节点向父亲返回的量和全局答案通常不是同一个量。先写叶子或空节点的返回值，再写父节点如何合并左右孩子，递归式才有边界基础。

\textbf{状态复用。} 把对子树的重复扫描压成返回值或队列状态；每个节点只合并孩子给出的稳定信息。本题落点是：规律是栈保存尚未访问但左侧已经处理完的祖先，BST 题里这个顺序会产生递增序列。手算时只改被新输入影响的那一小部分，而不是回头重算完整候选。

\textbf{指针和字段。} 状态字段要能说清：节点指针、传入约束或返回值、当前答案、层号或路径状态；空节点的返回值必须和状态定义匹配。手算时按这个协议跟踪：拿根 2、左 1、右 3 手算，中序结果必须是 1, 2, 3。显式栈做法先一路压到最左，弹出 1 后访问，再回到 2。然后按协议复查：手算时从叶子开始写返回值，或者从根开始写传入约束。不要只写访问顺序，要写每条边上传递的信息。

\textbf{代码落点。} 递归版本先处理空节点；迭代版本的栈元素要带完整状态；全负值、重复值和极端深树要单独测。关键代码行必须能逐句翻译成“旧状态怎样变成新状态”，否则就只是模板。

\textbf{正确性。} 证明从子树归纳或层序不变量开始：孩子状态正确时父节点合并正确；若用队列，当前轮队列必须正好保存当前层节点。

\textbf{边界实验。} 先用空树、单节点、链式左树、链式右树做反例检查。实验时覆盖空树、单节点、链式左树、链式右树，尤其关注空节点、单节点、链式树和全负值。递归版本和显式栈版本可以互相对拍。

\textbf{复杂度变化。} 暴力在每个节点重复扫描子树或反复寻找层边界可能退化到平方级；后序汇总、显式栈或层序遍历让每个节点处理常数次，通常是线性时间。

\textbf{迁移追问。} 如果题目改成“Morris 遍历可做到常数额外空间，但会临时改树结构”，先判断原状态是否仍然覆盖新答案集合，再决定是微调边界、改 value 语义，还是换算法。

\noindent\textbf{LeetCode 98 验证二叉搜索树。}

\textbf{题目说明。} LeetCode 98 验证二叉搜索树 的题面入口要先还原成具体任务：BST 约束是全局上下界，不是只比较父子节点。

\textbf{样例入口。} 拿 [5, 1, 4, null, null, 3, 6] 手算，节点 3 小于根 5 却在右子树里，只比较父子 4 和 3 会误判。

\textbf{暴力基线。} 错误暴力常只检查当前节点大于左孩子、小于右孩子；正确慢版本要对每个节点扫描整棵左子树最大值和右子树最小值。

\textbf{暴力过程。} 以 [5, 1, 4, null, null, 3, 6] 为例，节点 4 的左孩子 3 小于 4、右孩子 6 大于 4，看局部没问题；但 3 在根 5 的右子树里，小于 5，所以整棵树不是 BST。

\textbf{重复工作。} 题面模型：BST 约束是全局上下界，不是只比较父子节点。暴力版的慢点要落到本题：浪费在每个节点重复扫描子树求上下界。真正需要传递的是祖先给当前子树限定的取值范围。后面的优化只消掉这类重复工作，不能改变答案集合。

\textbf{优化条件。} 本题的关键观察是：中序严格递增或显式传递上下界都可以；中序版本要保存前一个访问值。这句话必须能翻译成可维护状态；如果题目条件改变后观察不再成立，就不能继续套原来的写法。

\textbf{相邻状态。} 规律是 BST 约束来自整条祖先上下界，或中序序列必须严格递增；重复值和 int 边界要按题意单独处理。把这个特例继续拆开看：节点向父亲返回的量和全局答案通常不是同一个量。先写叶子或空节点的返回值，再写父节点如何合并左右孩子，递归式才有边界基础。

\textbf{状态复用。} 自顶向下传入 lower 和 upper，当前节点值必须满足 lower<val<upper；去左子树时 upper 变成当前值，去右子树时 lower 变成当前值。手算时只改被新输入影响的那一小部分，而不是回头重算完整候选。

\textbf{指针和字段。} 状态字段要能说清：node 是当前节点，lower/upper 是祖先约束；空节点返回合法。手算时按这个协议跟踪：到样例中右子树节点 3 时，它的 lower 应该是 5、upper 是 4 或来自路径约束，3 不满足根给的下界。

\textbf{代码落点。} 关键是用 long long 或 optional 边界，避免节点值等于 int 最小/最大时哨兵冲突；重复值是否允许按题意处理，本题严格不允许。关键代码行必须能逐句翻译成“旧状态怎样变成新状态”，否则就只是模板。

\textbf{正确性。} BST 要求左子树所有值小于根、右子树所有值大于根；上下界把所有祖先约束都带到当前节点，覆盖了非局部违规。

\textbf{边界实验。} 先用重复值、int 最小最大值、局部合法但全局非法做反例检查。实验时覆盖重复值、int 最小最大值、局部合法但全局非法，尤其关注空节点、单节点、链式树和全负值。递归版本和显式栈版本可以互相对拍。

\textbf{复杂度变化。} 扫描子树慢版本最坏 O(\ensuremath{n^{2}})，上下界 DFS/BFS 每节点一次 O(n)，空间 O(h) 或显式栈 O(h)。

\textbf{迁移追问。} 如果题目改成“若允许重复值，需要明确重复放左还是放右”，先判断原状态是否仍然覆盖新答案集合，再决定是微调边界、改 value 语义，还是换算法。

\noindent\textbf{LeetCode 102 二叉树层序遍历。}

\textbf{题目说明。} LeetCode 102 二叉树层序遍历 的题面入口要先还原成具体任务：队列在每轮开始时保存当前层所有节点。

\textbf{样例入口。} 拿根 3、左 9、右 20 手算，队列第一轮只有 3，第二轮才是 9 和 20。若不固定本层 size，新入队的下一层节点会混进当前层。

\textbf{暴力基线。} 慢版本可以先给每个节点计算深度，再按深度分组；或者每层都从根重新扫描一遍收集该层节点。

\textbf{暴力过程。} 以根 3、左右孩子 9 和 20 为例，第一层只有 3，第二层是 9、20。若队列不固定当前层长度，20 的孩子可能混进同一层。

\textbf{重复工作。} 题面模型：队列在每轮开始时保存当前层所有节点。暴力版的慢点要落到本题：浪费在重复从根找层，或者层边界不清导致混层。队列天然按访问顺序保存下一批节点，只要固定本轮队列大小就是当前层。后面的优化只消掉这类重复工作，不能改变答案集合。

\textbf{优化条件。} 本题的关键观察是：记录 level size 固定层边界，避免下一层节点混入当前层。这句话必须能翻译成可维护状态；如果题目条件改变后观察不再成立，就不能继续套原来的写法。

\textbf{相邻状态。} 规律是每轮开始记录队列长度，这个长度就是当前层边界。把这个特例继续拆开看：节点向父亲返回的量和全局答案通常不是同一个量。先写叶子或空节点的返回值，再写父节点如何合并左右孩子，递归式才有边界基础。

\textbf{状态复用。} BFS 队列保存待处理节点。每轮先读取 level\_size，只弹出这 level\_size 个节点并收集值，同时把它们的孩子加入队尾。手算时只改被新输入影响的那一小部分，而不是回头重算完整候选。

\textbf{指针和字段。} 状态字段要能说清：queue 是待访问节点队列，level\_size 是当前层节点数，level\_values 是本层输出。手算时按这个协议跟踪：队列初始 [3]，level\_size=1，弹出 3 后压入 9 和 20；下一轮 level\_size=2，只处理 9 和 20。

\textbf{代码落点。} 关键是 for 循环次数使用进入本层前的 queue.size()，不要在循环条件里直接用动态变化的 queue.size()。关键代码行必须能逐句翻译成“旧状态怎样变成新状态”，否则就只是模板。

\textbf{正确性。} BFS 入队顺序保证父层节点先于子层节点；固定 level\_size 后，本轮处理的恰好是同一深度的全部节点。

\textbf{边界实验。} 先用空树、只有根、最后一层不满、左右顺序做反例检查。实验时覆盖空树、只有根、最后一层不满、左右顺序，尤其关注空节点、单节点、链式树和全负值。递归版本和显式栈版本可以互相对拍。

\textbf{复杂度变化。} 重复扫描每层最坏 O(nh)，BFS 每节点入队出队一次 O(n)，队列空间 O(width)。

\textbf{迁移追问。} 如果题目改成“锯齿形层序只是在每层输出顺序上增加方向控制”，先判断原状态是否仍然覆盖新答案集合，再决定是微调边界、改 value 语义，还是换算法。

\noindent\textbf{LeetCode 103 二叉树锯齿形层序遍历。}

\textbf{题目说明。} LeetCode 103 二叉树锯齿形层序遍历 的题面入口要先还原成具体任务：BFS 层边界不变，变化的是当前层写入结果的方向。

\textbf{样例入口。} 拿三层树手算，第一层从左到右，第二层从右到左，第三层再从左到右；BFS 层边界不变，只是写入结果的位置反向。

\textbf{暴力基线。} 慢版本在每个节点重新扫描子树或路径，先求出局部答案。重复扫描暴露了真正状态：每个节点只需要从孩子或父亲那里拿到少量信息。放到本题就是：先按“BFS 层边界不变，变化的是当前层写入结果的方向”完整试慢版本；样例里已经能看到“规律是层序遍历负责分层，方向标志只影响当前层数组的填充顺序”，所以优化前必须先能说清慢版本枚举了哪些候选。

\textbf{暴力过程。} 拿三层树手算，第一层从左到右，第二层从右到左，第三层再从左到右；BFS 层边界不变，只是写入结果的位置反向。

\textbf{重复工作。} 题面模型：BFS 层边界不变，变化的是当前层写入结果的方向。暴力版的慢点要落到本题：慢版本会反复确认“BFS 层边界不变，变化的是当前层写入结果的方向”里的层边界。样例规律是“规律是层序遍历负责分层，方向标志只影响当前层数组的填充顺序”，说明队列当前轮的长度就是稳定边界；不先固定边界，下一层节点会混入当前层。后面的优化只消掉这类重复工作，不能改变答案集合。

\textbf{优化条件。} 本题的关键观察是：可以层内反转，也可以预分配数组按方向写入目标下标。这句话必须能翻译成可维护状态；如果题目条件改变后观察不再成立，就不能继续套原来的写法。

\textbf{相邻状态。} 规律是层序遍历负责分层，方向标志只影响当前层数组的填充顺序。把这个特例继续拆开看：节点向父亲返回的量和全局答案通常不是同一个量。先写叶子或空节点的返回值，再写父节点如何合并左右孩子，递归式才有边界基础。

\textbf{状态复用。} 把对子树的重复扫描压成返回值或队列状态；每个节点只合并孩子给出的稳定信息。本题落点是：规律是层序遍历负责分层，方向标志只影响当前层数组的填充顺序。手算时只改被新输入影响的那一小部分，而不是回头重算完整候选。

\textbf{指针和字段。} 状态字段要能说清：节点指针、传入约束或返回值、当前答案、层号或路径状态；空节点的返回值必须和状态定义匹配。手算时按这个协议跟踪：拿三层树手算，第一层从左到右，第二层从右到左，第三层再从左到右；BFS 层边界不变，只是写入结果的位置反向。然后按协议复查：手算时从叶子开始写返回值，或者从根开始写传入约束。不要只写访问顺序，要写每条边上传递的信息。

\textbf{代码落点。} 递归版本先处理空节点；迭代版本的栈元素要带完整状态；全负值、重复值和极端深树要单独测。关键代码行必须能逐句翻译成“旧状态怎样变成新状态”，否则就只是模板。

\textbf{正确性。} 证明从子树归纳或层序不变量开始：孩子状态正确时父节点合并正确；若用队列，当前轮队列必须正好保存当前层节点。

\textbf{边界实验。} 先用单层、偶数层、奇数层、空树做反例检查。实验时覆盖单层、偶数层、奇数层、空树，尤其关注空节点、单节点、链式树和全负值。递归版本和显式栈版本可以互相对拍。

\textbf{复杂度变化。} 暴力在每个节点重复扫描子树或反复寻找层边界可能退化到平方级；后序汇总、显式栈或层序遍历让每个节点处理常数次，通常是线性时间。

\textbf{迁移追问。} 如果题目改成“若要求从底向上输出，收集后反转层列表即可”，先判断原状态是否仍然覆盖新答案集合，再决定是微调边界、改 value 语义，还是换算法。

\noindent\textbf{LeetCode 104 二叉树最大深度。}

\textbf{题目说明。} LeetCode 104 二叉树最大深度 的题面入口要先还原成具体任务：深度可以看成层数，也可以作为栈中携带的状态。

\textbf{样例入口。} 拿链式树 1->2->3 手算，最大深度是 3；空树是 0。递归返回左右子树深度的较大值加一，显式栈则把节点和当前深度一起压栈。

\textbf{暴力基线。} 慢版本在每个节点重新扫描子树或路径，先求出局部答案。重复扫描暴露了真正状态：每个节点只需要从孩子或父亲那里拿到少量信息。放到本题就是：先按“深度可以看成层数，也可以作为栈中携带的状态”完整试慢版本；样例里已经能看到“规律是状态必须携带深度，不能只记录访问过哪些节点”，所以优化前必须先能说清慢版本枚举了哪些候选。

\textbf{暴力过程。} 拿链式树 1->2->3 手算，最大深度是 3；空树是 0。递归返回左右子树深度的较大值加一，显式栈则把节点和当前深度一起压栈。

\textbf{重复工作。} 题面模型：深度可以看成层数，也可以作为栈中携带的状态。暴力版的慢点要落到本题：慢版本会在树上重复确认“深度可以看成层数，也可以作为栈中携带的状态”。样例规律是“规律是状态必须携带深度，不能只记录访问过哪些节点”，说明父子之间真正需要传递的是高度、上下界、命中状态、层号或两类选择值，而不是反复扫描整棵子树。后面的优化只消掉这类重复工作，不能改变答案集合。

\textbf{优化条件。} 本题的关键观察是：显式栈保存节点和深度，避免极端链式树递归栈过深。这句话必须能翻译成可维护状态；如果题目条件改变后观察不再成立，就不能继续套原来的写法。

\textbf{相邻状态。} 规律是状态必须携带深度，不能只记录访问过哪些节点。把这个特例继续拆开看：节点向父亲返回的量和全局答案通常不是同一个量。先写叶子或空节点的返回值，再写父节点如何合并左右孩子，递归式才有边界基础。

\textbf{状态复用。} 把对子树的重复扫描压成返回值或队列状态；每个节点只合并孩子给出的稳定信息。本题落点是：规律是状态必须携带深度，不能只记录访问过哪些节点。手算时只改被新输入影响的那一小部分，而不是回头重算完整候选。

\textbf{指针和字段。} 状态字段要能说清：节点指针、传入约束或返回值、当前答案、层号或路径状态；空节点的返回值必须和状态定义匹配。手算时按这个协议跟踪：拿链式树 1->2->3 手算，最大深度是 3；空树是 0。递归返回左右子树深度的较大值加一，显式栈则把节点和当前深度一起压栈。然后按协议复查：手算时从叶子开始写返回值，或者从根开始写传入约束。不要只写访问顺序，要写每条边上传递的信息。

\textbf{代码落点。} 递归版本先处理空节点；迭代版本的栈元素要带完整状态；全负值、重复值和极端深树要单独测。关键代码行必须能逐句翻译成“旧状态怎样变成新状态”，否则就只是模板。

\textbf{正确性。} 证明从子树归纳或层序不变量开始：孩子状态正确时父节点合并正确；若用队列，当前轮队列必须正好保存当前层节点。

\textbf{边界实验。} 先用空树返回零、根深度是一、链式树做反例检查。实验时覆盖空树返回零、根深度是一、链式树，尤其关注空节点、单节点、链式树和全负值。递归版本和显式栈版本可以互相对拍。

\textbf{复杂度变化。} 暴力在每个节点重复扫描子树或反复寻找层边界可能退化到平方级；后序汇总、显式栈或层序遍历让每个节点处理常数次，通常是线性时间。

\textbf{迁移追问。} 如果题目改成“最小深度不能简单取左右最小，因为空子树不构成叶子路径”，先判断原状态是否仍然覆盖新答案集合，再决定是微调边界、改 value 语义，还是换算法。

\noindent\textbf{LeetCode 105 前序中序构造二叉树。}

\textbf{题目说明。} LeetCode 105 前序中序构造二叉树 的题面入口要先还原成具体任务：前序给根，中序把左右子树切开。

\textbf{样例入口。} 拿 preorder=[3, 9, 20, 15, 7]、inorder=[9, 3, 15, 20, 7] 手算，前序第一个 3 是根，中序中 3 左边是左子树，右边是右子树。

\textbf{暴力基线。} 慢版本在每个节点重新扫描子树或路径，先求出局部答案。重复扫描暴露了真正状态：每个节点只需要从孩子或父亲那里拿到少量信息。放到本题就是：先按“前序给根，中序把左右子树切开”完整试慢版本；样例里已经能看到“规律是用哈希表快速找根在中序的位置，再用左右子树长度切分前序区间”，所以优化前必须先能说清慢版本枚举了哪些候选。

\textbf{暴力过程。} 拿 preorder=[3, 9, 20, 15, 7]、inorder=[9, 3, 15, 20, 7] 手算，前序第一个 3 是根，中序中 3 左边是左子树，右边是右子树。

\textbf{重复工作。} 题面模型：前序给根，中序把左右子树切开。暴力版的慢点要落到本题：慢版本会在树上重复确认“前序给根，中序把左右子树切开”。样例规律是“规律是用哈希表快速找根在中序的位置，再用左右子树长度切分前序区间”，说明父子之间真正需要传递的是高度、上下界、命中状态、层号或两类选择值，而不是反复扫描整棵子树。后面的优化只消掉这类重复工作，不能改变答案集合。

\textbf{优化条件。} 本题的关键观察是：用哈希表记录中序下标，避免每次线性寻找根位置。这句话必须能翻译成可维护状态；如果题目条件改变后观察不再成立，就不能继续套原来的写法。

\textbf{相邻状态。} 规律是用哈希表快速找根在中序的位置，再用左右子树长度切分前序区间。把这个特例继续拆开看：节点向父亲返回的量和全局答案通常不是同一个量。先写叶子或空节点的返回值，再写父节点如何合并左右孩子，递归式才有边界基础。

\textbf{状态复用。} 把对子树的重复扫描压成返回值或队列状态；每个节点只合并孩子给出的稳定信息。本题落点是：规律是用哈希表快速找根在中序的位置，再用左右子树长度切分前序区间。手算时只改被新输入影响的那一小部分，而不是回头重算完整候选。

\textbf{指针和字段。} 状态字段要能说清：节点指针、传入约束或返回值、当前答案、层号或路径状态；空节点的返回值必须和状态定义匹配。手算时按这个协议跟踪：拿 preorder=[3, 9, 20, 15, 7]、inorder=[9, 3, 15, 20, 7] 手算，前序第一个 3 是根，中序中 3 左边是左子树，右边是右子树。然后按协议复查：手算时从叶子开始写返回值，或者从根开始写传入约束。不要只写访问顺序，要写每条边上传递的信息。

\textbf{代码落点。} 递归版本先处理空节点；迭代版本的栈元素要带完整状态；全负值、重复值和极端深树要单独测。关键代码行必须能逐句翻译成“旧状态怎样变成新状态”，否则就只是模板。

\textbf{正确性。} 证明从子树归纳或层序不变量开始：孩子状态正确时父节点合并正确；若用队列，当前轮队列必须正好保存当前层节点。

\textbf{边界实验。} 先用节点值唯一、空子树、左右子树长度计算、输入不合法做反例检查。实验时覆盖节点值唯一、空子树、左右子树长度计算、输入不合法，尤其关注空节点、单节点、链式树和全负值。递归版本和显式栈版本可以互相对拍。

\textbf{复杂度变化。} 暴力在每个节点重复扫描子树或反复寻找层边界可能退化到平方级；后序汇总、显式栈或层序遍历让每个节点处理常数次，通常是线性时间。

\textbf{迁移追问。} 如果题目改成“后序中序构造类似，只是根来自后序末尾”，先判断原状态是否仍然覆盖新答案集合，再决定是微调边界、改 value 语义，还是换算法。

\noindent\textbf{LeetCode 106 中序后序构造二叉树。}

\textbf{题目说明。} LeetCode 106 中序后序构造二叉树 的题面入口要先还原成具体任务：后序末尾是根，中序位置确定左右子树规模。

\textbf{样例入口。} 拿 inorder=[9, 3, 15, 20, 7]、postorder=[9, 15, 7, 20, 3] 手算，后序最后一个 3 是根，中序把左右子树切开。

\textbf{暴力基线。} 慢版本在每个节点重新扫描子树或路径，先求出局部答案。重复扫描暴露了真正状态：每个节点只需要从孩子或父亲那里拿到少量信息。放到本题就是：先按“后序末尾是根，中序位置确定左右子树规模”完整试慢版本；样例里已经能看到“规律是根来自后序末尾，左右区间长度必须同步更新；全左链和全右链最容易暴露边界错位”，所以优化前必须先能说清慢版本枚举了哪些候选。

\textbf{暴力过程。} 拿 inorder=[9, 3, 15, 20, 7]、postorder=[9, 15, 7, 20, 3] 手算，后序最后一个 3 是根，中序把左右子树切开。

\textbf{重复工作。} 题面模型：后序末尾是根，中序位置确定左右子树规模。暴力版的慢点要落到本题：慢版本会在树上重复确认“后序末尾是根，中序位置确定左右子树规模”。样例规律是“规律是根来自后序末尾，左右区间长度必须同步更新；全左链和全右链最容易暴露边界错位”，说明父子之间真正需要传递的是高度、上下界、命中状态、层号或两类选择值，而不是反复扫描整棵子树。后面的优化只消掉这类重复工作，不能改变答案集合。

\textbf{优化条件。} 本题的关键观察是：构造时要小心后序左右区间边界，避免切分错位。这句话必须能翻译成可维护状态；如果题目条件改变后观察不再成立，就不能继续套原来的写法。

\textbf{相邻状态。} 规律是根来自后序末尾，左右区间长度必须同步更新；全左链和全右链最容易暴露边界错位。把这个特例继续拆开看：节点向父亲返回的量和全局答案通常不是同一个量。先写叶子或空节点的返回值，再写父节点如何合并左右孩子，递归式才有边界基础。

\textbf{状态复用。} 把对子树的重复扫描压成返回值或队列状态；每个节点只合并孩子给出的稳定信息。本题落点是：规律是根来自后序末尾，左右区间长度必须同步更新；全左链和全右链最容易暴露边界错位。手算时只改被新输入影响的那一小部分，而不是回头重算完整候选。

\textbf{指针和字段。} 状态字段要能说清：节点指针、传入约束或返回值、当前答案、层号或路径状态；空节点的返回值必须和状态定义匹配。手算时按这个协议跟踪：拿 inorder=[9, 3, 15, 20, 7]、postorder=[9, 15, 7, 20, 3] 手算，后序最后一个 3 是根，中序把左右子树切开。然后按协议复查：手算时从叶子开始写返回值，或者从根开始写传入约束。不要只写访问顺序，要写每条边上传递的信息。

\textbf{代码落点。} 递归版本先处理空节点；迭代版本的栈元素要带完整状态；全负值、重复值和极端深树要单独测。关键代码行必须能逐句翻译成“旧状态怎样变成新状态”，否则就只是模板。

\textbf{正确性。} 证明从子树归纳或层序不变量开始：孩子状态正确时父节点合并正确；若用队列，当前轮队列必须正好保存当前层节点。

\textbf{边界实验。} 先用空树、单节点、全左链、全右链做反例检查。实验时覆盖空树、单节点、全左链、全右链，尤其关注空节点、单节点、链式树和全负值。递归版本和显式栈版本可以互相对拍。

\textbf{复杂度变化。} 暴力在每个节点重复扫描子树或反复寻找层边界可能退化到平方级；后序汇总、显式栈或层序遍历让每个节点处理常数次，通常是线性时间。

\textbf{迁移追问。} 如果题目改成“若给前序和后序，二叉树通常不能唯一确定，除非额外限制”，先判断原状态是否仍然覆盖新答案集合，再决定是微调边界、改 value 语义，还是换算法。

\noindent\textbf{LeetCode 108 有序数组转二叉搜索树。}

\textbf{题目说明。} LeetCode 108 有序数组转二叉搜索树 的题面入口要先还原成具体任务：有序数组中点作为根能让左右规模接近。

\textbf{样例入口。} 拿 [-10, -3, 0, 5, 9] 手算，选 0 做根，左边两个数构成左子树，右边两个数构成右子树，高度接近平衡。

\textbf{暴力基线。} 慢版本在每个节点重新扫描子树或路径，先求出局部答案。重复扫描暴露了真正状态：每个节点只需要从孩子或父亲那里拿到少量信息。放到本题就是：先按“有序数组中点作为根能让左右规模接近”完整试慢版本；样例里已经能看到“规律是有序数组的中点天然把值域分成 BST 左右两半，空区间返回空节点”，所以优化前必须先能说清慢版本枚举了哪些候选。

\textbf{暴力过程。} 拿 [-10, -3, 0, 5, 9] 手算，选 0 做根，左边两个数构成左子树，右边两个数构成右子树，高度接近平衡。

\textbf{重复工作。} 题面模型：有序数组中点作为根能让左右规模接近。暴力版的慢点要落到本题：慢版本会在树上重复确认“有序数组中点作为根能让左右规模接近”。样例规律是“规律是有序数组的中点天然把值域分成 BST 左右两半，空区间返回空节点”，说明父子之间真正需要传递的是高度、上下界、命中状态、层号或两类选择值，而不是反复扫描整棵子树。后面的优化只消掉这类重复工作，不能改变答案集合。

\textbf{优化条件。} 本题的关键观察是：每次选择区间中点，左右子区间构造左右子树。这句话必须能翻译成可维护状态；如果题目条件改变后观察不再成立，就不能继续套原来的写法。

\textbf{相邻状态。} 规律是有序数组的中点天然把值域分成 BST 左右两半，空区间返回空节点。把这个特例继续拆开看：节点向父亲返回的量和全局答案通常不是同一个量。先写叶子或空节点的返回值，再写父节点如何合并左右孩子，递归式才有边界基础。

\textbf{状态复用。} 把对子树的重复扫描压成返回值或队列状态；每个节点只合并孩子给出的稳定信息。本题落点是：规律是有序数组的中点天然把值域分成 BST 左右两半，空区间返回空节点。手算时只改被新输入影响的那一小部分，而不是回头重算完整候选。

\textbf{指针和字段。} 状态字段要能说清：节点指针、传入约束或返回值、当前答案、层号或路径状态；空节点的返回值必须和状态定义匹配。手算时按这个协议跟踪：拿 [-10, -3, 0, 5, 9] 手算，选 0 做根，左边两个数构成左子树，右边两个数构成右子树，高度接近平衡。然后按协议复查：手算时从叶子开始写返回值，或者从根开始写传入约束。不要只写访问顺序，要写每条边上传递的信息。

\textbf{代码落点。} 递归版本先处理空节点；迭代版本的栈元素要带完整状态；全负值、重复值和极端深树要单独测。关键代码行必须能逐句翻译成“旧状态怎样变成新状态”，否则就只是模板。

\textbf{正确性。} 证明从子树归纳或层序不变量开始：孩子状态正确时父节点合并正确；若用队列，当前轮队列必须正好保存当前层节点。

\textbf{边界实验。} 先用空区间、偶数长度选左中还是右中、高度平衡定义做反例检查。实验时覆盖空区间、偶数长度选左中还是右中、高度平衡定义，尤其关注空节点、单节点、链式树和全负值。递归版本和显式栈版本可以互相对拍。

\textbf{复杂度变化。} 暴力在每个节点重复扫描子树或反复寻找层边界可能退化到平方级；后序汇总、显式栈或层序遍历让每个节点处理常数次，通常是线性时间。

\textbf{迁移追问。} 如果题目改成“工程上递归可读，若严格避免递归可用队列保存区间和父指针”，先判断原状态是否仍然覆盖新答案集合，再决定是微调边界、改 value 语义，还是换算法。

\noindent\textbf{LeetCode 110 平衡二叉树。}

\textbf{题目说明。} LeetCode 110 平衡二叉树 的题面入口要先还原成具体任务：每个节点需要知道子树高度以及是否已经失衡。

\textbf{样例入口。} 拿链式树 1->2->3 手算，根的左右高度差为 2，已经不平衡；若子树已经失衡，父节点没必要继续精算高度。

\textbf{暴力基线。} 慢版本在每个节点重新扫描子树或路径，先求出局部答案。重复扫描暴露了真正状态：每个节点只需要从孩子或父亲那里拿到少量信息。放到本题就是：先按“每个节点需要知道子树高度以及是否已经失衡”完整试慢版本；样例里已经能看到“规律是后序返回高度或失败标记，空树高度为 0，左右差正好 1 仍合法”，所以优化前必须先能说清慢版本枚举了哪些候选。

\textbf{暴力过程。} 拿链式树 1->2->3 手算，根的左右高度差为 2，已经不平衡；若子树已经失衡，父节点没必要继续精算高度。

\textbf{重复工作。} 题面模型：每个节点需要知道子树高度以及是否已经失衡。暴力版的慢点要落到本题：慢版本会在树上重复确认“每个节点需要知道子树高度以及是否已经失衡”。样例规律是“规律是后序返回高度或失败标记，空树高度为 0，左右差正好 1 仍合法”，说明父子之间真正需要传递的是高度、上下界、命中状态、层号或两类选择值，而不是反复扫描整棵子树。后面的优化只消掉这类重复工作，不能改变答案集合。

\textbf{优化条件。} 本题的关键观察是：后序汇总时一旦子树失衡就向上返回失败标记，避免重复算高度。这句话必须能翻译成可维护状态；如果题目条件改变后观察不再成立，就不能继续套原来的写法。

\textbf{相邻状态。} 规律是后序返回高度或失败标记，空树高度为 0，左右差正好 1 仍合法。把这个特例继续拆开看：节点向父亲返回的量和全局答案通常不是同一个量。先写叶子或空节点的返回值，再写父节点如何合并左右孩子，递归式才有边界基础。

\textbf{状态复用。} 把对子树的重复扫描压成返回值或队列状态；每个节点只合并孩子给出的稳定信息。本题落点是：规律是后序返回高度或失败标记，空树高度为 0，左右差正好 1 仍合法。手算时只改被新输入影响的那一小部分，而不是回头重算完整候选。

\textbf{指针和字段。} 状态字段要能说清：节点指针、传入约束或返回值、当前答案、层号或路径状态；空节点的返回值必须和状态定义匹配。手算时按这个协议跟踪：拿链式树 1->2->3 手算，根的左右高度差为 2，已经不平衡；若子树已经失衡，父节点没必要继续精算高度。然后按协议复查：手算时从叶子开始写返回值，或者从根开始写传入约束。不要只写访问顺序，要写每条边上传递的信息。

\textbf{代码落点。} 递归版本先处理空节点；迭代版本的栈元素要带完整状态；全负值、重复值和极端深树要单独测。关键代码行必须能逐句翻译成“旧状态怎样变成新状态”，否则就只是模板。

\textbf{正确性。} 证明从子树归纳或层序不变量开始：孩子状态正确时父节点合并正确；若用队列，当前轮队列必须正好保存当前层节点。

\textbf{边界实验。} 先用空树、单节点、左右高度差正好一、链式树做反例检查。实验时覆盖空树、单节点、左右高度差正好一、链式树，尤其关注空节点、单节点、链式树和全负值。递归版本和显式栈版本可以互相对拍。

\textbf{复杂度变化。} 暴力在每个节点重复扫描子树或反复寻找层边界可能退化到平方级；后序汇总、显式栈或层序遍历让每个节点处理常数次，通常是线性时间。

\textbf{迁移追问。} 如果题目改成“可以把返回值设计成 optional 高度，空表示不平衡”，先判断原状态是否仍然覆盖新答案集合，再决定是微调边界、改 value 语义，还是换算法。

\noindent\textbf{LeetCode 111 二叉树最小深度。}

\textbf{题目说明。} LeetCode 111 二叉树最小深度 的题面入口要先还原成具体任务：最小深度是到最近叶子的节点数，不是左右深度简单取最小。

\textbf{样例入口。} 拿根 1 只有左孩子 2 手算，最小深度是 2，不是 min(左深度1, 右深度0)+1。

\textbf{暴力基线。} 慢版本在每个节点重新扫描子树或路径，先求出局部答案。重复扫描暴露了真正状态：每个节点只需要从孩子或父亲那里拿到少量信息。放到本题就是：先按“最小深度是到最近叶子的节点数，不是左右深度简单取最小”完整试慢版本；样例里已经能看到“规律是空子树不能当作叶子路径；BFS 第一次遇到叶子最直接，因为层次递增”，所以优化前必须先能说清慢版本枚举了哪些候选。

\textbf{暴力过程。} 拿根 1 只有左孩子 2 手算，最小深度是 2，不是 min(左深度1, 右深度0)+1。

\textbf{重复工作。} 题面模型：最小深度是到最近叶子的节点数，不是左右深度简单取最小。暴力版的慢点要落到本题：慢版本会反复确认“最小深度是到最近叶子的节点数，不是左右深度简单取最小”里的层边界。样例规律是“规律是空子树不能当作叶子路径；BFS 第一次遇到叶子最直接，因为层次递增”，说明队列当前轮的长度就是稳定边界；不先固定边界，下一层节点会混入当前层。后面的优化只消掉这类重复工作，不能改变答案集合。

\textbf{优化条件。} 本题的关键观察是：BFS 第一次遇到叶子就是最小深度，因为按层扩展。这句话必须能翻译成可维护状态；如果题目条件改变后观察不再成立，就不能继续套原来的写法。

\textbf{相邻状态。} 规律是空子树不能当作叶子路径；BFS 第一次遇到叶子最直接，因为层次递增。把这个特例继续拆开看：节点向父亲返回的量和全局答案通常不是同一个量。先写叶子或空节点的返回值，再写父节点如何合并左右孩子，递归式才有边界基础。

\textbf{状态复用。} 把对子树的重复扫描压成返回值或队列状态；每个节点只合并孩子给出的稳定信息。本题落点是：规律是空子树不能当作叶子路径；BFS 第一次遇到叶子最直接，因为层次递增。手算时只改被新输入影响的那一小部分，而不是回头重算完整候选。

\textbf{指针和字段。} 状态字段要能说清：节点指针、传入约束或返回值、当前答案、层号或路径状态；空节点的返回值必须和状态定义匹配。手算时按这个协议跟踪：拿根 1 只有左孩子 2 手算，最小深度是 2，不是 min(左深度1, 右深度0)+1。然后按协议复查：手算时从叶子开始写返回值，或者从根开始写传入约束。不要只写访问顺序，要写每条边上传递的信息。

\textbf{代码落点。} 递归版本先处理空节点；迭代版本的栈元素要带完整状态；全负值、重复值和极端深树要单独测。关键代码行必须能逐句翻译成“旧状态怎样变成新状态”，否则就只是模板。

\textbf{正确性。} 证明从子树归纳或层序不变量开始：孩子状态正确时父节点合并正确；若用队列，当前轮队列必须正好保存当前层节点。

\textbf{边界实验。} 先用只有左子树、只有右子树、根为空、根是叶子做反例检查。实验时覆盖只有左子树、只有右子树、根为空、根是叶子，尤其关注空节点、单节点、链式树和全负值。递归版本和显式栈版本可以互相对拍。

\textbf{复杂度变化。} 暴力在每个节点重复扫描子树或反复寻找层边界可能退化到平方级；后序汇总、显式栈或层序遍历让每个节点处理常数次，通常是线性时间。

\textbf{迁移追问。} 如果题目改成“DFS 写法必须特殊处理空子树不能作为最小路径”，先判断原状态是否仍然覆盖新答案集合，再决定是微调边界、改 value 语义，还是换算法。

\noindent\textbf{LeetCode 112 路径总和。}

\textbf{题目说明。} LeetCode 112 路径总和 的题面入口要先还原成具体任务：路径必须从根到叶子，状态是剩余目标值。

\textbf{样例入口。} 拿路径 5->4->11、target=20 手算，沿途把目标减成 15、11、0，到叶子时正好为 0 才成功；非叶子提前等于目标不算。

\textbf{暴力基线。} 慢版本在每个节点重新扫描子树或路径，先求出局部答案。重复扫描暴露了真正状态：每个节点只需要从孩子或父亲那里拿到少量信息。放到本题就是：先按“路径必须从根到叶子，状态是剩余目标值”完整试慢版本；样例里已经能看到“规律是状态是剩余目标值，终止条件必须同时满足叶子节点”，所以优化前必须先能说清慢版本枚举了哪些候选。

\textbf{暴力过程。} 拿路径 5->4->11、target=20 手算，沿途把目标减成 15、11、0，到叶子时正好为 0 才成功；非叶子提前等于目标不算。

\textbf{重复工作。} 题面模型：路径必须从根到叶子，状态是剩余目标值。暴力版的慢点要落到本题：慢版本会在树上重复确认“路径必须从根到叶子，状态是剩余目标值”。样例规律是“规律是状态是剩余目标值，终止条件必须同时满足叶子节点”，说明父子之间真正需要传递的是高度、上下界、命中状态、层号或两类选择值，而不是反复扫描整棵子树。后面的优化只消掉这类重复工作，不能改变答案集合。

\textbf{优化条件。} 本题的关键观察是：沿路径减去当前节点值，到叶子时检查剩余是否等于叶子值。这句话必须能翻译成可维护状态；如果题目条件改变后观察不再成立，就不能继续套原来的写法。

\textbf{相邻状态。} 规律是状态是剩余目标值，终止条件必须同时满足叶子节点。把这个特例继续拆开看：节点向父亲返回的量和全局答案通常不是同一个量。先写叶子或空节点的返回值，再写父节点如何合并左右孩子，递归式才有边界基础。

\textbf{状态复用。} 把对子树的重复扫描压成返回值或队列状态；每个节点只合并孩子给出的稳定信息。本题落点是：规律是状态是剩余目标值，终止条件必须同时满足叶子节点。手算时只改被新输入影响的那一小部分，而不是回头重算完整候选。

\textbf{指针和字段。} 状态字段要能说清：节点指针、传入约束或返回值、当前答案、层号或路径状态；空节点的返回值必须和状态定义匹配。手算时按这个协议跟踪：拿路径 5->4->11、target=20 手算，沿途把目标减成 15、11、0，到叶子时正好为 0 才成功；非叶子提前等于目标不算。然后按协议复查：手算时从叶子开始写返回值，或者从根开始写传入约束。不要只写访问顺序，要写每条边上传递的信息。

\textbf{代码落点。} 递归版本先处理空节点；迭代版本的栈元素要带完整状态；全负值、重复值和极端深树要单独测。关键代码行必须能逐句翻译成“旧状态怎样变成新状态”，否则就只是模板。

\textbf{正确性。} 证明从子树归纳或层序不变量开始：孩子状态正确时父节点合并正确；若用队列，当前轮队列必须正好保存当前层节点。

\textbf{边界实验。} 先用负数节点、目标为零、非叶子提前达到目标、空树做反例检查。实验时覆盖负数节点、目标为零、非叶子提前达到目标、空树，尤其关注空节点、单节点、链式树和全负值。递归版本和显式栈版本可以互相对拍。

\textbf{复杂度变化。} 暴力在每个节点重复扫描子树或反复寻找层边界可能退化到平方级；后序汇总、显式栈或层序遍历让每个节点处理常数次，通常是线性时间。

\textbf{迁移追问。} 如果题目改成“若要求列出所有路径，需要保存路径数组并在回退时弹出”，先判断原状态是否仍然覆盖新答案集合，再决定是微调边界、改 value 语义，还是换算法。

\noindent\textbf{LeetCode 113 路径总和 II。}

\textbf{题目说明。} LeetCode 113 路径总和 II 的题面入口要先还原成具体任务：相比判定题，状态还要保存当前路径。

\textbf{样例入口。} 拿 target=22 和路径 5->4->11->2 手算，到叶子时剩余为 0，需要复制当前路径；回退后要弹出 2，继续找别的分支。

\textbf{暴力基线。} 慢版本在每个节点重新扫描子树或路径，先求出局部答案。重复扫描暴露了真正状态：每个节点只需要从孩子或父亲那里拿到少量信息。放到本题就是：先按“相比判定题，状态还要保存当前路径”完整试慢版本；样例里已经能看到“规律是判定题多了路径数组，保存答案必须复制，不能保存同一个可变容器引用”，所以优化前必须先能说清慢版本枚举了哪些候选。

\textbf{暴力过程。} 拿 target=22 和路径 5->4->11->2 手算，到叶子时剩余为 0，需要复制当前路径；回退后要弹出 2，继续找别的分支。

\textbf{重复工作。} 题面模型：相比判定题，状态还要保存当前路径。暴力版的慢点要落到本题：慢版本会在树上重复确认“相比判定题，状态还要保存当前路径”。样例规律是“规律是判定题多了路径数组，保存答案必须复制，不能保存同一个可变容器引用”，说明父子之间真正需要传递的是高度、上下界、命中状态、层号或两类选择值，而不是反复扫描整棵子树。后面的优化只消掉这类重复工作，不能改变答案集合。

\textbf{优化条件。} 本题的关键观察是：到叶子且和满足时复制路径；回溯时恢复路径。这句话必须能翻译成可维护状态；如果题目条件改变后观察不再成立，就不能继续套原来的写法。

\textbf{相邻状态。} 规律是判定题多了路径数组，保存答案必须复制，不能保存同一个可变容器引用。把这个特例继续拆开看：节点向父亲返回的量和全局答案通常不是同一个量。先写叶子或空节点的返回值，再写父节点如何合并左右孩子，递归式才有边界基础。

\textbf{状态复用。} 把对子树的重复扫描压成返回值或队列状态；每个节点只合并孩子给出的稳定信息。本题落点是：规律是判定题多了路径数组，保存答案必须复制，不能保存同一个可变容器引用。手算时只改被新输入影响的那一小部分，而不是回头重算完整候选。

\textbf{指针和字段。} 状态字段要能说清：节点指针、传入约束或返回值、当前答案、层号或路径状态；空节点的返回值必须和状态定义匹配。手算时按这个协议跟踪：拿 target=22 和路径 5->4->11->2 手算，到叶子时剩余为 0，需要复制当前路径；回退后要弹出 2，继续找别的分支。然后按协议复查：手算时从叶子开始写返回值，或者从根开始写传入约束。不要只写访问顺序，要写每条边上传递的信息。

\textbf{代码落点。} 递归版本先处理空节点；迭代版本的栈元素要带完整状态；全负值、重复值和极端深树要单独测。关键代码行必须能逐句翻译成“旧状态怎样变成新状态”，否则就只是模板。

\textbf{正确性。} 证明从子树归纳或层序不变量开始：孩子状态正确时父节点合并正确；若用队列，当前轮队列必须正好保存当前层节点。

\textbf{边界实验。} 先用多个答案、负数、空树、路径数组复制时机做反例检查。实验时覆盖多个答案、负数、空树、路径数组复制时机，尤其关注空节点、单节点、链式树和全负值。递归版本和显式栈版本可以互相对拍。

\textbf{复杂度变化。} 暴力在每个节点重复扫描子树或反复寻找层边界可能退化到平方级；后序汇总、显式栈或层序遍历让每个节点处理常数次，通常是线性时间。

\textbf{迁移追问。} 如果题目改成“若路径不要求从根到叶，模型会变成前缀和加哈希”，先判断原状态是否仍然覆盖新答案集合，再决定是微调边界、改 value 语义，还是换算法。

\noindent\textbf{LeetCode 114 二叉树展开为链表。}

\textbf{题目说明。} LeetCode 114 二叉树展开为链表 的题面入口要先还原成具体任务：目标是按前序顺序把树原地改成右指针链。

\textbf{样例入口。} 拿 1 的左子树 2 和右子树 5 手算，前序链应该是 1->2->...->5；若直接把左子树接到右边而不保存原右子树，5 会丢失。

\textbf{暴力基线。} 慢版本在每个节点重新扫描子树或路径，先求出局部答案。重复扫描暴露了真正状态：每个节点只需要从孩子或父亲那里拿到少量信息。放到本题就是：先按“目标是按前序顺序把树原地改成右指针链”完整试慢版本；样例里已经能看到“规律是按前序顺序重连，左指针置空，原右子树必须接到左子树展开后的尾部”，所以优化前必须先能说清慢版本枚举了哪些候选。

\textbf{暴力过程。} 拿 1 的左子树 2 和右子树 5 手算，前序链应该是 1->2->...->5；若直接把左子树接到右边而不保存原右子树，5 会丢失。

\textbf{重复工作。} 题面模型：目标是按前序顺序把树原地改成右指针链。暴力版的慢点要落到本题：慢版本会在树上重复确认“目标是按前序顺序把树原地改成右指针链”。样例规律是“规律是按前序顺序重连，左指针置空，原右子树必须接到左子树展开后的尾部”，说明父子之间真正需要传递的是高度、上下界、命中状态、层号或两类选择值，而不是反复扫描整棵子树。后面的优化只消掉这类重复工作，不能改变答案集合。

\textbf{优化条件。} 本题的关键观察是：显式栈按右后左先入顺序处理，逐个接到前驱右侧。这句话必须能翻译成可维护状态；如果题目条件改变后观察不再成立，就不能继续套原来的写法。

\textbf{相邻状态。} 规律是按前序顺序重连，左指针置空，原右子树必须接到左子树展开后的尾部。把这个特例继续拆开看：节点向父亲返回的量和全局答案通常不是同一个量。先写叶子或空节点的返回值，再写父节点如何合并左右孩子，递归式才有边界基础。

\textbf{状态复用。} 把对子树的重复扫描压成返回值或队列状态；每个节点只合并孩子给出的稳定信息。本题落点是：规律是按前序顺序重连，左指针置空，原右子树必须接到左子树展开后的尾部。手算时只改被新输入影响的那一小部分，而不是回头重算完整候选。

\textbf{指针和字段。} 状态字段要能说清：节点指针、传入约束或返回值、当前答案、层号或路径状态；空节点的返回值必须和状态定义匹配。手算时按这个协议跟踪：拿 1 的左子树 2 和右子树 5 手算，前序链应该是 1->2->...->5；若直接把左子树接到右边而不保存原右子树，5 会丢失。然后按协议复查：手算时从叶子开始写返回值，或者从根开始写传入约束。不要只写访问顺序，要写每条边上传递的信息。

\textbf{代码落点。} 递归版本先处理空节点；迭代版本的栈元素要带完整状态；全负值、重复值和极端深树要单独测。关键代码行必须能逐句翻译成“旧状态怎样变成新状态”，否则就只是模板。

\textbf{正确性。} 证明从子树归纳或层序不变量开始：孩子状态正确时父节点合并正确；若用队列，当前轮队列必须正好保存当前层节点。

\textbf{边界实验。} 先用空树、单节点、原有右子树不能丢、左指针置空做反例检查。实验时覆盖空树、单节点、原有右子树不能丢、左指针置空，尤其关注空节点、单节点、链式树和全负值。递归版本和显式栈版本可以互相对拍。

\textbf{复杂度变化。} 暴力在每个节点重复扫描子树或反复寻找层边界可能退化到平方级；后序汇总、显式栈或层序遍历让每个节点处理常数次，通常是线性时间。

\textbf{迁移追问。} 如果题目改成“也可用寻找左子树最右节点的原地方法”，先判断原状态是否仍然覆盖新答案集合，再决定是微调边界、改 value 语义，还是换算法。

\noindent\textbf{LeetCode 124 二叉树最大路径和。}

\textbf{题目说明。} LeetCode 124 二叉树最大路径和 的题面入口要先还原成具体任务：路径可以从任意节点到任意节点，但不能分叉向父节点贡献两边。

\textbf{样例入口。} 拿根 -10、左 9、右 20 手算，经过 20 的路径可以同时吃左增益 15 和右增益 7 更新全局答案，但向父亲只能贡献一条单边路径。

\textbf{暴力基线。} 暴力可以枚举每个节点作为路径最高拐点，再分别求左右两侧向下最大路径和。

\textbf{暴力过程。} 以根 -10、左 9、右 20、20 的孩子 15 和 7 为例，最大路径是 15->20->7，和为 42，不需要经过根。

\textbf{重复工作。} 题面模型：路径可以从任意节点到任意节点，但不能分叉向父节点贡献两边。暴力版的慢点要落到本题：浪费在对每个拐点重复求子树贡献。每个节点只需要向父亲返回一条单边最大增益，同时用左右两边更新全局答案。后面的优化只消掉这类重复工作，不能改变答案集合。

\textbf{优化条件。} 本题的关键观察是：每个节点向父亲贡献单边最大增益，全局答案可用左右两边加当前节点更新。这句话必须能翻译成可维护状态；如果题目条件改变后观察不再成立，就不能继续套原来的写法。

\textbf{相邻状态。} 规律是返回给父亲的是单边最大增益，全局答案才允许左右两边加当前节点。把这个特例继续拆开看：节点向父亲返回的量和全局答案通常不是同一个量。先写叶子或空节点的返回值，再写父节点如何合并左右孩子，递归式才有边界基础。

\textbf{状态复用。} 后序遍历。gain(node)=node.val+max(0, left\_gain, right\_gain)，全局答案用 node.val+max(0, left\_gain)+max(0, right\_gain) 更新。手算时只改被新输入影响的那一小部分，而不是回头重算完整候选。

\textbf{指针和字段。} 状态字段要能说清：left\_gain/right\_gain 是孩子能向当前节点贡献的单边路径，answer 是允许在当前节点拐弯的全局最大路径和。手算时按这个协议跟踪：节点 20 收到 15 和 7 的增益，向父亲只能返回 20+15=35，但全局答案可用 15+20+7=42 更新。

\textbf{代码落点。} 关键是负增益取 0，表示不接这条子路径；answer 初始化为很小值，不能初始化为 0，否则全负树会错。关键代码行必须能逐句翻译成“旧状态怎样变成新状态”，否则就只是模板。

\textbf{正确性。} 一条简单路径在某个最高节点处最多连接左右两条向下路径；向父亲延伸时只能选择其中一条，否则会分叉。

\textbf{边界实验。} 先用全负数、单节点、左右增益为负时取零、答案初始值做反例检查。实验时覆盖全负数、单节点、左右增益为负时取零、答案初始值，尤其关注空节点、单节点、链式树和全负值。递归版本和显式栈版本可以互相对拍。

\textbf{复杂度变化。} 暴力枚举拐点并求贡献最坏 O(\ensuremath{n^{2}})，后序状态 O(n) 时间，空间 O(h)。

\textbf{迁移追问。} 如果题目改成“若路径必须从根到叶，状态和答案位置完全不同”，先判断原状态是否仍然覆盖新答案集合，再决定是微调边界、改 value 语义，还是换算法。

\noindent\textbf{LeetCode 199 二叉树右视图。}

\textbf{题目说明。} LeetCode 199 二叉树右视图 的题面入口要先还原成具体任务：每层从右侧能看到的就是该层最后访问或最先访问的右侧节点。

\textbf{样例入口。} 拿根 1、左 2、右 3 手算，从右侧看第一层是 1，第二层是 3；若右边缺失，才会看到左侧节点。

\textbf{暴力基线。} 慢版本在每个节点重新扫描子树或路径，先求出局部答案。重复扫描暴露了真正状态：每个节点只需要从孩子或父亲那里拿到少量信息。放到本题就是：先按“每层从右侧能看到的就是该层最后访问或最先访问的右侧节点”完整试慢版本；样例里已经能看到“规律是 BFS 每层最后一个节点，或 DFS 优先右子树并第一次到达某层时记录”，所以优化前必须先能说清慢版本枚举了哪些候选。

\textbf{暴力过程。} 拿根 1、左 2、右 3 手算，从右侧看第一层是 1，第二层是 3；若右边缺失，才会看到左侧节点。

\textbf{重复工作。} 题面模型：每层从右侧能看到的就是该层最后访问或最先访问的右侧节点。暴力版的慢点要落到本题：慢版本会反复确认“每层从右侧能看到的就是该层最后访问或最先访问的右侧节点”里的层边界。样例规律是“规律是 BFS 每层最后一个节点，或 DFS 优先右子树并第一次到达某层时记录”，说明队列当前轮的长度就是稳定边界；不先固定边界，下一层节点会混入当前层。后面的优化只消掉这类重复工作，不能改变答案集合。

\textbf{优化条件。} 本题的关键观察是：BFS 固定层边界，取每层最后一个；或 DFS 按右优先首次到达深度。这句话必须能翻译成可维护状态；如果题目条件改变后观察不再成立，就不能继续套原来的写法。

\textbf{相邻状态。} 规律是 BFS 每层最后一个节点，或 DFS 优先右子树并第一次到达某层时记录。把这个特例继续拆开看：节点向父亲返回的量和全局答案通常不是同一个量。先写叶子或空节点的返回值，再写父节点如何合并左右孩子，递归式才有边界基础。

\textbf{状态复用。} 把对子树的重复扫描压成返回值或队列状态；每个节点只合并孩子给出的稳定信息。本题落点是：规律是 BFS 每层最后一个节点，或 DFS 优先右子树并第一次到达某层时记录。手算时只改被新输入影响的那一小部分，而不是回头重算完整候选。

\textbf{指针和字段。} 状态字段要能说清：节点指针、传入约束或返回值、当前答案、层号或路径状态；空节点的返回值必须和状态定义匹配。手算时按这个协议跟踪：拿根 1、左 2、右 3 手算，从右侧看第一层是 1，第二层是 3；若右边缺失，才会看到左侧节点。然后按协议复查：手算时从叶子开始写返回值，或者从根开始写传入约束。不要只写访问顺序，要写每条边上传递的信息。

\textbf{代码落点。} 递归版本先处理空节点；迭代版本的栈元素要带完整状态；全负值、重复值和极端深树要单独测。关键代码行必须能逐句翻译成“旧状态怎样变成新状态”，否则就只是模板。

\textbf{正确性。} 证明从子树归纳或层序不变量开始：孩子状态正确时父节点合并正确；若用队列，当前轮队列必须正好保存当前层节点。

\textbf{边界实验。} 先用空树、左链、右链、左右交错做反例检查。实验时覆盖空树、左链、右链、左右交错，尤其关注空节点、单节点、链式树和全负值。递归版本和显式栈版本可以互相对拍。

\textbf{复杂度变化。} 暴力在每个节点重复扫描子树或反复寻找层边界可能退化到平方级；后序汇总、显式栈或层序遍历让每个节点处理常数次，通常是线性时间。

\textbf{迁移追问。} 如果题目改成“若求左视图，只改变层内取值方向”，先判断原状态是否仍然覆盖新答案集合，再决定是微调边界、改 value 语义，还是换算法。

\noindent\textbf{LeetCode 208 实现 Trie。}

\textbf{题目说明。} LeetCode 208 实现 Trie 的题面入口要先还原成具体任务：前缀树把字符串公共前缀共享成路径。

\textbf{样例入口。} 拿插入 apple 再查 app 手算，app 的路径存在，但如果 p 节点没有终止标记，search(app) 应为假，startsWith(app) 才为真。

\textbf{暴力基线。} 慢版本在每个节点重新扫描子树或路径，先求出局部答案。重复扫描暴露了真正状态：每个节点只需要从孩子或父亲那里拿到少量信息。放到本题就是：先按“前缀树把字符串公共前缀共享成路径”完整试慢版本；样例里已经能看到“规律是 Trie 节点既要有孩子边，也要有单词终止标记”，所以优化前必须先能说清慢版本枚举了哪些候选。

\textbf{暴力过程。} 拿插入 apple 再查 app 手算，app 的路径存在，但如果 p 节点没有终止标记，search(app) 应为假，startsWith(app) 才为真。

\textbf{重复工作。} 题面模型：前缀树把字符串公共前缀共享成路径。暴力版的慢点要落到本题：慢版本会在树上重复确认“前缀树把字符串公共前缀共享成路径”。样例规律是“规律是 Trie 节点既要有孩子边，也要有单词终止标记”，说明父子之间真正需要传递的是高度、上下界、命中状态、层号或两类选择值，而不是反复扫描整棵子树。后面的优化只消掉这类重复工作，不能改变答案集合。

\textbf{优化条件。} 本题的关键观察是：插入和查询都逐字符沿边移动，不存在的边按需要创建。这句话必须能翻译成可维护状态；如果题目条件改变后观察不再成立，就不能继续套原来的写法。

\textbf{相邻状态。} 规律是 Trie 节点既要有孩子边，也要有单词终止标记。把这个特例继续拆开看：节点向父亲返回的量和全局答案通常不是同一个量。先写叶子或空节点的返回值，再写父节点如何合并左右孩子，递归式才有边界基础。

\textbf{状态复用。} 把对子树的重复扫描压成返回值或队列状态；每个节点只合并孩子给出的稳定信息。本题落点是：规律是 Trie 节点既要有孩子边，也要有单词终止标记。手算时只改被新输入影响的那一小部分，而不是回头重算完整候选。

\textbf{指针和字段。} 状态字段要能说清：节点指针、传入约束或返回值、当前答案、层号或路径状态；空节点的返回值必须和状态定义匹配。手算时按这个协议跟踪：拿插入 apple 再查 app 手算，app 的路径存在，但如果 p 节点没有终止标记，search(app) 应为假，startsWith(app) 才为真。然后按协议复查：手算时从叶子开始写返回值，或者从根开始写传入约束。不要只写访问顺序，要写每条边上传递的信息。

\textbf{代码落点。} 递归版本先处理空节点；迭代版本的栈元素要带完整状态；全负值、重复值和极端深树要单独测。关键代码行必须能逐句翻译成“旧状态怎样变成新状态”，否则就只是模板。

\textbf{正确性。} 证明从子树归纳或层序不变量开始：孩子状态正确时父节点合并正确；若用队列，当前轮队列必须正好保存当前层节点。

\textbf{边界实验。} 先用空字符串约定、单词结束标记、前缀存在不等于单词存在做反例检查。实验时覆盖空字符串约定、单词结束标记、前缀存在不等于单词存在，尤其关注空节点、单节点、链式树和全负值。递归版本和显式栈版本可以互相对拍。

\textbf{复杂度变化。} 暴力在每个节点重复扫描子树或反复寻找层边界可能退化到平方级；后序汇总、显式栈或层序遍历让每个节点处理常数次，通常是线性时间。

\textbf{迁移追问。} 如果题目改成“若字符集很大，用哈希子节点；小写字母用数组更快”，先判断原状态是否仍然覆盖新答案集合，再决定是微调边界、改 value 语义，还是换算法。

\noindent\textbf{LeetCode 211 添加与搜索单词。}

\textbf{题目说明。} LeetCode 211 添加与搜索单词 的题面入口要先还原成具体任务：点号通配符让查询可能分叉，但前缀树仍能剪掉不匹配前缀。

\textbf{样例入口。} 拿 add bad、dad、mad 后查 .ad 手算，点号可以走 b、d、m 三条边；查 b.. 时后两个点继续在子树里展开。

\textbf{暴力基线。} 慢版本在每个节点重新扫描子树或路径，先求出局部答案。重复扫描暴露了真正状态：每个节点只需要从孩子或父亲那里拿到少量信息。放到本题就是：先按“点号通配符让查询可能分叉，但前缀树仍能剪掉不匹配前缀”完整试慢版本；样例里已经能看到“规律是普通字符沿唯一边走，点号触发 DFS 枚举孩子，终点仍要检查单词结束标记”，所以优化前必须先能说清慢版本枚举了哪些候选。

\textbf{暴力过程。} 拿 add bad、dad、mad 后查 .ad 手算，点号可以走 b、d、m 三条边；查 b.. 时后两个点继续在子树里展开。

\textbf{重复工作。} 题面模型：点号通配符让查询可能分叉，但前缀树仍能剪掉不匹配前缀。暴力版的慢点要落到本题：慢版本会在树上重复确认“点号通配符让查询可能分叉，但前缀树仍能剪掉不匹配前缀”。样例规律是“规律是普通字符沿唯一边走，点号触发 DFS 枚举孩子，终点仍要检查单词结束标记”，说明父子之间真正需要传递的是高度、上下界、命中状态、层号或两类选择值，而不是反复扫描整棵子树。后面的优化只消掉这类重复工作，不能改变答案集合。

\textbf{优化条件。} 本题的关键观察是：普通字符沿唯一边走，点号枚举当前节点所有孩子。这句话必须能翻译成可维护状态；如果题目条件改变后观察不再成立，就不能继续套原来的写法。

\textbf{相邻状态。} 规律是普通字符沿唯一边走，点号触发 DFS 枚举孩子，终点仍要检查单词结束标记。把这个特例继续拆开看：节点向父亲返回的量和全局答案通常不是同一个量。先写叶子或空节点的返回值，再写父节点如何合并左右孩子，递归式才有边界基础。

\textbf{状态复用。} 把对子树的重复扫描压成返回值或队列状态；每个节点只合并孩子给出的稳定信息。本题落点是：规律是普通字符沿唯一边走，点号触发 DFS 枚举孩子，终点仍要检查单词结束标记。手算时只改被新输入影响的那一小部分，而不是回头重算完整候选。

\textbf{指针和字段。} 状态字段要能说清：节点指针、传入约束或返回值、当前答案、层号或路径状态；空节点的返回值必须和状态定义匹配。手算时按这个协议跟踪：拿 add bad、dad、mad 后查 .ad 手算，点号可以走 b、d、m 三条边；查 b.. 时后两个点继续在子树里展开。然后按协议复查：手算时从叶子开始写返回值，或者从根开始写传入约束。不要只写访问顺序，要写每条边上传递的信息。

\textbf{代码落点。} 递归版本先处理空节点；迭代版本的栈元素要带完整状态；全负值、重复值和极端深树要单独测。关键代码行必须能逐句翻译成“旧状态怎样变成新状态”，否则就只是模板。

\textbf{正确性。} 证明从子树归纳或层序不变量开始：孩子状态正确时父节点合并正确；若用队列，当前轮队列必须正好保存当前层节点。

\textbf{边界实验。} 先用连续通配符、空结果、单词结束标记、字符集大小做反例检查。实验时覆盖连续通配符、空结果、单词结束标记、字符集大小，尤其关注空节点、单节点、链式树和全负值。递归版本和显式栈版本可以互相对拍。

\textbf{复杂度变化。} 暴力在每个节点重复扫描子树或反复寻找层边界可能退化到平方级；后序汇总、显式栈或层序遍历让每个节点处理常数次，通常是线性时间。

\textbf{迁移追问。} 如果题目改成“大量通配符会接近遍历整棵 Trie，需要规模约束”，先判断原状态是否仍然覆盖新答案集合，再决定是微调边界、改 value 语义，还是换算法。

\noindent\textbf{LeetCode 226 翻转二叉树。}

\textbf{题目说明。} LeetCode 226 翻转二叉树 的题面入口要先还原成具体任务：每个节点交换左右孩子即可，遍历顺序不影响最终结构。

\textbf{样例入口。} 拿根 4、左 2、右 7 手算，翻转后左右孩子交换，2 和 7 的子树也要递归交换。

\textbf{暴力基线。} 慢版本在每个节点重新扫描子树或路径，先求出局部答案。重复扫描暴露了真正状态：每个节点只需要从孩子或父亲那里拿到少量信息。放到本题就是：先按“每个节点交换左右孩子即可，遍历顺序不影响最终结构”完整试慢版本；样例里已经能看到“规律是每个节点只负责交换自己的左右指针，子树正确性由递归返回；空节点直接返回”，所以优化前必须先能说清慢版本枚举了哪些候选。

\textbf{暴力过程。} 拿根 4、左 2、右 7 手算，翻转后左右孩子交换，2 和 7 的子树也要递归交换。

\textbf{重复工作。} 题面模型：每个节点交换左右孩子即可，遍历顺序不影响最终结构。暴力版的慢点要落到本题：慢版本会反复确认“每个节点交换左右孩子即可，遍历顺序不影响最终结构”里的层边界。样例规律是“规律是每个节点只负责交换自己的左右指针，子树正确性由递归返回；空节点直接返回”，说明队列当前轮的长度就是稳定边界；不先固定边界，下一层节点会混入当前层。后面的优化只消掉这类重复工作，不能改变答案集合。

\textbf{优化条件。} 本题的关键观察是：显式队列或栈逐个节点交换，避免递归深度。这句话必须能翻译成可维护状态；如果题目条件改变后观察不再成立，就不能继续套原来的写法。

\textbf{相邻状态。} 规律是每个节点只负责交换自己的左右指针，子树正确性由递归返回；空节点直接返回。把这个特例继续拆开看：节点向父亲返回的量和全局答案通常不是同一个量。先写叶子或空节点的返回值，再写父节点如何合并左右孩子，递归式才有边界基础。

\textbf{状态复用。} 把对子树的重复扫描压成返回值或队列状态；每个节点只合并孩子给出的稳定信息。本题落点是：规律是每个节点只负责交换自己的左右指针，子树正确性由递归返回；空节点直接返回。手算时只改被新输入影响的那一小部分，而不是回头重算完整候选。

\textbf{指针和字段。} 状态字段要能说清：节点指针、传入约束或返回值、当前答案、层号或路径状态；空节点的返回值必须和状态定义匹配。手算时按这个协议跟踪：拿根 4、左 2、右 7 手算，翻转后左右孩子交换，2 和 7 的子树也要递归交换。然后按协议复查：手算时从叶子开始写返回值，或者从根开始写传入约束。不要只写访问顺序，要写每条边上传递的信息。

\textbf{代码落点。} 递归版本先处理空节点；迭代版本的栈元素要带完整状态；全负值、重复值和极端深树要单独测。关键代码行必须能逐句翻译成“旧状态怎样变成新状态”，否则就只是模板。

\textbf{正确性。} 证明从子树归纳或层序不变量开始：孩子状态正确时父节点合并正确；若用队列，当前轮队列必须正好保存当前层节点。

\textbf{边界实验。} 先用空树、单节点、只有一侧子树、重复值做反例检查。实验时覆盖空树、单节点、只有一侧子树、重复值，尤其关注空节点、单节点、链式树和全负值。递归版本和显式栈版本可以互相对拍。

\textbf{复杂度变化。} 暴力在每个节点重复扫描子树或反复寻找层边界可能退化到平方级；后序汇总、显式栈或层序遍历让每个节点处理常数次，通常是线性时间。

\textbf{迁移追问。} 如果题目改成“这题简单但适合训练树节点指针修改”，先判断原状态是否仍然覆盖新答案集合，再决定是微调边界、改 value 语义，还是换算法。

\noindent\textbf{LeetCode 230 二叉搜索树中第 K 小。}

\textbf{题目说明。} LeetCode 230 二叉搜索树中第 K 小 的题面入口要先还原成具体任务：BST 中序遍历按升序访问节点。

\textbf{样例入口。} 拿 BST [3, 1, 4, null, 2]、k=1 手算，中序访问顺序是 1, 2, 3, 4，第一个就是答案。

\textbf{暴力基线。} 慢版本在每个节点重新扫描子树或路径，先求出局部答案。重复扫描暴露了真正状态：每个节点只需要从孩子或父亲那里拿到少量信息。放到本题就是：先按“BST 中序遍历按升序访问节点”完整试慢版本；样例里已经能看到“规律是 BST 的中序严格递增，计数到第 k 个即可停止；重复值是否允许要按题意定义”，所以优化前必须先能说清慢版本枚举了哪些候选。

\textbf{暴力过程。} 拿 BST [3, 1, 4, null, 2]、k=1 手算，中序访问顺序是 1, 2, 3, 4，第一个就是答案。

\textbf{重复工作。} 题面模型：BST 中序遍历按升序访问节点。暴力版的慢点要落到本题：慢版本会在树上重复确认“BST 中序遍历按升序访问节点”。样例规律是“规律是 BST 的中序严格递增，计数到第 k 个即可停止；重复值是否允许要按题意定义”，说明父子之间真正需要传递的是高度、上下界、命中状态、层号或两类选择值，而不是反复扫描整棵子树。后面的优化只消掉这类重复工作，不能改变答案集合。

\textbf{优化条件。} 本题的关键观察是：显式栈中序，访问第 k 个节点时返回。这句话必须能翻译成可维护状态；如果题目条件改变后观察不再成立，就不能继续套原来的写法。

\textbf{相邻状态。} 规律是 BST 的中序严格递增，计数到第 k 个即可停止；重复值是否允许要按题意定义。把这个特例继续拆开看：节点向父亲返回的量和全局答案通常不是同一个量。先写叶子或空节点的返回值，再写父节点如何合并左右孩子，递归式才有边界基础。

\textbf{状态复用。} 把对子树的重复扫描压成返回值或队列状态；每个节点只合并孩子给出的稳定信息。本题落点是：规律是 BST 的中序严格递增，计数到第 k 个即可停止；重复值是否允许要按题意定义。手算时只改被新输入影响的那一小部分，而不是回头重算完整候选。

\textbf{指针和字段。} 状态字段要能说清：节点指针、传入约束或返回值、当前答案、层号或路径状态；空节点的返回值必须和状态定义匹配。手算时按这个协议跟踪：拿 BST [3, 1, 4, null, 2]、k=1 手算，中序访问顺序是 1, 2, 3, 4，第一个就是答案。然后按协议复查：手算时从叶子开始写返回值，或者从根开始写传入约束。不要只写访问顺序，要写每条边上传递的信息。

\textbf{代码落点。} 递归版本先处理空节点；迭代版本的栈元素要带完整状态；全负值、重复值和极端深树要单独测。关键代码行必须能逐句翻译成“旧状态怎样变成新状态”，否则就只是模板。

\textbf{正确性。} 证明从子树归纳或层序不变量开始：孩子状态正确时父节点合并正确；若用队列，当前轮队列必须正好保存当前层节点。

\textbf{边界实验。} 先用k 为一、k 为节点数、树不平衡、是否有重复值做反例检查。实验时覆盖k 为一、k 为节点数、树不平衡、是否有重复值，尤其关注空节点、单节点、链式树和全负值。递归版本和显式栈版本可以互相对拍。

\textbf{复杂度变化。} 暴力在每个节点重复扫描子树或反复寻找层边界可能退化到平方级；后序汇总、显式栈或层序遍历让每个节点处理常数次，通常是线性时间。

\textbf{迁移追问。} 如果题目改成“若频繁查询，可在节点维护子树大小成为顺序统计树”，先判断原状态是否仍然覆盖新答案集合，再决定是微调边界、改 value 语义，还是换算法。

\noindent\textbf{LeetCode 235 二叉搜索树最近公共祖先。}

\textbf{题目说明。} LeetCode 235 二叉搜索树最近公共祖先 的题面入口要先还原成具体任务：BST 的值域关系能决定两个目标在当前节点的哪一侧。

\textbf{样例入口。} 拿 BST 根 6，p=2，q=8 手算，2 在左侧、8 在右侧，根 6 就是最近公共祖先；若 p 和 q 都小于根，就继续去左子树。

\textbf{暴力基线。} 慢版本在每个节点重新扫描子树或路径，先求出局部答案。重复扫描暴露了真正状态：每个节点只需要从孩子或父亲那里拿到少量信息。放到本题就是：先按“BST 的值域关系能决定两个目标在当前节点的哪一侧”完整试慢版本；样例里已经能看到“规律是 BST 的值域能直接决定搜索方向”，所以优化前必须先能说清慢版本枚举了哪些候选。

\textbf{暴力过程。} 拿 BST 根 6，p=2，q=8 手算，2 在左侧、8 在右侧，根 6 就是最近公共祖先；若 p 和 q 都小于根，就继续去左子树。

\textbf{重复工作。} 题面模型：BST 的值域关系能决定两个目标在当前节点的哪一侧。暴力版的慢点要落到本题：慢版本会在树上重复确认“BST 的值域关系能决定两个目标在当前节点的哪一侧”。样例规律是“规律是 BST 的值域能直接决定搜索方向”，说明父子之间真正需要传递的是高度、上下界、命中状态、层号或两类选择值，而不是反复扫描整棵子树。后面的优化只消掉这类重复工作，不能改变答案集合。

\textbf{优化条件。} 本题的关键观察是：两个值都小去左，都大去右，否则当前节点就是分叉点。这句话必须能翻译成可维护状态；如果题目条件改变后观察不再成立，就不能继续套原来的写法。

\textbf{相邻状态。} 规律是 BST 的值域能直接决定搜索方向。把这个特例继续拆开看：节点向父亲返回的量和全局答案通常不是同一个量。先写叶子或空节点的返回值，再写父节点如何合并左右孩子，递归式才有边界基础。

\textbf{状态复用。} 把对子树的重复扫描压成返回值或队列状态；每个节点只合并孩子给出的稳定信息。本题落点是：规律是 BST 的值域能直接决定搜索方向。手算时只改被新输入影响的那一小部分，而不是回头重算完整候选。

\textbf{指针和字段。} 状态字段要能说清：节点指针、传入约束或返回值、当前答案、层号或路径状态；空节点的返回值必须和状态定义匹配。手算时按这个协议跟踪：拿 BST 根 6，p=2，q=8 手算，2 在左侧、8 在右侧，根 6 就是最近公共祖先；若 p 和 q 都小于根，就继续去左子树。然后按协议复查：手算时从叶子开始写返回值，或者从根开始写传入约束。不要只写访问顺序，要写每条边上传递的信息。

\textbf{代码落点。} 递归版本先处理空节点；迭代版本的栈元素要带完整状态；全负值、重复值和极端深树要单独测。关键代码行必须能逐句翻译成“旧状态怎样变成新状态”，否则就只是模板。

\textbf{正确性。} 证明从子树归纳或层序不变量开始：孩子状态正确时父节点合并正确；若用队列，当前轮队列必须正好保存当前层节点。

\textbf{边界实验。} 先用一个节点是另一个祖先、目标顺序、重复值约定做反例检查。实验时覆盖一个节点是另一个祖先、目标顺序、重复值约定，尤其关注空节点、单节点、链式树和全负值。递归版本和显式栈版本可以互相对拍。

\textbf{复杂度变化。} 暴力在每个节点重复扫描子树或反复寻找层边界可能退化到平方级；后序汇总、显式栈或层序遍历让每个节点处理常数次，通常是线性时间。

\textbf{迁移追问。} 如果题目改成“普通二叉树不能用值域剪枝，需要父指针或后序状态”，先判断原状态是否仍然覆盖新答案集合，再决定是微调边界、改 value 语义，还是换算法。

\noindent\textbf{LeetCode 236 二叉树最近公共祖先。}

\textbf{题目说明。} LeetCode 236 二叉树最近公共祖先 的题面入口要先还原成具体任务：普通树没有有序性，只能依靠祖先关系。

\textbf{样例入口。} 拿普通二叉树根 3，p=5、q=1 手算，左右子树分别找到一个目标，根 3 就是最近公共祖先；若两个目标都在左子树，答案应由左子树返回。

\textbf{暴力基线。} 慢版本在每个节点重新扫描子树或路径，先求出局部答案。重复扫描暴露了真正状态：每个节点只需要从孩子或父亲那里拿到少量信息。放到本题就是：先按“普通树没有有序性，只能依靠祖先关系”完整试慢版本；样例里已经能看到“规律是后序返回是否找到目标或祖先，不能用 BST 的大小关系”，所以优化前必须先能说清慢版本枚举了哪些候选。

\textbf{暴力过程。} 拿普通二叉树根 3，p=5、q=1 手算，左右子树分别找到一个目标，根 3 就是最近公共祖先；若两个目标都在左子树，答案应由左子树返回。

\textbf{重复工作。} 题面模型：普通树没有有序性，只能依靠祖先关系。暴力版的慢点要落到本题：慢版本会在树上重复确认“普通树没有有序性，只能依靠祖先关系”。样例规律是“规律是后序返回是否找到目标或祖先，不能用 BST 的大小关系”，说明父子之间真正需要传递的是高度、上下界、命中状态、层号或两类选择值，而不是反复扫描整棵子树。后面的优化只消掉这类重复工作，不能改变答案集合。

\textbf{优化条件。} 本题的关键观察是：父指针集合或后序返回目标命中状态都可行。这句话必须能翻译成可维护状态；如果题目条件改变后观察不再成立，就不能继续套原来的写法。

\textbf{相邻状态。} 规律是后序返回是否找到目标或祖先，不能用 BST 的大小关系。把这个特例继续拆开看：节点向父亲返回的量和全局答案通常不是同一个量。先写叶子或空节点的返回值，再写父节点如何合并左右孩子，递归式才有边界基础。

\textbf{状态复用。} 把对子树的重复扫描压成返回值或队列状态；每个节点只合并孩子给出的稳定信息。本题落点是：规律是后序返回是否找到目标或祖先，不能用 BST 的大小关系。手算时只改被新输入影响的那一小部分，而不是回头重算完整候选。

\textbf{指针和字段。} 状态字段要能说清：节点指针、传入约束或返回值、当前答案、层号或路径状态；空节点的返回值必须和状态定义匹配。手算时按这个协议跟踪：拿普通二叉树根 3，p=5、q=1 手算，左右子树分别找到一个目标，根 3 就是最近公共祖先；若两个目标都在左子树，答案应由左子树返回。然后按协议复查：手算时从叶子开始写返回值，或者从根开始写传入约束。不要只写访问顺序，要写每条边上传递的信息。

\textbf{代码落点。} 递归版本先处理空节点；迭代版本的栈元素要带完整状态；全负值、重复值和极端深树要单独测。关键代码行必须能逐句翻译成“旧状态怎样变成新状态”，否则就只是模板。

\textbf{正确性。} 证明从子树归纳或层序不变量开始：孩子状态正确时父节点合并正确；若用队列，当前轮队列必须正好保存当前层节点。

\textbf{边界实验。} 先用目标不存在、p 等于 q、一个是另一个祖先、比较节点地址做反例检查。实验时覆盖目标不存在、p 等于 q、一个是另一个祖先、比较节点地址，尤其关注空节点、单节点、链式树和全负值。递归版本和显式栈版本可以互相对拍。

\textbf{复杂度变化。} 暴力在每个节点重复扫描子树或反复寻找层边界可能退化到平方级；后序汇总、显式栈或层序遍历让每个节点处理常数次，通常是线性时间。

\textbf{迁移追问。} 如果题目改成“多次 LCA 查询可预处理倍增祖先表”，先判断原状态是否仍然覆盖新答案集合，再决定是微调边界、改 value 语义，还是换算法。

\noindent\textbf{LeetCode 337 打家劫舍 III。}

\textbf{题目说明。} LeetCode 337 打家劫舍 III 的题面入口要先还原成具体任务：树上相邻节点不能同时选，每个节点需要偷和不偷两个状态。

\textbf{样例入口。} 拿树 [3, 2, 3, null, 3, null, 1] 手算，偷根 3 时不能偷两个孩子，只能接孙子；不偷根时可以分别取左右孩子的最优。

\textbf{暴力基线。} 慢版本在每个节点重新扫描子树或路径，先求出局部答案。重复扫描暴露了真正状态：每个节点只需要从孩子或父亲那里拿到少量信息。放到本题就是：先按“树上相邻节点不能同时选，每个节点需要偷和不偷两个状态”完整试慢版本；样例里已经能看到“规律是每个节点返回两个状态：偷当前和不偷当前，父节点再组合”，所以优化前必须先能说清慢版本枚举了哪些候选。

\textbf{暴力过程。} 拿树 [3, 2, 3, null, 3, null, 1] 手算，偷根 3 时不能偷两个孩子，只能接孙子；不偷根时可以分别取左右孩子的最优。

\textbf{重复工作。} 题面模型：树上相邻节点不能同时选，每个节点需要偷和不偷两个状态。暴力版的慢点要落到本题：慢版本会在树上重复确认“树上相邻节点不能同时选，每个节点需要偷和不偷两个状态”。样例规律是“规律是每个节点返回两个状态：偷当前和不偷当前，父节点再组合”，说明父子之间真正需要传递的是高度、上下界、命中状态、层号或两类选择值，而不是反复扫描整棵子树。后面的优化只消掉这类重复工作，不能改变答案集合。

\textbf{优化条件。} 本题的关键观察是：后序汇总：偷当前则孩子不能偷，不偷当前则孩子可偷可不偷。这句话必须能翻译成可维护状态；如果题目条件改变后观察不再成立，就不能继续套原来的写法。

\textbf{相邻状态。} 规律是每个节点返回两个状态：偷当前和不偷当前，父节点再组合。把这个特例继续拆开看：节点向父亲返回的量和全局答案通常不是同一个量。先写叶子或空节点的返回值，再写父节点如何合并左右孩子，递归式才有边界基础。

\textbf{状态复用。} 把对子树的重复扫描压成返回值或队列状态；每个节点只合并孩子给出的稳定信息。本题落点是：规律是每个节点返回两个状态：偷当前和不偷当前，父节点再组合。手算时只改被新输入影响的那一小部分，而不是回头重算完整候选。

\textbf{指针和字段。} 状态字段要能说清：节点指针、传入约束或返回值、当前答案、层号或路径状态；空节点的返回值必须和状态定义匹配。手算时按这个协议跟踪：拿树 [3, 2, 3, null, 3, null, 1] 手算，偷根 3 时不能偷两个孩子，只能接孙子；不偷根时可以分别取左右孩子的最优。然后按协议复查：手算时从叶子开始写返回值，或者从根开始写传入约束。不要只写访问顺序，要写每条边上传递的信息。

\textbf{代码落点。} 递归版本先处理空节点；迭代版本的栈元素要带完整状态；全负值、重复值和极端深树要单独测。关键代码行必须能逐句翻译成“旧状态怎样变成新状态”，否则就只是模板。

\textbf{正确性。} 证明从子树归纳或层序不变量开始：孩子状态正确时父节点合并正确；若用队列，当前轮队列必须正好保存当前层节点。

\textbf{边界实验。} 先用空树、负值是否可能、链式树、递归深度做反例检查。实验时覆盖空树、负值是否可能、链式树、递归深度，尤其关注空节点、单节点、链式树和全负值。递归版本和显式栈版本可以互相对拍。

\textbf{复杂度变化。} 暴力在每个节点重复扫描子树或反复寻找层边界可能退化到平方级；后序汇总、显式栈或层序遍历让每个节点处理常数次，通常是线性时间。

\textbf{迁移追问。} 如果题目改成“可用显式栈后序保存节点状态，避免调用栈”，先判断原状态是否仍然覆盖新答案集合，再决定是微调边界、改 value 语义，还是换算法。

\noindent\textbf{LeetCode 450 删除二叉搜索树中的节点。}

\textbf{题目说明。} LeetCode 450 删除二叉搜索树中的节点 的题面入口要先还原成具体任务：删除节点后仍要保持 BST 中序有序。

\textbf{样例入口。} 拿 BST 删除 3 且 3 有左右孩子手算，不能直接丢掉整棵子树；可用右子树最小节点接替 3，再在右子树删除这个后继。

\textbf{暴力基线。} 慢版本在每个节点重新扫描子树或路径，先求出局部答案。重复扫描暴露了真正状态：每个节点只需要从孩子或父亲那里拿到少量信息。放到本题就是：先按“删除节点后仍要保持 BST 中序有序”完整试慢版本；样例里已经能看到“规律是删除分三类：叶子、单孩子、双孩子，双孩子要用前驱或后继保持 BST 顺序”，所以优化前必须先能说清慢版本枚举了哪些候选。

\textbf{暴力过程。} 拿 BST 删除 3 且 3 有左右孩子手算，不能直接丢掉整棵子树；可用右子树最小节点接替 3，再在右子树删除这个后继。

\textbf{重复工作。} 题面模型：删除节点后仍要保持 BST 中序有序。暴力版的慢点要落到本题：慢版本会在树上重复确认“删除节点后仍要保持 BST 中序有序”。样例规律是“规律是删除分三类：叶子、单孩子、双孩子，双孩子要用前驱或后继保持 BST 顺序”，说明父子之间真正需要传递的是高度、上下界、命中状态、层号或两类选择值，而不是反复扫描整棵子树。后面的优化只消掉这类重复工作，不能改变答案集合。

\textbf{优化条件。} 本题的关键观察是：若有两个孩子，用后继或前驱替换当前值，再删除那个后继节点。这句话必须能翻译成可维护状态；如果题目条件改变后观察不再成立，就不能继续套原来的写法。

\textbf{相邻状态。} 规律是删除分三类：叶子、单孩子、双孩子，双孩子要用前驱或后继保持 BST 顺序。把这个特例继续拆开看：节点向父亲返回的量和全局答案通常不是同一个量。先写叶子或空节点的返回值，再写父节点如何合并左右孩子，递归式才有边界基础。

\textbf{状态复用。} 把对子树的重复扫描压成返回值或队列状态；每个节点只合并孩子给出的稳定信息。本题落点是：规律是删除分三类：叶子、单孩子、双孩子，双孩子要用前驱或后继保持 BST 顺序。手算时只改被新输入影响的那一小部分，而不是回头重算完整候选。

\textbf{指针和字段。} 状态字段要能说清：节点指针、传入约束或返回值、当前答案、层号或路径状态；空节点的返回值必须和状态定义匹配。手算时按这个协议跟踪：拿 BST 删除 3 且 3 有左右孩子手算，不能直接丢掉整棵子树；可用右子树最小节点接替 3，再在右子树删除这个后继。然后按协议复查：手算时从叶子开始写返回值，或者从根开始写传入约束。不要只写访问顺序，要写每条边上传递的信息。

\textbf{代码落点。} 递归版本先处理空节点；迭代版本的栈元素要带完整状态；全负值、重复值和极端深树要单独测。关键代码行必须能逐句翻译成“旧状态怎样变成新状态”，否则就只是模板。

\textbf{正确性。} 证明从子树归纳或层序不变量开始：孩子状态正确时父节点合并正确；若用队列，当前轮队列必须正好保存当前层节点。

\textbf{边界实验。} 先用删除叶子、一个孩子、两个孩子、删除根节点做反例检查。实验时覆盖删除叶子、一个孩子、两个孩子、删除根节点，尤其关注空节点、单节点、链式树和全负值。递归版本和显式栈版本可以互相对拍。

\textbf{复杂度变化。} 暴力在每个节点重复扫描子树或反复寻找层边界可能退化到平方级；后序汇总、显式栈或层序遍历让每个节点处理常数次，通常是线性时间。

\textbf{迁移追问。} 如果题目改成“工程实现要清楚是移动值还是重连节点”，先判断原状态是否仍然覆盖新答案集合，再决定是微调边界、改 value 语义，还是换算法。

\noindent\textbf{LeetCode 543 二叉树直径。}

\textbf{题目说明。} LeetCode 543 二叉树直径 的题面入口要先还原成具体任务：直径经过某节点时等于左高加右高，但向父亲只能贡献单边高度。

\textbf{样例入口。} 拿树 1 的左右子树高度分别为 2 和 1 手算，经过根的直径是 3 条边；但向父节点只能返回单侧高度。

\textbf{暴力基线。} 慢版本在每个节点重新扫描子树或路径，先求出局部答案。重复扫描暴露了真正状态：每个节点只需要从孩子或父亲那里拿到少量信息。放到本题就是：先按“直径经过某节点时等于左高加右高，但向父亲只能贡献单边高度”完整试慢版本；样例里已经能看到“规律是后序返回高度，全局答案用 left\_height+right\_height 更新”，所以优化前必须先能说清慢版本枚举了哪些候选。

\textbf{暴力过程。} 拿树 1 的左右子树高度分别为 2 和 1 手算，经过根的直径是 3 条边；但向父节点只能返回单侧高度。

\textbf{重复工作。} 题面模型：直径经过某节点时等于左高加右高，但向父亲只能贡献单边高度。暴力版的慢点要落到本题：慢版本会在树上重复确认“直径经过某节点时等于左高加右高，但向父亲只能贡献单边高度”。样例规律是“规律是后序返回高度，全局答案用 left\_height+right\_height 更新”，说明父子之间真正需要传递的是高度、上下界、命中状态、层号或两类选择值，而不是反复扫描整棵子树。后面的优化只消掉这类重复工作，不能改变答案集合。

\textbf{优化条件。} 本题的关键观察是：后序遍历同时更新全局直径并返回高度。这句话必须能翻译成可维护状态；如果题目条件改变后观察不再成立，就不能继续套原来的写法。

\textbf{相邻状态。} 规律是后序返回高度，全局答案用 left\_height+right\_height 更新。把这个特例继续拆开看：节点向父亲返回的量和全局答案通常不是同一个量。先写叶子或空节点的返回值，再写父节点如何合并左右孩子，递归式才有边界基础。

\textbf{状态复用。} 把对子树的重复扫描压成返回值或队列状态；每个节点只合并孩子给出的稳定信息。本题落点是：规律是后序返回高度，全局答案用 left\_height+right\_height 更新。手算时只改被新输入影响的那一小部分，而不是回头重算完整候选。

\textbf{指针和字段。} 状态字段要能说清：节点指针、传入约束或返回值、当前答案、层号或路径状态；空节点的返回值必须和状态定义匹配。手算时按这个协议跟踪：拿树 1 的左右子树高度分别为 2 和 1 手算，经过根的直径是 3 条边；但向父节点只能返回单侧高度。然后按协议复查：手算时从叶子开始写返回值，或者从根开始写传入约束。不要只写访问顺序，要写每条边上传递的信息。

\textbf{代码落点。} 递归版本先处理空节点；迭代版本的栈元素要带完整状态；全负值、重复值和极端深树要单独测。关键代码行必须能逐句翻译成“旧状态怎样变成新状态”，否则就只是模板。

\textbf{正确性。} 证明从子树归纳或层序不变量开始：孩子状态正确时父节点合并正确；若用队列，当前轮队列必须正好保存当前层节点。

\textbf{边界实验。} 先用空树、单节点、链式树、边数还是节点数定义做反例检查。实验时覆盖空树、单节点、链式树、边数还是节点数定义，尤其关注空节点、单节点、链式树和全负值。递归版本和显式栈版本可以互相对拍。

\textbf{复杂度变化。} 暴力在每个节点重复扫描子树或反复寻找层边界可能退化到平方级；后序汇总、显式栈或层序遍历让每个节点处理常数次，通常是线性时间。

\textbf{迁移追问。} 如果题目改成“最大路径和是同型题，但贡献值和全局更新有负数处理”，先判断原状态是否仍然覆盖新答案集合，再决定是微调边界、改 value 语义，还是换算法。

\noindent\textbf{LeetCode 834 树中距离之和。}

\textbf{题目说明。} LeetCode 834 树中距离之和 的题面入口要先还原成具体任务：每个节点作为根的距离和可以由相邻换根关系转移。

\textbf{样例入口。} 拿 n=6 的样例树手算，以 0 为根先算子树大小和距离和；把根从 0 换到 2 时，2 子树内 4 个点距离都减一，外部 2 个点距离都加一。

\textbf{暴力基线。} 慢版本在每个节点重新扫描子树或路径，先求出局部答案。重复扫描暴露了真正状态：每个节点只需要从孩子或父亲那里拿到少量信息。放到本题就是：先按“每个节点作为根的距离和可以由相邻换根关系转移”完整试慢版本；样例里已经能看到“规律是先求一个根的答案，再沿边按“子树内减少、子树外增加”的关系换根”，所以优化前必须先能说清慢版本枚举了哪些候选。

\textbf{暴力过程。} 拿 n=6 的样例树手算，以 0 为根先算子树大小和距离和；把根从 0 换到 2 时，2 子树内 4 个点距离都减一，外部 2 个点距离都加一。

\textbf{重复工作。} 题面模型：每个节点作为根的距离和可以由相邻换根关系转移。暴力版的慢点要落到本题：慢版本会在树上重复确认“每个节点作为根的距离和可以由相邻换根关系转移”。样例规律是“规律是先求一个根的答案，再沿边按“子树内减少、子树外增加”的关系换根”，说明父子之间真正需要传递的是高度、上下界、命中状态、层号或两类选择值，而不是反复扫描整棵子树。后面的优化只消掉这类重复工作，不能改变答案集合。

\textbf{优化条件。} 本题的关键观察是：先后序计算子树大小和根节点答案，再从父到子换根：子树内距离减一，子树外距离加一。这句话必须能翻译成可维护状态；如果题目条件改变后观察不再成立，就不能继续套原来的写法。

\textbf{相邻状态。} 规律是先求一个根的答案，再沿边按“子树内减少、子树外增加”的关系换根。把这个特例继续拆开看：节点向父亲返回的量和全局答案通常不是同一个量。先写叶子或空节点的返回值，再写父节点如何合并左右孩子，递归式才有边界基础。

\textbf{状态复用。} 把对子树的重复扫描压成返回值或队列状态；每个节点只合并孩子给出的稳定信息。本题落点是：规律是先求一个根的答案，再沿边按“子树内减少、子树外增加”的关系换根。手算时只改被新输入影响的那一小部分，而不是回头重算完整候选。

\textbf{指针和字段。} 状态字段要能说清：节点指针、传入约束或返回值、当前答案、层号或路径状态；空节点的返回值必须和状态定义匹配。手算时按这个协议跟踪：拿 n=6 的样例树手算，以 0 为根先算子树大小和距离和；把根从 0 换到 2 时，2 子树内 4 个点距离都减一，外部 2 个点距离都加一。然后按协议复查：手算时从叶子开始写返回值，或者从根开始写传入约束。不要只写访问顺序，要写每条边上传递的信息。

\textbf{代码落点。} 递归版本先处理空节点；迭代版本的栈元素要带完整状态；全负值、重复值和极端深树要单独测。关键代码行必须能逐句翻译成“旧状态怎样变成新状态”，否则就只是模板。

\textbf{正确性。} 证明从子树归纳或层序不变量开始：孩子状态正确时父节点合并正确；若用队列，当前轮队列必须正好保存当前层节点。

\textbf{边界实验。} 先用单节点、链式树、星形树、n 很大递归栈、边是无向需建双向邻接做反例检查。实验时覆盖单节点、链式树、星形树、n 很大递归栈、边是无向需建双向邻接，尤其关注空节点、单节点、链式树和全负值。递归版本和显式栈版本可以互相对拍。

\textbf{复杂度变化。} 暴力在每个节点重复扫描子树或反复寻找层边界可能退化到平方级；后序汇总、显式栈或层序遍历让每个节点处理常数次，通常是线性时间。

\textbf{迁移追问。} 如果题目改成“若边有权，换根公式要把 size 项乘以边权”，先判断原状态是否仍然覆盖新答案集合，再决定是微调边界、改 value 语义，还是换算法。

\noindent\textbf{LeetCode 1483 树节点的第 K 个祖先。}

\textbf{题目说明。} LeetCode 1483 树节点的第 K 个祖先 的题面入口要先还原成具体任务：第 k 个祖先可以按 k 的二进制拆成若干 \ensuremath{2^{j}} 级跳跃。

\textbf{样例入口。} 拿一棵树中查询 node=5 的第 6 个祖先手算，6 的二进制是 110，可以先跳 2 级再跳 4 级，而不是一步步向上走 6 次。

\textbf{暴力基线。} 慢版本在每个节点重新扫描子树或路径，先求出局部答案。重复扫描暴露了真正状态：每个节点只需要从孩子或父亲那里拿到少量信息。放到本题就是：先按“第 k 个祖先可以按 k 的二进制拆成若干 \ensuremath{2^{j}} 级跳跃”完整试慢版本；样例里已经能看到“规律是预处理 up[node][j] 表示 \ensuremath{2^{j}} 祖先，查询时按 k 的二进制位跳跃”，所以优化前必须先能说清慢版本枚举了哪些候选。

\textbf{暴力过程。} 拿一棵树中查询 node=5 的第 6 个祖先手算，6 的二进制是 110，可以先跳 2 级再跳 4 级，而不是一步步向上走 6 次。

\textbf{重复工作。} 题面模型：第 k 个祖先可以按 k 的二进制拆成若干 \ensuremath{2^{j}} 级跳跃。暴力版的慢点要落到本题：慢版本会在树上重复确认“第 k 个祖先可以按 k 的二进制拆成若干 \ensuremath{2^{j}} 级跳跃”。样例规律是“规律是预处理 up[node][j] 表示 \ensuremath{2^{j}} 祖先，查询时按 k 的二进制位跳跃”，说明父子之间真正需要传递的是高度、上下界、命中状态、层号或两类选择值，而不是反复扫描整棵子树。后面的优化只消掉这类重复工作，不能改变答案集合。

\textbf{优化条件。} 本题的关键观察是：预处理 up[node][j] 表示 node 的 \ensuremath{2^{j}} 祖先，查询时按 k 的二进制位向上跳。这句话必须能翻译成可维护状态；如果题目条件改变后观察不再成立，就不能继续套原来的写法。

\textbf{相邻状态。} 规律是预处理 up[node][j] 表示 \ensuremath{2^{j}} 祖先，查询时按 k 的二进制位跳跃。把这个特例继续拆开看：节点向父亲返回的量和全局答案通常不是同一个量。先写叶子或空节点的返回值，再写父节点如何合并左右孩子，递归式才有边界基础。

\textbf{状态复用。} 把对子树的重复扫描压成返回值或队列状态；每个节点只合并孩子给出的稳定信息。本题落点是：规律是预处理 up[node][j] 表示 \ensuremath{2^{j}} 祖先，查询时按 k 的二进制位跳跃。手算时只改被新输入影响的那一小部分，而不是回头重算完整候选。

\textbf{指针和字段。} 状态字段要能说清：节点指针、传入约束或返回值、当前答案、层号或路径状态；空节点的返回值必须和状态定义匹配。手算时按这个协议跟踪：拿一棵树中查询 node=5 的第 6 个祖先手算，6 的二进制是 110，可以先跳 2 级再跳 4 级，而不是一步步向上走 6 次。然后按协议复查：手算时从叶子开始写返回值，或者从根开始写传入约束。不要只写访问顺序，要写每条边上传递的信息。

\textbf{代码落点。} 递归版本先处理空节点；迭代版本的栈元素要带完整状态；全负值、重复值和极端深树要单独测。关键代码行必须能逐句翻译成“旧状态怎样变成新状态”，否则就只是模板。

\textbf{正确性。} 证明从子树归纳或层序不变量开始：孩子状态正确时父节点合并正确；若用队列，当前轮队列必须正好保存当前层节点。

\textbf{边界实验。} 先用根节点祖先不存在、k 为零、树很深、查询很多、表大小做反例检查。实验时覆盖根节点祖先不存在、k 为零、树很深、查询很多、表大小，尤其关注空节点、单节点、链式树和全负值。递归版本和显式栈版本可以互相对拍。

\textbf{复杂度变化。} 暴力在每个节点重复扫描子树或反复寻找层边界可能退化到平方级；后序汇总、显式栈或层序遍历让每个节点处理常数次，通常是线性时间。

\textbf{迁移追问。} 如果题目改成“LCA 也使用倍增表，只是同时提升两个节点”，先判断原状态是否仍然覆盖新答案集合，再决定是微调边界、改 value 语义，还是换算法。

\subsection{实验复盘}

追问通常会问返回值语义、空节点初值、全负数和极端深树。请准备把递归思想改写成显式栈或队列的口头方案。

复盘本节时先从 LeetCode 94 二叉树中序遍历、LeetCode 98 验证二叉搜索树、LeetCode 102 二叉树层序遍历 开始：每道题都先手写一个能覆盖全部候选的暴力版，再指出优化结构具体保存了哪一段历史信息，最后用相邻案例比较为什么不能互换解法。

\topic{图建模、BFS 与 DFS}
这一节先看 LeetCode 126 单词接龙 II、LeetCode 127 单词接龙、LeetCode 130 被围绕的区域 这几道 LeetCode 题，再整理同一题型的分流和推导。阅读顺序不是先背原理，而是先看案例的暴力瓶颈，再把重复出现的观察抽象成方法。

\subsection{原理和适用条件}

以 LeetCode 126 单词接龙 II、LeetCode 127 单词接龙、LeetCode 130 被围绕的区域 为主线：LeetCode 126 单词接龙 II：模型是“不仅要最短转换长度，还要输出所有最短转换路径。”，优化抓手是“BFS 只建立最短层级和前驱列表，同一层的多个前驱都保留，最后再从终点回溯路径。”；LeetCode 127 单词接龙：模型是“每个单词是节点，一次修改一个字符形成边。”，优化抓手是“通配模式桶加速邻居查询，BFS 第一次到达终点就是最短转换长度。”；LeetCode 130 被围绕的区域：模型是“只有不连通到边界的 O 才会被翻转。”，优化抓手是“从边界 O 出发标记安全区域，再翻转剩余 O。”。这些案例共同推出的规律是：先定义节点、边和状态，再决定 BFS、DFS、拓扑或并查集。图题失败常常不是搜索模板错，而是状态维度不完整或邻居生成过慢。 方法名只是复盘后的标签，真正要掌握的是从题面约束、暴力失败、状态设计到边界反例的推导链。

\begin{principlebox}{图示：从 LeetCode 案例到 图建模、BFS 与 DFS推导链}
\begin{tabular}{@{}p{0.18\linewidth}p{0.74\linewidth}@{}}
暴力基线 & 枚举候选、完整模拟或逐个检查，先保证覆盖全部可能答案。\\
瓶颈定位 & 慢点来自重复访问和邻居查询过慢。若节点很多但边很少，用邻接矩阵会浪费空间；若每次找邻居都扫描全体节点，BFS 本身再正确也会超时。\\
结构选择 & 无权最短步数用 BFS，连通块和状态检测用 DFS 或显式栈，有向无环依赖用拓扑排序。多源问题把所有源同时入队，状态带资源的问题不能只按位置去重。\\
不变量 & 访问标记通常在入队或入栈时设置，保证一个状态只被加入一次。BFS 队列按距离递增；拓扑排序队列保存当前入度为零的节点。\\
代码落地 & 图节点连续时用 \code{std::vector<std::vector<int>>}；节点是字符串时先编号或用哈希表。网格方向数组用 \code{constexpr std::array}，边界判断集中写。\\
\end{tabular}
\end{principlebox}

\subsection{LeetCode 案例分流}

\begin{principlebox}{从 LeetCode 案例分流：图建模、BFS 与 DFS}
\textbf{连通块。}
代表案例：LeetCode 200 岛屿数量、LeetCode 695 岛屿最大面积、LeetCode 130 被围绕区域、LeetCode 417 太平洋大西洋水流。先看这些题的题面条件：节点之间只有可达关系，答案是块数量、块面积或安全区域。由此推出的处理方式是：扫描未访问节点后启动 BFS/DFS 标记整块。风险点：边界反向搜索常比正向判断更简单。

\textbf{无权最短路。}
代表案例：LeetCode 127 单词接龙、LeetCode 752 打开转盘锁、LeetCode 994 腐烂橘子、LeetCode 1091 二进制矩阵最短路径。先看这些题的题面条件：每条边代价相同，问最少步数或最短转换长度。由此推出的处理方式是：BFS 按层扩展，第一次到达即最短。风险点：访问标记应覆盖完整状态维度。

\textbf{依赖和拓扑。}
代表案例：LeetCode 207 课程表、LeetCode 210 课程表 II、LeetCode 1462 课程表 IV。先看这些题的题面条件：有向边表达先后关系，需要判环或输出顺序。由此推出的处理方式是：入度队列或三色 DFS 判断有向环。风险点：边方向会影响输出顺序和可达预处理。

\textbf{状态图。}
代表案例：LeetCode 864 获取所有钥匙最短路径、LeetCode 1293 网格障碍消除最短路。先看这些题的题面条件：位置之外还带钥匙、资源、访问集合或时间。由此推出的处理方式是：把附加资源编码进状态，visited 不能只按位置。风险点：状态维度不完整会错误剪枝。

\end{principlebox}

\subsection{案例矩阵}

本节案例覆盖 LeetCode 126 单词接龙 II、LeetCode 127 单词接龙、LeetCode 130 被围绕的区域、LeetCode 133 克隆图、LeetCode 200 岛屿数量、LeetCode 207 课程表、LeetCode 210 课程表 II、LeetCode 417 太平洋大西洋水流问题、LeetCode 695 岛屿的最大面积、LeetCode 752 打开转盘锁、LeetCode 785 判断二分图、LeetCode 815 公交路线、LeetCode 847 访问所有节点的最短路径、LeetCode 864 获取所有钥匙的最短路径、LeetCode 886 可能的二分法、LeetCode 994 腐烂的橘子、LeetCode 1091 二进制矩阵中的最短路径、LeetCode 1192 查找集群内的关键连接、LeetCode 1293 网格中的最短路径、LeetCode 1462 课程表 IV、LeetCode 1584 连接所有点的最小费用、LeetCode 1820 最多邀请的个数。阅读时不要按题号背答案，而要先判断题面落在哪个分流场景：查询目标是什么，输入约束给了什么单调性或等价关系，暴力慢在哪里，优化结构保存了什么状态。

\begin{principlebox}{案例矩阵}
\textbf{LeetCode 126 单词接龙 II。}
不仅要最短转换长度，还要输出所有最短转换路径。单点风险：终点不在词表、同层多父节点、本层访问不能提前删除等长前驱、输出规模可能很大。

\textbf{LeetCode 127 单词接龙。}
每个单词是节点，一次修改一个字符形成边。单点风险：终点不在词表、起点和终点相同、桶重复扫描、单词长度一致。

\textbf{LeetCode 130 被围绕的区域。}
只有不连通到边界的 O 才会被翻转。单点风险：全边界、无 O、单行单列、内部岛屿。

\textbf{LeetCode 133 克隆图。}
每个原节点需要唯一对应一个新节点，边按原图连接。单点风险：空图、自环、重边、环结构、不能按节点值唯一假设。

\textbf{LeetCode 200 岛屿数量。}
陆地格子是节点，四方向相邻是边，岛屿是连通块。单点风险：空网格、全水、全陆、斜对角不连通、是否允许修改输入。

\textbf{LeetCode 207 课程表。}
课程是节点，先修关系是有向边，问题是判断是否有环。单点风险：边方向、孤立课程、环、自依赖、重复边。

\textbf{LeetCode 210 课程表 II。}
在判环基础上还要输出一个可行拓扑序。单点风险：有环返回空数组、多个可行顺序、边方向。

\textbf{LeetCode 417 太平洋大西洋水流问题。}
从每个格子向海搜索很慢，反向从海边沿非递减高度搜索。单点风险：单行单列、平台高度、边界重复入队、方向条件反向。

\textbf{LeetCode 695 岛屿的最大面积。}
面积就是一个连通块中陆地格子的数量。单点风险：全水、全陆、细长岛、是否允许修改输入。

\textbf{LeetCode 752 打开转盘锁。}
四位密码是状态，每次转动一位形成边。单点风险：起点死亡、目标就是起点、数字环绕、死亡集合很大。

\textbf{LeetCode 785 判断二分图。}
二分图要求每条边两端颜色相反。单点风险：非连通图、自环、奇环、孤立点、重复边。

\textbf{LeetCode 815 公交路线。}
公交路线和站点构成二分关系，换乘次数比站点步数更重要。单点风险：起点等于终点、某站经过路线很多、路线重复访问、目标站不可达。

\textbf{LeetCode 847 访问所有节点的最短路径。}
访问集合影响未来能力，同一节点配不同 mask 不是同一状态。单点风险：单节点、完全图、星形图、重复经过节点、visited 必须包含 mask。

\textbf{LeetCode 864 获取所有钥匙的最短路径。}
位置相同但钥匙集合不同，未来能开的门不同。单点风险：起点旁边有门、钥匙顺序、不可达、钥匙数量决定掩码大小。

\textbf{LeetCode 886 可能的二分法。}
互相讨厌的人必须分到不同组，本质是图二染色。单点风险：孤立点、非连通图、奇环、重复 dislike、自身 dislike。

\textbf{LeetCode 994 腐烂的橘子。}
所有初始腐烂橘子同时作为 BFS 源点。单点风险：没有新鲜橘子、没有腐烂源、隔离新鲜橘子、空格阻断。

\textbf{LeetCode 1091 二进制矩阵中的最短路径。}
八方向网格无权边，答案是从左上到右下的最少格子数。单点风险：起点或终点阻塞、单格、八方向越界、无路可达。

\textbf{LeetCode 1192 查找集群内的关键连接。}
无向图中的桥是删除后会断开连通性的边。单点风险：重边、根节点、图保证连通、无向边需要编号跳过反向边、递归栈。

\textbf{LeetCode 1293 网格中的最短路径。}
走到同一格时，剩余可消除障碍数量不同会影响未来可达性。单点风险：起点终点、k 很大、障碍密集、二维 visited 错误剪枝。

\textbf{LeetCode 1462 课程表 IV。}
大量先修关系查询需要预处理课程之间的可达性。单点风险：课程之间没有关系、自身是否算先修、查询很多、图中按题意应为 DAG。

\textbf{LeetCode 1584 连接所有点的最小费用。}
完全图上连接所有点的最小总成本是最小生成树问题，边权为曼哈顿距离。单点风险：单点、重复坐标、完全图边数、曼哈顿距离、Prim 更新外部点最小连边。

\textbf{LeetCode 1820 最多邀请的个数。}
男生和女生构成二分图，目标是最大一对一匹配数量。单点风险：某一侧为空、多个男生争同一女生、贪心失败、访问标记按轮重置、稀疏边。

\end{principlebox}

\subsection{共性落点}

\begin{principlebox}{本组共性落点}
\textbf{慢版本。} 暴力做法常从每个节点重新搜索，或者在图中不加访问标记地反复扩展状态。图题第一步不是选 DFS 还是 BFS，而是明确节点和边。
\textbf{瓶颈。} 慢点来自重复访问和邻居查询过慢。若节点很多但边很少，用邻接矩阵会浪费空间；若每次找邻居都扫描全体节点，BFS 本身再正确也会超时。
\textbf{状态字段。} 状态编号、邻接生成方式、访问标记、距离或层数、队列内容；若状态含资源，visited 不能只按位置记录。
\textbf{不变量。} 访问标记通常在入队或入栈时设置，保证一个状态只被加入一次。BFS 队列按距离递增；拓扑排序队列保存当前入度为零的节点。
\textbf{实现检查。} 入队时标记还是出队时标记必须一致；方向数组集中定义；多源 BFS 要一次性把所有源点放入初始队列。
\textbf{C++ 落地。} 图节点连续时用 \code{std::vector<std::vector<int>>}；节点是字符串时先编号或用哈希表。网格方向数组用 \code{constexpr std::array}，边界判断集中写。
\textbf{证明抓手。} 证明从可达性开始：搜索覆盖所有合法边能到达的状态，访问标记只删除重复状态而不删除不同未来能力。
\end{principlebox}

\subsection{案例推导}

\noindent\textbf{LeetCode 126 单词接龙 II。}

\textbf{题目说明。} LeetCode 126 单词接龙 II 的题面入口要先还原成具体任务：不仅要最短转换长度，还要输出所有最短转换路径。

\textbf{样例入口。} 拿 hit 到 cog 的小词表手算，同一层 hot 可以通向 dot 和 lot，后续 dog/log 都可能到 cog。若某个单词在本层第一次见到就立刻删掉，会丢掉另一条同长度父路径。

\textbf{暴力基线。} 慢版本先按题意画出完整状态图，用普通搜索跑通可达性。这个过程用于确认不仅要最短转换长度，还要输出所有最短转换路径的节点和边，之后再讨论访问标记、层次或代价顺序。放到本题就是：先按“不仅要最短转换长度，还要输出所有最短转换路径”完整试慢版本；样例里已经能看到“规律是 BFS 建最短层级，前驱列表保留同层多父节点，最后再回溯所有路径”，所以优化前必须先能说清慢版本枚举了哪些候选。

\textbf{暴力过程。} 拿 hit 到 cog 的小词表手算，同一层 hot 可以通向 dot 和 lot，后续 dog/log 都可能到 cog。若某个单词在本层第一次见到就立刻删掉，会丢掉另一条同长度父路径。

\textbf{重复工作。} 题面模型：不仅要最短转换长度，还要输出所有最短转换路径。暴力版的慢点要落到本题：慢版本会围绕“不仅要最短转换长度，还要输出所有最短转换路径”反复展开同一批状态。样例规律是“规律是 BFS 建最短层级，前驱列表保留同层多父节点，最后再回溯所有路径”，说明必须把节点、边、访问标记和资源维度一次建清；同一状态不应重复入队，不同资源状态也不能误合并。后面的优化只消掉这类重复工作，不能改变答案集合。

\textbf{优化条件。} 本题的关键观察是：BFS 只建立最短层级和前驱列表，同一层的多个前驱都保留，最后再从终点回溯路径。这句话必须能翻译成可维护状态；如果题目条件改变后观察不再成立，就不能继续套原来的写法。

\textbf{相邻状态。} 规律是 BFS 建最短层级，前驱列表保留同层多父节点，最后再回溯所有路径。把这个特例继续拆开看：状态节点不一定等于原始位置，还可能包含钥匙、剩余资源、时间或访问集合。visited 只能合并未来能力完全相同的状态，否则会把合法路径剪掉。

\textbf{状态复用。} 把重复搜索压成访问标记、距离数组或拓扑状态；邻居生成也要从全量扫描变成结构化枚举。本题落点是：规律是 BFS 建最短层级，前驱列表保留同层多父节点，最后再回溯所有路径。手算时只改被新输入影响的那一小部分，而不是回头重算完整候选。

\textbf{指针和字段。} 状态字段要能说清：状态编号、邻接生成方式、访问标记、距离或层数、队列内容；若状态含资源，visited 不能只按位置记录。手算时按这个协议跟踪：拿 hit 到 cog 的小词表手算，同一层 hot 可以通向 dot 和 lot，后续 dog/log 都可能到 cog。若某个单词在本层第一次见到就立刻删掉，会丢掉另一条同长度父路径。然后按协议复查：手算时画队列或栈，状态必须包含题目需要的所有资源。入队时标记还是出队时标记，要和去重语义一致。

\textbf{代码落点。} 入队时标记还是出队时标记必须一致；方向数组集中定义；多源 BFS 要一次性把所有源点放入初始队列。关键代码行必须能逐句翻译成“旧状态怎样变成新状态”，否则就只是模板。

\textbf{正确性。} 证明从可达性开始：搜索覆盖所有合法边能到达的状态，访问标记只删除重复状态而不删除不同未来能力。

\textbf{边界实验。} 先用终点不在词表、同层多父节点、本层访问不能提前删除等长前驱、输出规模可能很大做反例检查。实验时覆盖终点不在词表、同层多父节点、本层访问不能提前删除等长前驱、输出规模可能很大，并打印入队时的完整状态。若同一位置但资源不同的状态被合并，很多图题会出现隐蔽错误。

\textbf{复杂度变化。} 暴力重复从多个起点搜索会反复遍历图；正确建模和访问标记后，BFS 或 DFS 通常按节点数加边数计算，状态图题则按完整状态数计算。

\textbf{迁移追问。} 如果题目改成“若只问最短长度，普通 BFS 或双向 BFS 更轻；若要路径，就必须保存前驱关系”，先判断原状态是否仍然覆盖新答案集合，再决定是微调边界、改 value 语义，还是换算法。

\noindent\textbf{LeetCode 127 单词接龙。}

\textbf{题目说明。} LeetCode 127 单词接龙 的题面入口要先还原成具体任务：每个单词是节点，一次修改一个字符形成边。

\textbf{样例入口。} 拿 hit -> hot -> dot -> dog -> cog 手算，每次只能改一个字符，所以 BFS 第一层是 hot，第二层才到 dot 或 lot。

\textbf{暴力基线。} 暴力可以把每个单词和所有其他单词比较一遍，差一个字符就连边，再从 beginWord 做 BFS。

\textbf{暴力过程。} 以 hit 到 cog 为例，hit 只能一步到 hot；hot 下一层到 dot 或 lot。第一次到 cog 的层数就是最短转换长度。

\textbf{重复工作。} 题面模型：每个单词是节点，一次修改一个字符形成边。暴力版的慢点要落到本题：浪费在找邻居时每次扫描整个词表。单词长度固定时，可以把 h*t、*ot 这类通配模式作为桶，快速找到只差一位的邻居。后面的优化只消掉这类重复工作，不能改变答案集合。

\textbf{优化条件。} 本题的关键观察是：通配模式桶加速邻居查询，BFS 第一次到达终点就是最短转换长度。这句话必须能翻译成可维护状态；如果题目条件改变后观察不再成立，就不能继续套原来的写法。

\textbf{相邻状态。} 规律是单词作为节点，通配桶如 h*t 把邻居快速找出；第一次到达终点就是最短转换长度。把这个特例继续拆开看：状态节点不一定等于原始位置，还可能包含钥匙、剩余资源、时间或访问集合。visited 只能合并未来能力完全相同的状态，否则会把合法路径剪掉。

\textbf{状态复用。} 预处理 pattern -> words 的映射，BFS 队列保存当前单词和步数，visited 防止重复入队。手算时只改被新输入影响的那一小部分，而不是回头重算完整候选。

\textbf{指针和字段。} 状态字段要能说清：word 是当前单词，pattern 是替换某一位后的通配键，distance 是转换层数。手算时按这个协议跟踪：hit 生成 *it、h*t、hi*，其中 h*t 桶找到 hot；hot 再通过 *ot 找到 dot 和 lot。

\textbf{代码落点。} 关键是访问标记在入队时设置。桶用完后可以清空，避免不同节点重复展开同一批邻居。关键代码行必须能逐句翻译成“旧状态怎样变成新状态”，否则就只是模板。

\textbf{正确性。} 每条转换边权都为 1，BFS 按层扩展；第一次到达 endWord 时经过的边数最少。

\textbf{边界实验。} 先用终点不在词表、起点和终点相同、桶重复扫描、单词长度一致做反例检查。实验时覆盖终点不在词表、起点和终点相同、桶重复扫描、单词长度一致，并打印入队时的完整状态。若同一位置但资源不同的状态被合并，很多图题会出现隐蔽错误。

\textbf{复杂度变化。} 两两建图 O(\ensuremath{N^{2}}L)，通配桶建图和 BFS 约 O(\ensuremath{NL^{2}}) 或按桶规模计，空间 O(NL)。

\textbf{迁移追问。} 如果题目改成“双向 BFS 可以进一步降低搜索空间”，先判断原状态是否仍然覆盖新答案集合，再决定是微调边界、改 value 语义，还是换算法。

\noindent\textbf{LeetCode 130 被围绕的区域。}

\textbf{题目说明。} LeetCode 130 被围绕的区域 的题面入口要先还原成具体任务：只有不连通到边界的 O 才会被翻转。

\textbf{样例入口。} 拿边界上有 O、内部也有 O 的棋盘手算，边界 O 以及和它连通的 O 不能翻转，真正要翻的是不接触边界的内部块。

\textbf{暴力基线。} 慢版本先按题意画出完整状态图，用普通搜索跑通可达性。这个过程用于确认只有不连通到边界的 O 才会被翻转的节点和边，之后再讨论访问标记、层次或代价顺序。放到本题就是：先按“只有不连通到边界的 O 才会被翻转”完整试慢版本；样例里已经能看到“规律是反向从边界 O 出发标记安全区，再翻转剩下的 O；单行单列全部接触边界”，所以优化前必须先能说清慢版本枚举了哪些候选。

\textbf{暴力过程。} 拿边界上有 O、内部也有 O 的棋盘手算，边界 O 以及和它连通的 O 不能翻转，真正要翻的是不接触边界的内部块。

\textbf{重复工作。} 题面模型：只有不连通到边界的 O 才会被翻转。暴力版的慢点要落到本题：慢版本会围绕“只有不连通到边界的 O 才会被翻转”反复展开同一批状态。样例规律是“规律是反向从边界 O 出发标记安全区，再翻转剩下的 O；单行单列全部接触边界”，说明必须把节点、边、访问标记和资源维度一次建清；同一状态不应重复入队，不同资源状态也不能误合并。后面的优化只消掉这类重复工作，不能改变答案集合。

\textbf{优化条件。} 本题的关键观察是：从边界 O 出发标记安全区域，再翻转剩余 O。这句话必须能翻译成可维护状态；如果题目条件改变后观察不再成立，就不能继续套原来的写法。

\textbf{相邻状态。} 规律是反向从边界 O 出发标记安全区，再翻转剩下的 O；单行单列全部接触边界。把这个特例继续拆开看：状态节点不一定等于原始位置，还可能包含钥匙、剩余资源、时间或访问集合。visited 只能合并未来能力完全相同的状态，否则会把合法路径剪掉。

\textbf{状态复用。} 把重复搜索压成访问标记、距离数组或拓扑状态；邻居生成也要从全量扫描变成结构化枚举。本题落点是：规律是反向从边界 O 出发标记安全区，再翻转剩下的 O；单行单列全部接触边界。手算时只改被新输入影响的那一小部分，而不是回头重算完整候选。

\textbf{指针和字段。} 状态字段要能说清：状态编号、邻接生成方式、访问标记、距离或层数、队列内容；若状态含资源，visited 不能只按位置记录。手算时按这个协议跟踪：拿边界上有 O、内部也有 O 的棋盘手算，边界 O 以及和它连通的 O 不能翻转，真正要翻的是不接触边界的内部块。然后按协议复查：手算时画队列或栈，状态必须包含题目需要的所有资源。入队时标记还是出队时标记，要和去重语义一致。

\textbf{代码落点。} 入队时标记还是出队时标记必须一致；方向数组集中定义；多源 BFS 要一次性把所有源点放入初始队列。关键代码行必须能逐句翻译成“旧状态怎样变成新状态”，否则就只是模板。

\textbf{正确性。} 证明从可达性开始：搜索覆盖所有合法边能到达的状态，访问标记只删除重复状态而不删除不同未来能力。

\textbf{边界实验。} 先用全边界、无 O、单行单列、内部岛屿做反例检查。实验时覆盖全边界、无 O、单行单列、内部岛屿，并打印入队时的完整状态。若同一位置但资源不同的状态被合并，很多图题会出现隐蔽错误。

\textbf{复杂度变化。} 暴力重复从多个起点搜索会反复遍历图；正确建模和访问标记后，BFS 或 DFS 通常按节点数加边数计算，状态图题则按完整状态数计算。

\textbf{迁移追问。} 如果题目改成“这是反向思维：找不能被翻转的区域比直接判断每块区域更简单”，先判断原状态是否仍然覆盖新答案集合，再决定是微调边界、改 value 语义，还是换算法。

\noindent\textbf{LeetCode 133 克隆图。}

\textbf{题目说明。} LeetCode 133 克隆图 的题面入口要先还原成具体任务：每个原节点需要唯一对应一个新节点，边按原图连接。

\textbf{样例入口。} 拿三角形 1-2-3-1 手算，克隆 1 时会遇到 2 和 3；克隆 2 时又会遇到 1，如果没有哈希表会无限复制。

\textbf{暴力基线。} 慢版本先按题意画出完整状态图，用普通搜索跑通可达性。这个过程用于确认每个原节点需要唯一对应一个新节点，边按原图连接的节点和边，之后再讨论访问标记、层次或代价顺序。放到本题就是：先按“每个原节点需要唯一对应一个新节点，边按原图连接”完整试慢版本；样例里已经能看到“规律是哈希表保存原节点到新节点的映射，先建节点再连邻居，环和自环都依赖这个映射”，所以优化前必须先能说清慢版本枚举了哪些候选。

\textbf{暴力过程。} 拿三角形 1-2-3-1 手算，克隆 1 时会遇到 2 和 3；克隆 2 时又会遇到 1，如果没有哈希表会无限复制。

\textbf{重复工作。} 题面模型：每个原节点需要唯一对应一个新节点，边按原图连接。暴力版的慢点要落到本题：慢版本会围绕“每个原节点需要唯一对应一个新节点，边按原图连接”反复展开同一批状态。样例规律是“规律是哈希表保存原节点到新节点的映射，先建节点再连邻居，环和自环都依赖这个映射”，说明必须把节点、边、访问标记和资源维度一次建清；同一状态不应重复入队，不同资源状态也不能误合并。后面的优化只消掉这类重复工作，不能改变答案集合。

\textbf{优化条件。} 本题的关键观察是：哈希表保存原节点到克隆节点的映射，BFS 或 DFS 遍历图。这句话必须能翻译成可维护状态；如果题目条件改变后观察不再成立，就不能继续套原来的写法。

\textbf{相邻状态。} 规律是哈希表保存原节点到新节点的映射，先建节点再连邻居，环和自环都依赖这个映射。把这个特例继续拆开看：状态节点不一定等于原始位置，还可能包含钥匙、剩余资源、时间或访问集合。visited 只能合并未来能力完全相同的状态，否则会把合法路径剪掉。

\textbf{状态复用。} 把重复搜索压成访问标记、距离数组或拓扑状态；邻居生成也要从全量扫描变成结构化枚举。本题落点是：规律是哈希表保存原节点到新节点的映射，先建节点再连邻居，环和自环都依赖这个映射。手算时只改被新输入影响的那一小部分，而不是回头重算完整候选。

\textbf{指针和字段。} 状态字段要能说清：状态编号、邻接生成方式、访问标记、距离或层数、队列内容；若状态含资源，visited 不能只按位置记录。手算时按这个协议跟踪：拿三角形 1-2-3-1 手算，克隆 1 时会遇到 2 和 3；克隆 2 时又会遇到 1，如果没有哈希表会无限复制。然后按协议复查：手算时画队列或栈，状态必须包含题目需要的所有资源。入队时标记还是出队时标记，要和去重语义一致。

\textbf{代码落点。} 入队时标记还是出队时标记必须一致；方向数组集中定义；多源 BFS 要一次性把所有源点放入初始队列。关键代码行必须能逐句翻译成“旧状态怎样变成新状态”，否则就只是模板。

\textbf{正确性。} 证明从可达性开始：搜索覆盖所有合法边能到达的状态，访问标记只删除重复状态而不删除不同未来能力。

\textbf{边界实验。} 先用空图、自环、重边、环结构、不能按节点值唯一假设做反例检查。实验时覆盖空图、自环、重边、环结构、不能按节点值唯一假设，并打印入队时的完整状态。若同一位置但资源不同的状态被合并，很多图题会出现隐蔽错误。

\textbf{复杂度变化。} 暴力重复从多个起点搜索会反复遍历图；正确建模和访问标记后，BFS 或 DFS 通常按节点数加边数计算，状态图题则按完整状态数计算。

\textbf{迁移追问。} 如果题目改成“若图节点值不唯一，必须用地址作为键”，先判断原状态是否仍然覆盖新答案集合，再决定是微调边界、改 value 语义，还是换算法。

\noindent\textbf{LeetCode 200 岛屿数量。}

\textbf{题目说明。} LeetCode 200 岛屿数量 的题面入口要先还原成具体任务：陆地格子是节点，四方向相邻是边，岛屿是连通块。

\textbf{样例入口。} 拿一个 2x2 小岛手算，从一个陆地出发能沿上下左右标记同一块岛。若不标记访问，会重复计数；若对角线也连通，会误合并。

\textbf{暴力基线。} 暴力可以对每个陆地格子单独搜索它所属的连通块，再用集合去重；这样同一块岛会被重复搜索很多次。

\textbf{暴力过程。} 以 2x2 全是陆地为例，从左上搜索会访问四个格子；如果之后从右上又重新搜索，就把同一座岛重复数了一遍。

\textbf{重复工作。} 题面模型：陆地格子是节点，四方向相邻是边，岛屿是连通块。暴力版的慢点要落到本题：浪费在已确认属于某个岛的陆地没有被标记。每块岛只需要在第一次遇到时完整淹没或标记。后面的优化只消掉这类重复工作，不能改变答案集合。

\textbf{优化条件。} 本题的关键观察是：扫描遇到未访问陆地就启动搜索并标记整块。这句话必须能翻译成可维护状态；如果题目条件改变后观察不再成立，就不能继续套原来的写法。

\textbf{相邻状态。} 规律是每发现一个未访问陆地，答案加一，并用 DFS/BFS 标记与它四连通的全部陆地。把这个特例继续拆开看：状态节点不一定等于原始位置，还可能包含钥匙、剩余资源、时间或访问集合。visited 只能合并未来能力完全相同的状态，否则会把合法路径剪掉。

\textbf{状态复用。} 扫描网格。遇到未访问陆地，答案加一，然后 BFS/DFS 把与它四方向连通的全部陆地标记为已访问。手算时只改被新输入影响的那一小部分，而不是回头重算完整候选。

\textbf{指针和字段。} 状态字段要能说清：row/col 是当前格子，visited 或原地改写表示已经归属某座岛，directions 是四方向邻居。手算时按这个协议跟踪：从左上陆地入队，依次扩展右边和下边陆地，最终整块连通区域都标记；扫描到这些格子时直接跳过。

\textbf{代码落点。} 关键是只看上下左右，不看对角线。边界判断集中写，访问标记在入队时设置，避免重复入队。关键代码行必须能逐句翻译成“旧状态怎样变成新状态”，否则就只是模板。

\textbf{正确性。} 一次搜索会访问且只访问与起点四连通的全部陆地；外层每次启动搜索都对应一块此前未计数的新岛。

\textbf{边界实验。} 先用空网格、全水、全陆、斜对角不连通、是否允许修改输入做反例检查。实验时覆盖空网格、全水、全陆、斜对角不连通、是否允许修改输入，并打印入队时的完整状态。若同一位置但资源不同的状态被合并，很多图题会出现隐蔽错误。

\textbf{复杂度变化。} 重复搜索最坏 O(\ensuremath{(mn)^{2}})，标记搜索 O(mn) 时间，空间 O(mn) 或递归/队列大小。

\textbf{迁移追问。} 如果题目改成“也可用并查集合并相邻陆地，适合动态岛屿扩展”，先判断原状态是否仍然覆盖新答案集合，再决定是微调边界、改 value 语义，还是换算法。

\noindent\textbf{LeetCode 207 课程表。}

\textbf{题目说明。} LeetCode 207 课程表 的题面入口要先还原成具体任务：课程是节点，先修关系是有向边，问题是判断是否有环。

\textbf{样例入口。} 拿课程 0->1 手算，入度为 0 的课程 0 先入队，学完后 1 的入度降为 0；若有 0->1 和 1->0，队列一开始为空。

\textbf{暴力基线。} 暴力可以不断扫描所有课程，找当前先修都已满足的课程；若一轮没有新课程可学，就判断有环。

\textbf{暴力过程。} 以 0->1 为例，课程 0 入度为 0，先学 0，随后 1 的入度变 0。若有 0->1 和 1->0，一开始没有入度为 0 的课。

\textbf{重复工作。} 题面模型：课程是节点，先修关系是有向边，问题是判断是否有环。暴力版的慢点要落到本题：浪费在每轮都重新检查所有课程的先修。入度正好表示还剩多少未满足依赖，学完一门课只影响它指向的后继。后面的优化只消掉这类重复工作，不能改变答案集合。

\textbf{优化条件。} 本题的关键观察是：入度拓扑排序不断学习当前无依赖课程。这句话必须能翻译成可维护状态；如果题目条件改变后观察不再成立，就不能继续套原来的写法。

\textbf{相邻状态。} 规律是拓扑排序用入度为 0 的节点推进，处理数量小于课程数说明有环。把这个特例继续拆开看：状态节点不一定等于原始位置，还可能包含钥匙、剩余资源、时间或访问集合。visited 只能合并未来能力完全相同的状态，否则会把合法路径剪掉。

\textbf{状态复用。} 构建邻接表和 indegree。把所有入度为 0 的课程入队，弹出课程时让后继入度减一，减到 0 就入队。手算时只改被新输入影响的那一小部分，而不是回头重算完整候选。

\textbf{指针和字段。} 状态字段要能说清：indegree[v] 是课程 v 尚未满足的先修数量，queue 保存当前可学习课程，visited\_count 是已处理数量。手算时按这个协议跟踪：0 出队后，遍历它的后继 1，indegree[1] 从 1 变 0，1 入队；最后处理数量等于课程数，说明无环。

\textbf{代码落点。} 关键是边方向要统一：prerequisite [a, b] 表示 b -> a。处理数量小于 numCourses 时返回 false。关键代码行必须能逐句翻译成“旧状态怎样变成新状态”，否则就只是模板。

\textbf{正确性。} 入度为 0 的节点没有未满足前置条件，可以安全学习；有向环中的节点永远不会全部降到入度 0。

\textbf{边界实验。} 先用边方向、孤立课程、环、自依赖、重复边做反例检查。实验时覆盖边方向、孤立课程、环、自依赖、重复边，并打印入队时的完整状态。若同一位置但资源不同的状态被合并，很多图题会出现隐蔽错误。

\textbf{复杂度变化。} 反复扫描最坏 O(VE)，拓扑排序 O(V+E) 时间和空间。

\textbf{迁移追问。} 如果题目改成“DFS 三色标记也能判环，但要明确访问状态”，先判断原状态是否仍然覆盖新答案集合，再决定是微调边界、改 value 语义，还是换算法。

\noindent\textbf{LeetCode 210 课程表 II。}

\textbf{题目说明。} LeetCode 210 课程表 II 的题面入口要先还原成具体任务：在判环基础上还要输出一个可行拓扑序。

\textbf{样例入口。} 拿 0->1、0->2 手算，课程 0 入度为 0，先输出 0，再把 1 和 2 的入度减到 0。若存在环，输出数量会少于课程数。

\textbf{暴力基线。} 慢版本先按题意画出完整状态图，用普通搜索跑通可达性。这个过程用于确认在判环基础上还要输出一个可行拓扑序的节点和边，之后再讨论访问标记、层次或代价顺序。放到本题就是：先按“在判环基础上还要输出一个可行拓扑序”完整试慢版本；样例里已经能看到“规律是课程表 II 比判定题多了输出顺序，队列弹出的顺序就是一种可行拓扑序”，所以优化前必须先能说清慢版本枚举了哪些候选。

\textbf{暴力过程。} 拿 0->1、0->2 手算，课程 0 入度为 0，先输出 0，再把 1 和 2 的入度减到 0。若存在环，输出数量会少于课程数。

\textbf{重复工作。} 题面模型：在判环基础上还要输出一个可行拓扑序。暴力版的慢点要落到本题：慢版本会围绕“在判环基础上还要输出一个可行拓扑序”反复展开同一批状态。样例规律是“规律是课程表 II 比判定题多了输出顺序，队列弹出的顺序就是一种可行拓扑序”，说明必须把节点、边、访问标记和资源维度一次建清；同一状态不应重复入队，不同资源状态也不能误合并。后面的优化只消掉这类重复工作，不能改变答案集合。

\textbf{优化条件。} 本题的关键观察是：每弹出一个入度为零课程，就把它加入顺序并删除出边。这句话必须能翻译成可维护状态；如果题目条件改变后观察不再成立，就不能继续套原来的写法。

\textbf{相邻状态。} 规律是课程表 II 比判定题多了输出顺序，队列弹出的顺序就是一种可行拓扑序。把这个特例继续拆开看：状态节点不一定等于原始位置，还可能包含钥匙、剩余资源、时间或访问集合。visited 只能合并未来能力完全相同的状态，否则会把合法路径剪掉。

\textbf{状态复用。} 把重复搜索压成访问标记、距离数组或拓扑状态；邻居生成也要从全量扫描变成结构化枚举。本题落点是：规律是课程表 II 比判定题多了输出顺序，队列弹出的顺序就是一种可行拓扑序。手算时只改被新输入影响的那一小部分，而不是回头重算完整候选。

\textbf{指针和字段。} 状态字段要能说清：状态编号、邻接生成方式、访问标记、距离或层数、队列内容；若状态含资源，visited 不能只按位置记录。手算时按这个协议跟踪：拿 0->1、0->2 手算，课程 0 入度为 0，先输出 0，再把 1 和 2 的入度减到 0。若存在环，输出数量会少于课程数。然后按协议复查：手算时画队列或栈，状态必须包含题目需要的所有资源。入队时标记还是出队时标记，要和去重语义一致。

\textbf{代码落点。} 入队时标记还是出队时标记必须一致；方向数组集中定义；多源 BFS 要一次性把所有源点放入初始队列。关键代码行必须能逐句翻译成“旧状态怎样变成新状态”，否则就只是模板。

\textbf{正确性。} 证明从可达性开始：搜索覆盖所有合法边能到达的状态，访问标记只删除重复状态而不删除不同未来能力。

\textbf{边界实验。} 先用有环返回空数组、多个可行顺序、边方向做反例检查。实验时覆盖有环返回空数组、多个可行顺序、边方向，并打印入队时的完整状态。若同一位置但资源不同的状态被合并，很多图题会出现隐蔽错误。

\textbf{复杂度变化。} 暴力重复从多个起点搜索会反复遍历图；正确建模和访问标记后，BFS 或 DFS 通常按节点数加边数计算，状态图题则按完整状态数计算。

\textbf{迁移追问。} 如果题目改成“若需要字典序最小拓扑序，把队列换成小顶堆”，先判断原状态是否仍然覆盖新答案集合，再决定是微调边界、改 value 语义，还是换算法。

\noindent\textbf{LeetCode 417 太平洋大西洋水流问题。}

\textbf{题目说明。} LeetCode 417 太平洋大西洋水流问题 的题面入口要先还原成具体任务：从每个格子向海搜索很慢，反向从海边沿非递减高度搜索。

\textbf{样例入口。} 拿高度矩阵 [[1, 2], [4, 3]] 手算，水从高处流向低处；反向从太平洋边界向内，只能走到更高或相等的格子。

\textbf{暴力基线。} 慢版本先按题意画出完整状态图，用普通搜索跑通可达性。这个过程用于确认从每个格子向海搜索很慢，反向从海边沿非递减高度搜索的节点和边，之后再讨论访问标记、层次或代价顺序。放到本题就是：先按“从每个格子向海搜索很慢，反向从海边沿非递减高度搜索”完整试慢版本；样例里已经能看到“规律是分别从两个海洋边界反向 BFS/DFS，两个访问集合的交集就是答案”，所以优化前必须先能说清慢版本枚举了哪些候选。

\textbf{暴力过程。} 拿高度矩阵 [[1, 2], [4, 3]] 手算，水从高处流向低处；反向从太平洋边界向内，只能走到更高或相等的格子。

\textbf{重复工作。} 题面模型：从每个格子向海搜索很慢，反向从海边沿非递减高度搜索。暴力版的慢点要落到本题：慢版本会围绕“从每个格子向海搜索很慢，反向从海边沿非递减高度搜索”反复展开同一批状态。样例规律是“规律是分别从两个海洋边界反向 BFS/DFS，两个访问集合的交集就是答案”，说明必须把节点、边、访问标记和资源维度一次建清；同一状态不应重复入队，不同资源状态也不能误合并。后面的优化只消掉这类重复工作，不能改变答案集合。

\textbf{优化条件。} 本题的关键观察是：分别标记能到太平洋和大西洋的格子，交集就是答案。这句话必须能翻译成可维护状态；如果题目条件改变后观察不再成立，就不能继续套原来的写法。

\textbf{相邻状态。} 规律是分别从两个海洋边界反向 BFS/DFS，两个访问集合的交集就是答案。把这个特例继续拆开看：状态节点不一定等于原始位置，还可能包含钥匙、剩余资源、时间或访问集合。visited 只能合并未来能力完全相同的状态，否则会把合法路径剪掉。

\textbf{状态复用。} 把重复搜索压成访问标记、距离数组或拓扑状态；邻居生成也要从全量扫描变成结构化枚举。本题落点是：规律是分别从两个海洋边界反向 BFS/DFS，两个访问集合的交集就是答案。手算时只改被新输入影响的那一小部分，而不是回头重算完整候选。

\textbf{指针和字段。} 状态字段要能说清：状态编号、邻接生成方式、访问标记、距离或层数、队列内容；若状态含资源，visited 不能只按位置记录。手算时按这个协议跟踪：拿高度矩阵 [[1, 2], [4, 3]] 手算，水从高处流向低处；反向从太平洋边界向内，只能走到更高或相等的格子。然后按协议复查：手算时画队列或栈，状态必须包含题目需要的所有资源。入队时标记还是出队时标记，要和去重语义一致。

\textbf{代码落点。} 入队时标记还是出队时标记必须一致；方向数组集中定义；多源 BFS 要一次性把所有源点放入初始队列。关键代码行必须能逐句翻译成“旧状态怎样变成新状态”，否则就只是模板。

\textbf{正确性。} 证明从可达性开始：搜索覆盖所有合法边能到达的状态，访问标记只删除重复状态而不删除不同未来能力。

\textbf{边界实验。} 先用单行单列、平台高度、边界重复入队、方向条件反向做反例检查。实验时覆盖单行单列、平台高度、边界重复入队、方向条件反向，并打印入队时的完整状态。若同一位置但资源不同的状态被合并，很多图题会出现隐蔽错误。

\textbf{复杂度变化。} 暴力重复从多个起点搜索会反复遍历图；正确建模和访问标记后，BFS 或 DFS 通常按节点数加边数计算，状态图题则按完整状态数计算。

\textbf{迁移追问。} 如果题目改成“反向搜索是很多可达性题的重要技巧”，先判断原状态是否仍然覆盖新答案集合，再决定是微调边界、改 value 语义，还是换算法。

\noindent\textbf{LeetCode 695 岛屿的最大面积。}

\textbf{题目说明。} LeetCode 695 岛屿的最大面积 的题面入口要先还原成具体任务：面积就是一个连通块中陆地格子的数量。

\textbf{样例入口。} 拿网格 [[1, 1, 0], [0, 1, 0]] 手算，从左上陆地出发能访问 3 个四连通格子；对角线不算连通。

\textbf{暴力基线。} 慢版本先按题意画出完整状态图，用普通搜索跑通可达性。这个过程用于确认面积就是一个连通块中陆地格子的数量的节点和边，之后再讨论访问标记、层次或代价顺序。放到本题就是：先按“面积就是一个连通块中陆地格子的数量”完整试慢版本；样例里已经能看到“规律是每发现未访问陆地就 DFS/BFS 统计面积并标记，最大值随每个连通块更新”，所以优化前必须先能说清慢版本枚举了哪些候选。

\textbf{暴力过程。} 拿网格 [[1, 1, 0], [0, 1, 0]] 手算，从左上陆地出发能访问 3 个四连通格子；对角线不算连通。

\textbf{重复工作。} 题面模型：面积就是一个连通块中陆地格子的数量。暴力版的慢点要落到本题：慢版本会围绕“面积就是一个连通块中陆地格子的数量”反复展开同一批状态。样例规律是“规律是每发现未访问陆地就 DFS/BFS 统计面积并标记，最大值随每个连通块更新”，说明必须把节点、边、访问标记和资源维度一次建清；同一状态不应重复入队，不同资源状态也不能误合并。后面的优化只消掉这类重复工作，不能改变答案集合。

\textbf{优化条件。} 本题的关键观察是：每发现新陆地就搜索完整连通块并计数，取最大值。这句话必须能翻译成可维护状态；如果题目条件改变后观察不再成立，就不能继续套原来的写法。

\textbf{相邻状态。} 规律是每发现未访问陆地就 DFS/BFS 统计面积并标记，最大值随每个连通块更新。把这个特例继续拆开看：状态节点不一定等于原始位置，还可能包含钥匙、剩余资源、时间或访问集合。visited 只能合并未来能力完全相同的状态，否则会把合法路径剪掉。

\textbf{状态复用。} 把重复搜索压成访问标记、距离数组或拓扑状态；邻居生成也要从全量扫描变成结构化枚举。本题落点是：规律是每发现未访问陆地就 DFS/BFS 统计面积并标记，最大值随每个连通块更新。手算时只改被新输入影响的那一小部分，而不是回头重算完整候选。

\textbf{指针和字段。} 状态字段要能说清：状态编号、邻接生成方式、访问标记、距离或层数、队列内容；若状态含资源，visited 不能只按位置记录。手算时按这个协议跟踪：拿网格 [[1, 1, 0], [0, 1, 0]] 手算，从左上陆地出发能访问 3 个四连通格子；对角线不算连通。然后按协议复查：手算时画队列或栈，状态必须包含题目需要的所有资源。入队时标记还是出队时标记，要和去重语义一致。

\textbf{代码落点。} 入队时标记还是出队时标记必须一致；方向数组集中定义；多源 BFS 要一次性把所有源点放入初始队列。关键代码行必须能逐句翻译成“旧状态怎样变成新状态”，否则就只是模板。

\textbf{正确性。} 证明从可达性开始：搜索覆盖所有合法边能到达的状态，访问标记只删除重复状态而不删除不同未来能力。

\textbf{边界实验。} 先用全水、全陆、细长岛、是否允许修改输入做反例检查。实验时覆盖全水、全陆、细长岛、是否允许修改输入，并打印入队时的完整状态。若同一位置但资源不同的状态被合并，很多图题会出现隐蔽错误。

\textbf{复杂度变化。} 暴力重复从多个起点搜索会反复遍历图；正确建模和访问标记后，BFS 或 DFS 通常按节点数加边数计算，状态图题则按完整状态数计算。

\textbf{迁移追问。} 如果题目改成“与岛屿数量相比，只是连通块聚合值不同”，先判断原状态是否仍然覆盖新答案集合，再决定是微调边界、改 value 语义，还是换算法。

\noindent\textbf{LeetCode 752 打开转盘锁。}

\textbf{题目说明。} LeetCode 752 打开转盘锁 的题面入口要先还原成具体任务：四位密码是状态，每次转动一位形成边。

\textbf{样例入口。} 拿锁 0000 到 0001 手算，每一步只能拨动一位加一或减一，0000 的邻居有 1000、9000、0100 等。

\textbf{暴力基线。} 慢版本先按题意画出完整状态图，用普通搜索跑通可达性。这个过程用于确认四位密码是状态，每次转动一位形成边的节点和边，之后再讨论访问标记、层次或代价顺序。放到本题就是：先按“四位密码是状态，每次转动一位形成边”完整试慢版本；样例里已经能看到“规律是状态是四位字符串，BFS 按步数扩展，deadend 入队前就要排除”，所以优化前必须先能说清慢版本枚举了哪些候选。

\textbf{暴力过程。} 拿锁 0000 到 0001 手算，每一步只能拨动一位加一或减一，0000 的邻居有 1000、9000、0100 等。

\textbf{重复工作。} 题面模型：四位密码是状态，每次转动一位形成边。暴力版的慢点要落到本题：慢版本会围绕“四位密码是状态，每次转动一位形成边”反复展开同一批状态。样例规律是“规律是状态是四位字符串，BFS 按步数扩展，deadend 入队前就要排除”，说明必须把节点、边、访问标记和资源维度一次建清；同一状态不应重复入队，不同资源状态也不能误合并。后面的优化只消掉这类重复工作，不能改变答案集合。

\textbf{优化条件。} 本题的关键观察是：BFS 从起点扩展，死亡状态不能入队，访问状态避免重复。这句话必须能翻译成可维护状态；如果题目条件改变后观察不再成立，就不能继续套原来的写法。

\textbf{相邻状态。} 规律是状态是四位字符串，BFS 按步数扩展，deadend 入队前就要排除。把这个特例继续拆开看：状态节点不一定等于原始位置，还可能包含钥匙、剩余资源、时间或访问集合。visited 只能合并未来能力完全相同的状态，否则会把合法路径剪掉。

\textbf{状态复用。} 把重复搜索压成访问标记、距离数组或拓扑状态；邻居生成也要从全量扫描变成结构化枚举。本题落点是：规律是状态是四位字符串，BFS 按步数扩展，deadend 入队前就要排除。手算时只改被新输入影响的那一小部分，而不是回头重算完整候选。

\textbf{指针和字段。} 状态字段要能说清：状态编号、邻接生成方式、访问标记、距离或层数、队列内容；若状态含资源，visited 不能只按位置记录。手算时按这个协议跟踪：拿锁 0000 到 0001 手算，每一步只能拨动一位加一或减一，0000 的邻居有 1000、9000、0100 等。然后按协议复查：手算时画队列或栈，状态必须包含题目需要的所有资源。入队时标记还是出队时标记，要和去重语义一致。

\textbf{代码落点。} 入队时标记还是出队时标记必须一致；方向数组集中定义；多源 BFS 要一次性把所有源点放入初始队列。关键代码行必须能逐句翻译成“旧状态怎样变成新状态”，否则就只是模板。

\textbf{正确性。} 证明从可达性开始：搜索覆盖所有合法边能到达的状态，访问标记只删除重复状态而不删除不同未来能力。

\textbf{边界实验。} 先用起点死亡、目标就是起点、数字环绕、死亡集合很大做反例检查。实验时覆盖起点死亡、目标就是起点、数字环绕、死亡集合很大，并打印入队时的完整状态。若同一位置但资源不同的状态被合并，很多图题会出现隐蔽错误。

\textbf{复杂度变化。} 暴力重复从多个起点搜索会反复遍历图；正确建模和访问标记后，BFS 或 DFS 通常按节点数加边数计算，状态图题则按完整状态数计算。

\textbf{迁移追问。} 如果题目改成“双向 BFS 能明显减少状态展开”，先判断原状态是否仍然覆盖新答案集合，再决定是微调边界、改 value 语义，还是换算法。

\noindent\textbf{LeetCode 785 判断二分图。}

\textbf{题目说明。} LeetCode 785 判断二分图 的题面入口要先还原成具体任务：二分图要求每条边两端颜色相反。

\textbf{样例入口。} 拿 graph=[[1, 2, 3], [0, 2], [0, 1, 3], [0, 2]] 手算，0 染红后 1、2、3 都要染蓝，但 1 和 2 之间还有边，蓝蓝相邻冲突，所以不是二分图。

\textbf{暴力基线。} 慢版本先按题意画出完整状态图，用普通搜索跑通可达性。这个过程用于确认二分图要求每条边两端颜色相反的节点和边，之后再讨论访问标记、层次或代价顺序。放到本题就是：先按“二分图要求每条边两端颜色相反”完整试慢版本；样例里已经能看到“规律是 BFS 染色时每条边两端必须异色，发现同色边立即失败”，所以优化前必须先能说清慢版本枚举了哪些候选。

\textbf{暴力过程。} 拿 graph=[[1, 2, 3], [0, 2], [0, 1, 3], [0, 2]] 手算，0 染红后 1、2、3 都要染蓝，但 1 和 2 之间还有边，蓝蓝相邻冲突，所以不是二分图。

\textbf{重复工作。} 题面模型：二分图要求每条边两端颜色相反。暴力版的慢点要落到本题：慢版本会围绕“二分图要求每条边两端颜色相反”反复展开同一批状态。样例规律是“规律是 BFS 染色时每条边两端必须异色，发现同色边立即失败”，说明必须把节点、边、访问标记和资源维度一次建清；同一状态不应重复入队，不同资源状态也不能误合并。后面的优化只消掉这类重复工作，不能改变答案集合。

\textbf{优化条件。} 本题的关键观察是：对每个未染色连通块 BFS 或 DFS 染色，遇到同色边立即失败。这句话必须能翻译成可维护状态；如果题目条件改变后观察不再成立，就不能继续套原来的写法。

\textbf{相邻状态。} 规律是 BFS 染色时每条边两端必须异色，发现同色边立即失败。把这个特例继续拆开看：状态节点不一定等于原始位置，还可能包含钥匙、剩余资源、时间或访问集合。visited 只能合并未来能力完全相同的状态，否则会把合法路径剪掉。

\textbf{状态复用。} 把重复搜索压成访问标记、距离数组或拓扑状态；邻居生成也要从全量扫描变成结构化枚举。本题落点是：规律是 BFS 染色时每条边两端必须异色，发现同色边立即失败。手算时只改被新输入影响的那一小部分，而不是回头重算完整候选。

\textbf{指针和字段。} 状态字段要能说清：状态编号、邻接生成方式、访问标记、距离或层数、队列内容；若状态含资源，visited 不能只按位置记录。手算时按这个协议跟踪：拿 graph=[[1, 2, 3], [0, 2], [0, 1, 3], [0, 2]] 手算，0 染红后 1、2、3 都要染蓝，但 1 和 2 之间还有边，蓝蓝相邻冲突，所以不是二分图。然后按协议复查：手算时画队列或栈，状态必须包含题目需要的所有资源。入队时标记还是出队时标记，要和去重语义一致。

\textbf{代码落点。} 入队时标记还是出队时标记必须一致；方向数组集中定义；多源 BFS 要一次性把所有源点放入初始队列。关键代码行必须能逐句翻译成“旧状态怎样变成新状态”，否则就只是模板。

\textbf{正确性。} 证明从可达性开始：搜索覆盖所有合法边能到达的状态，访问标记只删除重复状态而不删除不同未来能力。

\textbf{边界实验。} 先用非连通图、自环、奇环、孤立点、重复边做反例检查。实验时覆盖非连通图、自环、奇环、孤立点、重复边，并打印入队时的完整状态。若同一位置但资源不同的状态被合并，很多图题会出现隐蔽错误。

\textbf{复杂度变化。} 暴力重复从多个起点搜索会反复遍历图；正确建模和访问标记后，BFS 或 DFS 通常按节点数加边数计算，状态图题则按完整状态数计算。

\textbf{迁移追问。} 如果题目改成“可能的二分法是同一染色模型，只是图由 dislike 关系构造”，先判断原状态是否仍然覆盖新答案集合，再决定是微调边界、改 value 语义，还是换算法。

\noindent\textbf{LeetCode 815 公交路线。}

\textbf{题目说明。} LeetCode 815 公交路线 的题面入口要先还原成具体任务：公交路线和站点构成二分关系，换乘次数比站点步数更重要。

\textbf{样例入口。} 拿起点站可坐路线 A，A 和 B 在某站相交，B 到终点手算。按站点一步步 BFS 会反复展开同一条路线的所有站；按路线 BFS，一次乘车层数就是换乘次数。

\textbf{暴力基线。} 慢版本先按题意画出完整状态图，用普通搜索跑通可达性。这个过程用于确认公交路线和站点构成二分关系，换乘次数比站点步数更重要的节点和边，之后再讨论访问标记、层次或代价顺序。放到本题就是：先按“公交路线和站点构成二分关系，换乘次数比站点步数更重要”完整试慢版本；样例里已经能看到“规律是访问过的路线不再展开，站点到路线表只用于找到下一批路线”，所以优化前必须先能说清慢版本枚举了哪些候选。

\textbf{暴力过程。} 拿起点站可坐路线 A，A 和 B 在某站相交，B 到终点手算。按站点一步步 BFS 会反复展开同一条路线的所有站；按路线 BFS，一次乘车层数就是换乘次数。

\textbf{重复工作。} 题面模型：公交路线和站点构成二分关系，换乘次数比站点步数更重要。暴力版的慢点要落到本题：慢版本会围绕“公交路线和站点构成二分关系，换乘次数比站点步数更重要”反复展开同一批状态。样例规律是“规律是访问过的路线不再展开，站点到路线表只用于找到下一批路线”，说明必须把节点、边、访问标记和资源维度一次建清；同一状态不应重复入队，不同资源状态也不能误合并。后面的优化只消掉这类重复工作，不能改变答案集合。

\textbf{优化条件。} 本题的关键观察是：用站点到路线列表的哈希表找可乘路线，以路线或乘车层数做 BFS，访问过的路线不再重复展开。这句话必须能翻译成可维护状态；如果题目条件改变后观察不再成立，就不能继续套原来的写法。

\textbf{相邻状态。} 规律是访问过的路线不再展开，站点到路线表只用于找到下一批路线。把这个特例继续拆开看：状态节点不一定等于原始位置，还可能包含钥匙、剩余资源、时间或访问集合。visited 只能合并未来能力完全相同的状态，否则会把合法路径剪掉。

\textbf{状态复用。} 把重复搜索压成访问标记、距离数组或拓扑状态；邻居生成也要从全量扫描变成结构化枚举。本题落点是：规律是访问过的路线不再展开，站点到路线表只用于找到下一批路线。手算时只改被新输入影响的那一小部分，而不是回头重算完整候选。

\textbf{指针和字段。} 状态字段要能说清：状态编号、邻接生成方式、访问标记、距离或层数、队列内容；若状态含资源，visited 不能只按位置记录。手算时按这个协议跟踪：拿起点站可坐路线 A，A 和 B 在某站相交，B 到终点手算。按站点一步步 BFS 会反复展开同一条路线的所有站；按路线 BFS，一次乘车层数就是换乘次数。然后按协议复查：手算时画队列或栈，状态必须包含题目需要的所有资源。入队时标记还是出队时标记，要和去重语义一致。

\textbf{代码落点。} 入队时标记还是出队时标记必须一致；方向数组集中定义；多源 BFS 要一次性把所有源点放入初始队列。关键代码行必须能逐句翻译成“旧状态怎样变成新状态”，否则就只是模板。

\textbf{正确性。} 证明从可达性开始：搜索覆盖所有合法边能到达的状态，访问标记只删除重复状态而不删除不同未来能力。

\textbf{边界实验。} 先用起点等于终点、某站经过路线很多、路线重复访问、目标站不可达做反例检查。实验时覆盖起点等于终点、某站经过路线很多、路线重复访问、目标站不可达，并打印入队时的完整状态。若同一位置但资源不同的状态被合并，很多图题会出现隐蔽错误。

\textbf{复杂度变化。} 暴力重复从多个起点搜索会反复遍历图；正确建模和访问标记后，BFS 或 DFS 通常按节点数加边数计算，状态图题则按完整状态数计算。

\textbf{迁移追问。} 如果题目改成“若把每个站点之间完全建边会爆炸，应利用路线这一层压缩邻居生成”，先判断原状态是否仍然覆盖新答案集合，再决定是微调边界、改 value 语义，还是换算法。

\noindent\textbf{LeetCode 847 访问所有节点的最短路径。}

\textbf{题目说明。} LeetCode 847 访问所有节点的最短路径 的题面入口要先还原成具体任务：访问集合影响未来能力，同一节点配不同 mask 不是同一状态。

\textbf{样例入口。} 拿 graph=[[1, 2, 3], [0], [0], [0]] 手算，从中心 0 出发访问三个叶子需要来回走，长度为 4；如果只按节点 visited，到 0 时会把访问了不同叶子的状态混掉。

\textbf{暴力基线。} 慢版本先按题意画出完整状态图，用普通搜索跑通可达性。这个过程用于确认访问集合影响未来能力，同一节点配不同 mask 不是同一状态的节点和边，之后再讨论访问标记、层次或代价顺序。放到本题就是：先按“访问集合影响未来能力，同一节点配不同 mask 不是同一状态”完整试慢版本；样例里已经能看到“规律是 BFS 状态必须是 (node, mask)，从所有起点多源出发，第一次到全集 mask 就是最短”，所以优化前必须先能说清慢版本枚举了哪些候选。

\textbf{暴力过程。} 拿 graph=[[1, 2, 3], [0], [0], [0]] 手算，从中心 0 出发访问三个叶子需要来回走，长度为 4；如果只按节点 visited，到 0 时会把访问了不同叶子的状态混掉。

\textbf{重复工作。} 题面模型：访问集合影响未来能力，同一节点配不同 mask 不是同一状态。暴力版的慢点要落到本题：慢版本会围绕“访问集合影响未来能力，同一节点配不同 mask 不是同一状态”反复展开同一批状态。样例规律是“规律是 BFS 状态必须是 (node, mask)，从所有起点多源出发，第一次到全集 mask 就是最短”，说明必须把节点、边、访问标记和资源维度一次建清；同一状态不应重复入队，不同资源状态也不能误合并。后面的优化只消掉这类重复工作，不能改变答案集合。

\textbf{优化条件。} 本题的关键观察是：从所有节点以各自单点 mask 多源入队，BFS 扩展 (node, mask)，第一次到全集 mask 即最短。这句话必须能翻译成可维护状态；如果题目条件改变后观察不再成立，就不能继续套原来的写法。

\textbf{相邻状态。} 规律是 BFS 状态必须是 (node, mask)，从所有起点多源出发，第一次到全集 mask 就是最短。把这个特例继续拆开看：状态节点不一定等于原始位置，还可能包含钥匙、剩余资源、时间或访问集合。visited 只能合并未来能力完全相同的状态，否则会把合法路径剪掉。

\textbf{状态复用。} 把重复搜索压成访问标记、距离数组或拓扑状态；邻居生成也要从全量扫描变成结构化枚举。本题落点是：规律是 BFS 状态必须是 (node, mask)，从所有起点多源出发，第一次到全集 mask 就是最短。手算时只改被新输入影响的那一小部分，而不是回头重算完整候选。

\textbf{指针和字段。} 状态字段要能说清：状态编号、邻接生成方式、访问标记、距离或层数、队列内容；若状态含资源，visited 不能只按位置记录。手算时按这个协议跟踪：拿 graph=[[1, 2, 3], [0], [0], [0]] 手算，从中心 0 出发访问三个叶子需要来回走，长度为 4；如果只按节点 visited，到 0 时会把访问了不同叶子的状态混掉。然后按协议复查：手算时画队列或栈，状态必须包含题目需要的所有资源。入队时标记还是出队时标记，要和去重语义一致。

\textbf{代码落点。} 入队时标记还是出队时标记必须一致；方向数组集中定义；多源 BFS 要一次性把所有源点放入初始队列。关键代码行必须能逐句翻译成“旧状态怎样变成新状态”，否则就只是模板。

\textbf{正确性。} 证明从可达性开始：搜索覆盖所有合法边能到达的状态，访问标记只删除重复状态而不删除不同未来能力。

\textbf{边界实验。} 先用单节点、完全图、星形图、重复经过节点、visited 必须包含 mask做反例检查。实验时覆盖单节点、完全图、星形图、重复经过节点、visited 必须包含 mask，并打印入队时的完整状态。若同一位置但资源不同的状态被合并，很多图题会出现隐蔽错误。

\textbf{复杂度变化。} 暴力重复从多个起点搜索会反复遍历图；正确建模和访问标记后，BFS 或 DFS 通常按节点数加边数计算，状态图题则按完整状态数计算。

\textbf{迁移追问。} 如果题目改成“若边有权，状态图要改用 Dijkstra”，先判断原状态是否仍然覆盖新答案集合，再决定是微调边界、改 value 语义，还是换算法。

\noindent\textbf{LeetCode 864 获取所有钥匙的最短路径。}

\textbf{题目说明。} LeetCode 864 获取所有钥匙的最短路径 的题面入口要先还原成具体任务：位置相同但钥匙集合不同，未来能开的门不同。

\textbf{样例入口。} 拿网格 @.a / \#.\# / b.A 手算，走到 a 后状态不是位置变了而是钥匙集合多了一把；同一位置带不同钥匙未来能力不同。

\textbf{暴力基线。} 慢版本先按题意画出完整状态图，用普通搜索跑通可达性。这个过程用于确认位置相同但钥匙集合不同，未来能开的门不同的节点和边，之后再讨论访问标记、层次或代价顺序。放到本题就是：先按“位置相同但钥匙集合不同，未来能开的门不同”完整试慢版本；样例里已经能看到“规律是 BFS 的 visited 必须按 (row, col, keymask) 记录，拿齐所有钥匙第一次出队就是最短步数”，所以优化前必须先能说清慢版本枚举了哪些候选。

\textbf{暴力过程。} 拿网格 @.a / \#.\# / b.A 手算，走到 a 后状态不是位置变了而是钥匙集合多了一把；同一位置带不同钥匙未来能力不同。

\textbf{重复工作。} 题面模型：位置相同但钥匙集合不同，未来能开的门不同。暴力版的慢点要落到本题：慢版本会围绕“位置相同但钥匙集合不同，未来能开的门不同”反复展开同一批状态。样例规律是“规律是 BFS 的 visited 必须按 (row, col, keymask) 记录，拿齐所有钥匙第一次出队就是最短步数”，说明必须把节点、边、访问标记和资源维度一次建清；同一状态不应重复入队，不同资源状态也不能误合并。后面的优化只消掉这类重复工作，不能改变答案集合。

\textbf{优化条件。} 本题的关键观察是：BFS 状态包含行、列和钥匙掩码，访问标记也必须包含掩码。这句话必须能翻译成可维护状态；如果题目条件改变后观察不再成立，就不能继续套原来的写法。

\textbf{相邻状态。} 规律是 BFS 的 visited 必须按 (row, col, keymask) 记录，拿齐所有钥匙第一次出队就是最短步数。把这个特例继续拆开看：状态节点不一定等于原始位置，还可能包含钥匙、剩余资源、时间或访问集合。visited 只能合并未来能力完全相同的状态，否则会把合法路径剪掉。

\textbf{状态复用。} 把重复搜索压成访问标记、距离数组或拓扑状态；邻居生成也要从全量扫描变成结构化枚举。本题落点是：规律是 BFS 的 visited 必须按 (row, col, keymask) 记录，拿齐所有钥匙第一次出队就是最短步数。手算时只改被新输入影响的那一小部分，而不是回头重算完整候选。

\textbf{指针和字段。} 状态字段要能说清：状态编号、邻接生成方式、访问标记、距离或层数、队列内容；若状态含资源，visited 不能只按位置记录。手算时按这个协议跟踪：拿网格 @.a / \#.\# / b.A 手算，走到 a 后状态不是位置变了而是钥匙集合多了一把；同一位置带不同钥匙未来能力不同。然后按协议复查：手算时画队列或栈，状态必须包含题目需要的所有资源。入队时标记还是出队时标记，要和去重语义一致。

\textbf{代码落点。} 入队时标记还是出队时标记必须一致；方向数组集中定义；多源 BFS 要一次性把所有源点放入初始队列。关键代码行必须能逐句翻译成“旧状态怎样变成新状态”，否则就只是模板。

\textbf{正确性。} 证明从可达性开始：搜索覆盖所有合法边能到达的状态，访问标记只删除重复状态而不删除不同未来能力。

\textbf{边界实验。} 先用起点旁边有门、钥匙顺序、不可达、钥匙数量决定掩码大小做反例检查。实验时覆盖起点旁边有门、钥匙顺序、不可达、钥匙数量决定掩码大小，并打印入队时的完整状态。若同一位置但资源不同的状态被合并，很多图题会出现隐蔽错误。

\textbf{复杂度变化。} 暴力重复从多个起点搜索会反复遍历图；正确建模和访问标记后，BFS 或 DFS 通常按节点数加边数计算，状态图题则按完整状态数计算。

\textbf{迁移追问。} 如果题目改成“这是状态图维度不能只按位置去重的典型题”，先判断原状态是否仍然覆盖新答案集合，再决定是微调边界、改 value 语义，还是换算法。

\noindent\textbf{LeetCode 886 可能的二分法。}

\textbf{题目说明。} LeetCode 886 可能的二分法 的题面入口要先还原成具体任务：互相讨厌的人必须分到不同组，本质是图二染色。

\textbf{样例入口。} 拿 dislikes=[[1, 2], [1, 3], [2, 4]] 手算，把 1 放 A，则 2 和 3 必须放 B，2 讨厌 4，所以 4 放 A，可行。若出现奇环，某条边两端会被迫同色。

\textbf{暴力基线。} 慢版本先按题意画出完整状态图，用普通搜索跑通可达性。这个过程用于确认互相讨厌的人必须分到不同组，本质是图二染色的节点和边，之后再讨论访问标记、层次或代价顺序。放到本题就是：先按“互相讨厌的人必须分到不同组，本质是图二染色”完整试慢版本；样例里已经能看到“规律是把 dislike 当无向边做二染色，发现同色冲突就不能二分”，所以优化前必须先能说清慢版本枚举了哪些候选。

\textbf{暴力过程。} 拿 dislikes=[[1, 2], [1, 3], [2, 4]] 手算，把 1 放 A，则 2 和 3 必须放 B，2 讨厌 4，所以 4 放 A，可行。若出现奇环，某条边两端会被迫同色。

\textbf{重复工作。} 题面模型：互相讨厌的人必须分到不同组，本质是图二染色。暴力版的慢点要落到本题：慢版本会围绕“互相讨厌的人必须分到不同组，本质是图二染色”反复展开同一批状态。样例规律是“规律是把 dislike 当无向边做二染色，发现同色冲突就不能二分”，说明必须把节点、边、访问标记和资源维度一次建清；同一状态不应重复入队，不同资源状态也不能误合并。后面的优化只消掉这类重复工作，不能改变答案集合。

\textbf{优化条件。} 本题的关键观察是：按 dislike 建无向图，对每个连通块 BFS 染色，发现同色冲突则不可行。这句话必须能翻译成可维护状态；如果题目条件改变后观察不再成立，就不能继续套原来的写法。

\textbf{相邻状态。} 规律是把 dislike 当无向边做二染色，发现同色冲突就不能二分。把这个特例继续拆开看：状态节点不一定等于原始位置，还可能包含钥匙、剩余资源、时间或访问集合。visited 只能合并未来能力完全相同的状态，否则会把合法路径剪掉。

\textbf{状态复用。} 把重复搜索压成访问标记、距离数组或拓扑状态；邻居生成也要从全量扫描变成结构化枚举。本题落点是：规律是把 dislike 当无向边做二染色，发现同色冲突就不能二分。手算时只改被新输入影响的那一小部分，而不是回头重算完整候选。

\textbf{指针和字段。} 状态字段要能说清：状态编号、邻接生成方式、访问标记、距离或层数、队列内容；若状态含资源，visited 不能只按位置记录。手算时按这个协议跟踪：拿 dislikes=[[1, 2], [1, 3], [2, 4]] 手算，把 1 放 A，则 2 和 3 必须放 B，2 讨厌 4，所以 4 放 A，可行。若出现奇环，某条边两端会被迫同色。然后按协议复查：手算时画队列或栈，状态必须包含题目需要的所有资源。入队时标记还是出队时标记，要和去重语义一致。

\textbf{代码落点。} 入队时标记还是出队时标记必须一致；方向数组集中定义；多源 BFS 要一次性把所有源点放入初始队列。关键代码行必须能逐句翻译成“旧状态怎样变成新状态”，否则就只是模板。

\textbf{正确性。} 证明从可达性开始：搜索覆盖所有合法边能到达的状态，访问标记只删除重复状态而不删除不同未来能力。

\textbf{边界实验。} 先用孤立点、非连通图、奇环、重复 dislike、自身 dislike做反例检查。实验时覆盖孤立点、非连通图、奇环、重复 dislike、自身 dislike，并打印入队时的完整状态。若同一位置但资源不同的状态被合并，很多图题会出现隐蔽错误。

\textbf{复杂度变化。} 暴力重复从多个起点搜索会反复遍历图；正确建模和访问标记后，BFS 或 DFS 通常按节点数加边数计算，状态图题则按完整状态数计算。

\textbf{迁移追问。} 如果题目改成“它和判断二分图同型，只是输入不是邻接表”，先判断原状态是否仍然覆盖新答案集合，再决定是微调边界、改 value 语义，还是换算法。

\noindent\textbf{LeetCode 994 腐烂的橘子。}

\textbf{题目说明。} LeetCode 994 腐烂的橘子 的题面入口要先还原成具体任务：所有初始腐烂橘子同时作为 BFS 源点。

\textbf{样例入口。} 拿 [[2, 1, 1], [1, 1, 0], [0, 1, 1]] 手算，所有腐烂橘子同时向四周扩散，第一分钟会感染相邻新鲜橘子。

\textbf{暴力基线。} 暴力按分钟模拟整张网格：每一分钟扫描所有格子，找到会被感染的新鲜橘子并统一变腐烂。

\textbf{暴力过程。} 在 [[2, 1, 1], [1, 1, 0], [0, 1, 1]] 中，初始腐烂橘子同时向四周扩散；第一分钟感染它上下左右的新鲜橘子。

\textbf{重复工作。} 题面模型：所有初始腐烂橘子同时作为 BFS 源点。暴力版的慢点要落到本题：浪费在每分钟都扫描全图。扩散只会从上一分钟刚腐烂的边界继续向外，队列正好保存这个边界。后面的优化只消掉这类重复工作，不能改变答案集合。

\textbf{优化条件。} 本题的关键观察是：按层扩散，分钟数等于最后一个新鲜橘子被首次访问的层数。这句话必须能翻译成可维护状态；如果题目条件改变后观察不再成立，就不能继续套原来的写法。

\textbf{相邻状态。} 规律是多源 BFS，初始所有腐烂橘子同时入队，层数就是分钟数。把这个特例继续拆开看：状态节点不一定等于原始位置，还可能包含钥匙、剩余资源、时间或访问集合。visited 只能合并未来能力完全相同的状态，否则会把合法路径剪掉。

\textbf{状态复用。} 多源 BFS。初始把所有腐烂橘子入队，并统计 fresh 数量；每弹出一层代表一分钟，感染相邻新鲜橘子并 fresh--。手算时只改被新输入影响的那一小部分，而不是回头重算完整候选。

\textbf{指针和字段。} 状态字段要能说清：queue 保存当前扩散边界，fresh 是剩余新鲜橘子数，minutes 是已经完成的层数，directions 是四方向。手算时按这个协议跟踪：第一层处理初始腐烂源，感染邻居入队；第二层只处理第一分钟新腐烂的橘子，层数自然对应分钟。

\textbf{代码落点。} 关键是所有初始腐烂橘子一次性入队。感染时立即把 grid[nr][nc] 改成 2，避免同一分钟被重复入队。关键代码行必须能逐句翻译成“旧状态怎样变成新状态”，否则就只是模板。

\textbf{正确性。} BFS 按层扩展，某个新鲜橘子第一次被访问的层数就是它能被腐烂的最早分钟；多源初始化表示所有源同时开始。

\textbf{边界实验。} 先用没有新鲜橘子、没有腐烂源、隔离新鲜橘子、空格阻断做反例检查。实验时覆盖没有新鲜橘子、没有腐烂源、隔离新鲜橘子、空格阻断，并打印入队时的完整状态。若同一位置但资源不同的状态被合并，很多图题会出现隐蔽错误。

\textbf{复杂度变化。} 逐分钟全图扫描最坏 O(mn*分钟数)，BFS 每格最多入队一次，O(mn) 时间、O(mn) 空间。

\textbf{迁移追问。} 如果题目改成“多源 BFS 也适用于最近零距离、地图扩散等题”，先判断原状态是否仍然覆盖新答案集合，再决定是微调边界、改 value 语义，还是换算法。

\noindent\textbf{LeetCode 1091 二进制矩阵中的最短路径。}

\textbf{题目说明。} LeetCode 1091 二进制矩阵中的最短路径 的题面入口要先还原成具体任务：八方向网格无权边，答案是从左上到右下的最少格子数。

\textbf{样例入口。} 拿 [[0, 1], [1, 0]] 手算，从左上可以斜着一步到右下，路径长度是 2；若起点或终点为 1 直接不可达。

\textbf{暴力基线。} 慢版本先按题意画出完整状态图，用普通搜索跑通可达性。这个过程用于确认八方向网格无权边，答案是从左上到右下的最少格子数的节点和边，之后再讨论访问标记、层次或代价顺序。放到本题就是：先按“八方向网格无权边，答案是从左上到右下的最少格子数”完整试慢版本；样例里已经能看到“规律是八方向 BFS，第一次到达终点就是最短路径长度”，所以优化前必须先能说清慢版本枚举了哪些候选。

\textbf{暴力过程。} 拿 [[0, 1], [1, 0]] 手算，从左上可以斜着一步到右下，路径长度是 2；若起点或终点为 1 直接不可达。

\textbf{重复工作。} 题面模型：八方向网格无权边，答案是从左上到右下的最少格子数。暴力版的慢点要落到本题：慢版本会围绕“八方向网格无权边，答案是从左上到右下的最少格子数”反复展开同一批状态。样例规律是“规律是八方向 BFS，第一次到达终点就是最短路径长度”，说明必须把节点、边、访问标记和资源维度一次建清；同一状态不应重复入队，不同资源状态也不能误合并。后面的优化只消掉这类重复工作，不能改变答案集合。

\textbf{优化条件。} 本题的关键观察是：BFS 入队时标记访问，层数或距离随状态保存。这句话必须能翻译成可维护状态；如果题目条件改变后观察不再成立，就不能继续套原来的写法。

\textbf{相邻状态。} 规律是八方向 BFS，第一次到达终点就是最短路径长度。把这个特例继续拆开看：状态节点不一定等于原始位置，还可能包含钥匙、剩余资源、时间或访问集合。visited 只能合并未来能力完全相同的状态，否则会把合法路径剪掉。

\textbf{状态复用。} 把重复搜索压成访问标记、距离数组或拓扑状态；邻居生成也要从全量扫描变成结构化枚举。本题落点是：规律是八方向 BFS，第一次到达终点就是最短路径长度。手算时只改被新输入影响的那一小部分，而不是回头重算完整候选。

\textbf{指针和字段。} 状态字段要能说清：状态编号、邻接生成方式、访问标记、距离或层数、队列内容；若状态含资源，visited 不能只按位置记录。手算时按这个协议跟踪：拿 [[0, 1], [1, 0]] 手算，从左上可以斜着一步到右下，路径长度是 2；若起点或终点为 1 直接不可达。然后按协议复查：手算时画队列或栈，状态必须包含题目需要的所有资源。入队时标记还是出队时标记，要和去重语义一致。

\textbf{代码落点。} 入队时标记还是出队时标记必须一致；方向数组集中定义；多源 BFS 要一次性把所有源点放入初始队列。关键代码行必须能逐句翻译成“旧状态怎样变成新状态”，否则就只是模板。

\textbf{正确性。} 证明从可达性开始：搜索覆盖所有合法边能到达的状态，访问标记只删除重复状态而不删除不同未来能力。

\textbf{边界实验。} 先用起点或终点阻塞、单格、八方向越界、无路可达做反例检查。实验时覆盖起点或终点阻塞、单格、八方向越界、无路可达，并打印入队时的完整状态。若同一位置但资源不同的状态被合并，很多图题会出现隐蔽错误。

\textbf{复杂度变化。} 暴力重复从多个起点搜索会反复遍历图；正确建模和访问标记后，BFS 或 DFS 通常按节点数加边数计算，状态图题则按完整状态数计算。

\textbf{迁移追问。} 如果题目改成“边权相同才能用 BFS，若每格有代价就要最短路”，先判断原状态是否仍然覆盖新答案集合，再决定是微调边界、改 value 语义，还是换算法。

\noindent\textbf{LeetCode 1192 查找集群内的关键连接。}

\textbf{题目说明。} LeetCode 1192 查找集群内的关键连接 的题面入口要先还原成具体任务：无向图中的桥是删除后会断开连通性的边。

\textbf{样例入口。} 拿 n=4、connections=[[0, 1], [1, 2], [2, 0], [1, 3]] 手算，0-1-2 形成环，删除其中任一条仍可互通；节点 3 只通过边 1-3 接入，删除它会断开。

\textbf{暴力基线。} 慢版本先按题意画出完整状态图，用普通搜索跑通可达性。这个过程用于确认无向图中的桥是删除后会断开连通性的边的节点和边，之后再讨论访问标记、层次或代价顺序。放到本题就是：先按“无向图中的桥是删除后会断开连通性的边”完整试慢版本；样例里已经能看到“规律是 DFS 中 low[child] 若大于 disc[parent]，说明 child 子树没有备用返祖通道，这条边就是桥”，所以优化前必须先能说清慢版本枚举了哪些候选。

\textbf{暴力过程。} 拿 n=4、connections=[[0, 1], [1, 2], [2, 0], [1, 3]] 手算，0-1-2 形成环，删除其中任一条仍可互通；节点 3 只通过边 1-3 接入，删除它会断开。

\textbf{重复工作。} 题面模型：无向图中的桥是删除后会断开连通性的边。暴力版的慢点要落到本题：慢版本会围绕“无向图中的桥是删除后会断开连通性的边”反复展开同一批状态。样例规律是“规律是 DFS 中 low[child] 若大于 disc[parent]，说明 child 子树没有备用返祖通道，这条边就是桥”，说明必须把节点、边、访问标记和资源维度一次建清；同一状态不应重复入队，不同资源状态也不能误合并。后面的优化只消掉这类重复工作，不能改变答案集合。

\textbf{优化条件。} 本题的关键观察是：DFS 记录 disc 和 low，若 low[child] 大于 disc[parent]，子树没有返祖通道，父子边是桥。这句话必须能翻译成可维护状态；如果题目条件改变后观察不再成立，就不能继续套原来的写法。

\textbf{相邻状态。} 规律是 DFS 中 low[child] 若大于 disc[parent]，说明 child 子树没有备用返祖通道，这条边就是桥。把这个特例继续拆开看：状态节点不一定等于原始位置，还可能包含钥匙、剩余资源、时间或访问集合。visited 只能合并未来能力完全相同的状态，否则会把合法路径剪掉。

\textbf{状态复用。} 把重复搜索压成访问标记、距离数组或拓扑状态；邻居生成也要从全量扫描变成结构化枚举。本题落点是：规律是 DFS 中 low[child] 若大于 disc[parent]，说明 child 子树没有备用返祖通道，这条边就是桥。手算时只改被新输入影响的那一小部分，而不是回头重算完整候选。

\textbf{指针和字段。} 状态字段要能说清：状态编号、邻接生成方式、访问标记、距离或层数、队列内容；若状态含资源，visited 不能只按位置记录。手算时按这个协议跟踪：拿 n=4、connections=[[0, 1], [1, 2], [2, 0], [1, 3]] 手算，0-1-2 形成环，删除其中任一条仍可互通；节点 3 只通过边 1-3 接入，删除它会断开。然后按协议复查：手算时画队列或栈，状态必须包含题目需要的所有资源。入队时标记还是出队时标记，要和去重语义一致。

\textbf{代码落点。} 入队时标记还是出队时标记必须一致；方向数组集中定义；多源 BFS 要一次性把所有源点放入初始队列。关键代码行必须能逐句翻译成“旧状态怎样变成新状态”，否则就只是模板。

\textbf{正确性。} 证明从可达性开始：搜索覆盖所有合法边能到达的状态，访问标记只删除重复状态而不删除不同未来能力。

\textbf{边界实验。} 先用重边、根节点、图保证连通、无向边需要编号跳过反向边、递归栈做反例检查。实验时覆盖重边、根节点、图保证连通、无向边需要编号跳过反向边、递归栈，并打印入队时的完整状态。若同一位置但资源不同的状态被合并，很多图题会出现隐蔽错误。

\textbf{复杂度变化。} 暴力重复从多个起点搜索会反复遍历图；正确建模和访问标记后，BFS 或 DFS 通常按节点数加边数计算，状态图题则按完整状态数计算。

\textbf{迁移追问。} 如果题目改成“若要找割点，判断条件和根节点处理不同”，先判断原状态是否仍然覆盖新答案集合，再决定是微调边界、改 value 语义，还是换算法。

\noindent\textbf{LeetCode 1293 网格中的最短路径。}

\textbf{题目说明。} LeetCode 1293 网格中的最短路径 的题面入口要先还原成具体任务：走到同一格时，剩余可消除障碍数量不同会影响未来可达性。

\textbf{样例入口。} 拿网格三列中间有障碍且 k=1 手算，走到同一格时，如果剩余消障次数不同，未来可走路径也不同。

\textbf{暴力基线。} 慢版本先按题意画出完整状态图，用普通搜索跑通可达性。这个过程用于确认走到同一格时，剩余可消除障碍数量不同会影响未来可达性的节点和边，之后再讨论访问标记、层次或代价顺序。放到本题就是：先按“走到同一格时，剩余可消除障碍数量不同会影响未来可达性”完整试慢版本；样例里已经能看到“规律是 BFS 的状态必须包含剩余 k，visited 不能只按位置记录”，所以优化前必须先能说清慢版本枚举了哪些候选。

\textbf{暴力过程。} 拿网格三列中间有障碍且 k=1 手算，走到同一格时，如果剩余消障次数不同，未来可走路径也不同。

\textbf{重复工作。} 题面模型：走到同一格时，剩余可消除障碍数量不同会影响未来可达性。暴力版的慢点要落到本题：慢版本会围绕“走到同一格时，剩余可消除障碍数量不同会影响未来可达性”反复展开同一批状态。样例规律是“规律是 BFS 的状态必须包含剩余 k，visited 不能只按位置记录”，说明必须把节点、边、访问标记和资源维度一次建清；同一状态不应重复入队，不同资源状态也不能误合并。后面的优化只消掉这类重复工作，不能改变答案集合。

\textbf{优化条件。} 本题的关键观察是：BFS 状态包含位置和剩余消除次数，visited 记录该维度或记录到达该格的最大剩余次数。这句话必须能翻译成可维护状态；如果题目条件改变后观察不再成立，就不能继续套原来的写法。

\textbf{相邻状态。} 规律是 BFS 的状态必须包含剩余 k，visited 不能只按位置记录。把这个特例继续拆开看：状态节点不一定等于原始位置，还可能包含钥匙、剩余资源、时间或访问集合。visited 只能合并未来能力完全相同的状态，否则会把合法路径剪掉。

\textbf{状态复用。} 把重复搜索压成访问标记、距离数组或拓扑状态；邻居生成也要从全量扫描变成结构化枚举。本题落点是：规律是 BFS 的状态必须包含剩余 k，visited 不能只按位置记录。手算时只改被新输入影响的那一小部分，而不是回头重算完整候选。

\textbf{指针和字段。} 状态字段要能说清：状态编号、邻接生成方式、访问标记、距离或层数、队列内容；若状态含资源，visited 不能只按位置记录。手算时按这个协议跟踪：拿网格三列中间有障碍且 k=1 手算，走到同一格时，如果剩余消障次数不同，未来可走路径也不同。然后按协议复查：手算时画队列或栈，状态必须包含题目需要的所有资源。入队时标记还是出队时标记，要和去重语义一致。

\textbf{代码落点。} 入队时标记还是出队时标记必须一致；方向数组集中定义；多源 BFS 要一次性把所有源点放入初始队列。关键代码行必须能逐句翻译成“旧状态怎样变成新状态”，否则就只是模板。

\textbf{正确性。} 证明从可达性开始：搜索覆盖所有合法边能到达的状态，访问标记只删除重复状态而不删除不同未来能力。

\textbf{边界实验。} 先用起点终点、k 很大、障碍密集、二维 visited 错误剪枝做反例检查。实验时覆盖起点终点、k 很大、障碍密集、二维 visited 错误剪枝，并打印入队时的完整状态。若同一位置但资源不同的状态被合并，很多图题会出现隐蔽错误。

\textbf{复杂度变化。} 暴力重复从多个起点搜索会反复遍历图；正确建模和访问标记后，BFS 或 DFS 通常按节点数加边数计算，状态图题则按完整状态数计算。

\textbf{迁移追问。} 如果题目改成“它和钥匙题一样说明状态维度不能丢”，先判断原状态是否仍然覆盖新答案集合，再决定是微调边界、改 value 语义，还是换算法。

\noindent\textbf{LeetCode 1462 课程表 IV。}

\textbf{题目说明。} LeetCode 1462 课程表 IV 的题面入口要先还原成具体任务：大量先修关系查询需要预处理课程之间的可达性。

\textbf{样例入口。} 拿 0->1->2 和查询 0 是否先修 2 手算，直接边只能回答 0->1、1->2，不能回答传递关系。若查询很多，每次 DFS 会重复走相同路径。

\textbf{暴力基线。} 慢版本先按题意画出完整状态图，用普通搜索跑通可达性。这个过程用于确认大量先修关系查询需要预处理课程之间的可达性的节点和边，之后再讨论访问标记、层次或代价顺序。放到本题就是：先按“大量先修关系查询需要预处理课程之间的可达性”完整试慢版本；样例里已经能看到“规律是先做可达性闭包或拓扑合并祖先集合，再 O(1) 回答查询”，所以优化前必须先能说清慢版本枚举了哪些候选。

\textbf{暴力过程。} 拿 0->1->2 和查询 0 是否先修 2 手算，直接边只能回答 0->1、1->2，不能回答传递关系。若查询很多，每次 DFS 会重复走相同路径。

\textbf{重复工作。} 题面模型：大量先修关系查询需要预处理课程之间的可达性。暴力版的慢点要落到本题：慢版本会围绕“大量先修关系查询需要预处理课程之间的可达性”反复展开同一批状态。样例规律是“规律是先做可达性闭包或拓扑合并祖先集合，再 O(1) 回答查询”，说明必须把节点、边、访问标记和资源维度一次建清；同一状态不应重复入队，不同资源状态也不能误合并。后面的优化只消掉这类重复工作，不能改变答案集合。

\textbf{优化条件。} 本题的关键观察是：可以用 Floyd 传递闭包、每个点 BFS，或拓扑 DP 合并祖先集合，查询时直接查可达表。这句话必须能翻译成可维护状态；如果题目条件改变后观察不再成立，就不能继续套原来的写法。

\textbf{相邻状态。} 规律是先做可达性闭包或拓扑合并祖先集合，再 O(1) 回答查询。把这个特例继续拆开看：状态节点不一定等于原始位置，还可能包含钥匙、剩余资源、时间或访问集合。visited 只能合并未来能力完全相同的状态，否则会把合法路径剪掉。

\textbf{状态复用。} 把重复搜索压成访问标记、距离数组或拓扑状态；邻居生成也要从全量扫描变成结构化枚举。本题落点是：规律是先做可达性闭包或拓扑合并祖先集合，再 O(1) 回答查询。手算时只改被新输入影响的那一小部分，而不是回头重算完整候选。

\textbf{指针和字段。} 状态字段要能说清：状态编号、邻接生成方式、访问标记、距离或层数、队列内容；若状态含资源，visited 不能只按位置记录。手算时按这个协议跟踪：拿 0->1->2 和查询 0 是否先修 2 手算，直接边只能回答 0->1、1->2，不能回答传递关系。若查询很多，每次 DFS 会重复走相同路径。然后按协议复查：手算时画队列或栈，状态必须包含题目需要的所有资源。入队时标记还是出队时标记，要和去重语义一致。

\textbf{代码落点。} 入队时标记还是出队时标记必须一致；方向数组集中定义；多源 BFS 要一次性把所有源点放入初始队列。关键代码行必须能逐句翻译成“旧状态怎样变成新状态”，否则就只是模板。

\textbf{正确性。} 证明从可达性开始：搜索覆盖所有合法边能到达的状态，访问标记只删除重复状态而不删除不同未来能力。

\textbf{边界实验。} 先用课程之间没有关系、自身是否算先修、查询很多、图中按题意应为 DAG做反例检查。实验时覆盖课程之间没有关系、自身是否算先修、查询很多、图中按题意应为 DAG，并打印入队时的完整状态。若同一位置但资源不同的状态被合并，很多图题会出现隐蔽错误。

\textbf{复杂度变化。} 暴力重复从多个起点搜索会反复遍历图；正确建模和访问标记后，BFS 或 DFS 通常按节点数加边数计算，状态图题则按完整状态数计算。

\textbf{迁移追问。} 如果题目改成“如果查询很少，逐次搜索可能更省预处理成本；多查询时闭包更稳”，先判断原状态是否仍然覆盖新答案集合，再决定是微调边界、改 value 语义，还是换算法。

\noindent\textbf{LeetCode 1584 连接所有点的最小费用。}

\textbf{题目说明。} LeetCode 1584 连接所有点的最小费用 的题面入口要先还原成具体任务：完全图上连接所有点的最小总成本是最小生成树问题，边权为曼哈顿距离。

\textbf{样例入口。} 拿 points=[[0, 0], [2, 2], [3, 10], [5, 2], [7, 0]] 手算，任意两点都有曼哈顿距离边，目标是选 n-1 条边连通且总成本最小。

\textbf{暴力基线。} 慢版本先按题意画出完整状态图，用普通搜索跑通可达性。这个过程用于确认完全图上连接所有点的最小总成本是最小生成树问题，边权为曼哈顿距离的节点和边，之后再讨论访问标记、层次或代价顺序。放到本题就是：先按“完全图上连接所有点的最小总成本是最小生成树问题，边权为曼哈顿距离”完整试慢版本；样例里已经能看到“规律是最小生成树问题，Prim 维护每个外部点到已连通集合的最小边，每次加入最便宜外部点”，所以优化前必须先能说清慢版本枚举了哪些候选。

\textbf{暴力过程。} 拿 points=[[0, 0], [2, 2], [3, 10], [5, 2], [7, 0]] 手算，任意两点都有曼哈顿距离边，目标是选 n-1 条边连通且总成本最小。

\textbf{重复工作。} 题面模型：完全图上连接所有点的最小总成本是最小生成树问题，边权为曼哈顿距离。暴力版的慢点要落到本题：慢版本会围绕“完全图上连接所有点的最小总成本是最小生成树问题，边权为曼哈顿距离”反复展开同一批状态。样例规律是“规律是最小生成树问题，Prim 维护每个外部点到已连通集合的最小边，每次加入最便宜外部点”，说明必须把节点、边、访问标记和资源维度一次建清；同一状态不应重复入队，不同资源状态也不能误合并。后面的优化只消掉这类重复工作，不能改变答案集合。

\textbf{优化条件。} 本题的关键观察是：Prim 从已连通点集扩展，每次选到外部点的最小距离；也可 Kruskal 显式建边排序。这句话必须能翻译成可维护状态；如果题目条件改变后观察不再成立，就不能继续套原来的写法。

\textbf{相邻状态。} 规律是最小生成树问题，Prim 维护每个外部点到已连通集合的最小边，每次加入最便宜外部点。把这个特例继续拆开看：状态节点不一定等于原始位置，还可能包含钥匙、剩余资源、时间或访问集合。visited 只能合并未来能力完全相同的状态，否则会把合法路径剪掉。

\textbf{状态复用。} 把重复搜索压成访问标记、距离数组或拓扑状态；邻居生成也要从全量扫描变成结构化枚举。本题落点是：规律是最小生成树问题，Prim 维护每个外部点到已连通集合的最小边，每次加入最便宜外部点。手算时只改被新输入影响的那一小部分，而不是回头重算完整候选。

\textbf{指针和字段。} 状态字段要能说清：状态编号、邻接生成方式、访问标记、距离或层数、队列内容；若状态含资源，visited 不能只按位置记录。手算时按这个协议跟踪：拿 points=[[0, 0], [2, 2], [3, 10], [5, 2], [7, 0]] 手算，任意两点都有曼哈顿距离边，目标是选 n-1 条边连通且总成本最小。然后按协议复查：手算时画队列或栈，状态必须包含题目需要的所有资源。入队时标记还是出队时标记，要和去重语义一致。

\textbf{代码落点。} 入队时标记还是出队时标记必须一致；方向数组集中定义；多源 BFS 要一次性把所有源点放入初始队列。关键代码行必须能逐句翻译成“旧状态怎样变成新状态”，否则就只是模板。

\textbf{正确性。} 证明从可达性开始：搜索覆盖所有合法边能到达的状态，访问标记只删除重复状态而不删除不同未来能力。

\textbf{边界实验。} 先用单点、重复坐标、完全图边数、曼哈顿距离、Prim 更新外部点最小连边做反例检查。实验时覆盖单点、重复坐标、完全图边数、曼哈顿距离、Prim 更新外部点最小连边，并打印入队时的完整状态。若同一位置但资源不同的状态被合并，很多图题会出现隐蔽错误。

\textbf{复杂度变化。} 暴力重复从多个起点搜索会反复遍历图；正确建模和访问标记后，BFS 或 DFS 通常按节点数加边数计算，状态图题则按完整状态数计算。

\textbf{迁移追问。} 如果题目改成“规模很大时要研究几何图优化，不能显式 O(\ensuremath{n^{2}}) 建边”，先判断原状态是否仍然覆盖新答案集合，再决定是微调边界、改 value 语义，还是换算法。

\noindent\textbf{LeetCode 1820 最多邀请的个数。}

\textbf{题目说明。} LeetCode 1820 最多邀请的个数 的题面入口要先还原成具体任务：男生和女生构成二分图，目标是最大一对一匹配数量。

\textbf{样例入口。} 拿 grid=[[1, 1, 1], [1, 0, 1], [0, 0, 1]] 手算，若第一个男生先匹配第一个女生，后面男生仍可能通过重新分配得到三对。贪心固定早期选择会挡住后续唯一选择。

\textbf{暴力基线。} 慢版本先按题意画出完整状态图，用普通搜索跑通可达性。这个过程用于确认男生和女生构成二分图，目标是最大一对一匹配数量的节点和边，之后再讨论访问标记、层次或代价顺序。放到本题就是：先按“男生和女生构成二分图，目标是最大一对一匹配数量”完整试慢版本；样例里已经能看到“规律是每次寻找增广路，找到未匹配右点后翻转路径上的匹配关系，使匹配数增加一”，所以优化前必须先能说清慢版本枚举了哪些候选。

\textbf{暴力过程。} 拿 grid=[[1, 1, 1], [1, 0, 1], [0, 0, 1]] 手算，若第一个男生先匹配第一个女生，后面男生仍可能通过重新分配得到三对。贪心固定早期选择会挡住后续唯一选择。

\textbf{重复工作。} 题面模型：男生和女生构成二分图，目标是最大一对一匹配数量。暴力版的慢点要落到本题：慢版本会围绕“男生和女生构成二分图，目标是最大一对一匹配数量”反复展开同一批状态。样例规律是“规律是每次寻找增广路，找到未匹配右点后翻转路径上的匹配关系，使匹配数增加一”，说明必须把节点、边、访问标记和资源维度一次建清；同一状态不应重复入队，不同资源状态也不能误合并。后面的优化只消掉这类重复工作，不能改变答案集合。

\textbf{优化条件。} 本题的关键观察是：对每个未匹配左点寻找增广路，找到未匹配右点后翻转路径，使匹配数加一。这句话必须能翻译成可维护状态；如果题目条件改变后观察不再成立，就不能继续套原来的写法。

\textbf{相邻状态。} 规律是每次寻找增广路，找到未匹配右点后翻转路径上的匹配关系，使匹配数增加一。把这个特例继续拆开看：状态节点不一定等于原始位置，还可能包含钥匙、剩余资源、时间或访问集合。visited 只能合并未来能力完全相同的状态，否则会把合法路径剪掉。

\textbf{状态复用。} 把重复搜索压成访问标记、距离数组或拓扑状态；邻居生成也要从全量扫描变成结构化枚举。本题落点是：规律是每次寻找增广路，找到未匹配右点后翻转路径上的匹配关系，使匹配数增加一。手算时只改被新输入影响的那一小部分，而不是回头重算完整候选。

\textbf{指针和字段。} 状态字段要能说清：状态编号、邻接生成方式、访问标记、距离或层数、队列内容；若状态含资源，visited 不能只按位置记录。手算时按这个协议跟踪：拿 grid=[[1, 1, 1], [1, 0, 1], [0, 0, 1]] 手算，若第一个男生先匹配第一个女生，后面男生仍可能通过重新分配得到三对。贪心固定早期选择会挡住后续唯一选择。然后按协议复查：手算时画队列或栈，状态必须包含题目需要的所有资源。入队时标记还是出队时标记，要和去重语义一致。

\textbf{代码落点。} 入队时标记还是出队时标记必须一致；方向数组集中定义；多源 BFS 要一次性把所有源点放入初始队列。关键代码行必须能逐句翻译成“旧状态怎样变成新状态”，否则就只是模板。

\textbf{正确性。} 证明从可达性开始：搜索覆盖所有合法边能到达的状态，访问标记只删除重复状态而不删除不同未来能力。

\textbf{边界实验。} 先用某一侧为空、多个男生争同一女生、贪心失败、访问标记按轮重置、稀疏边做反例检查。实验时覆盖某一侧为空、多个男生争同一女生、贪心失败、访问标记按轮重置、稀疏边，并打印入队时的完整状态。若同一位置但资源不同的状态被合并，很多图题会出现隐蔽错误。

\textbf{复杂度变化。} 暴力重复从多个起点搜索会反复遍历图；正确建模和访问标记后，BFS 或 DFS 通常按节点数加边数计算，状态图题则按完整状态数计算。

\textbf{迁移追问。} 如果题目改成“若边有容量或权重，升级为最大流或带权匹配”，先判断原状态是否仍然覆盖新答案集合，再决定是微调边界、改 value 语义，还是换算法。

\subsection{实验复盘}

追问通常会问节点和边的建模、访问标记时机、BFS 为什么最短。请准备说明状态是否还包含资源，为什么不能只按位置去重。

复盘本节时先从 LeetCode 126 单词接龙 II、LeetCode 127 单词接龙、LeetCode 130 被围绕的区域 开始：每道题都先手写一个能覆盖全部候选的暴力版，再指出优化结构具体保存了哪一段历史信息，最后用相邻案例比较为什么不能互换解法。

\topic{最短路与带权图}
这一节先看 LeetCode 505 迷宫 II、LeetCode 743 网络延迟时间、LeetCode 778 水位上升的泳池中游泳 这几道 LeetCode 题，再整理同一题型的分流和推导。阅读顺序不是先背原理，而是先看案例的暴力瓶颈，再把重复出现的观察抽象成方法。

\subsection{原理和适用条件}

以 LeetCode 505 迷宫 II、LeetCode 743 网络延迟时间、LeetCode 778 水位上升的泳池中游泳 为主线：LeetCode 505 迷宫 II：模型是“小球一旦选择方向就会滚到墙前停止，边权是滚动经过的格子数。”，优化抓手是“邻居生成不是走一步，而是沿四个方向模拟滚动到停点，再用 Dijkstra 按累计距离松弛。”；LeetCode 743 网络延迟时间：模型是“有向正权图从起点到所有节点的最短路最大值。”，优化抓手是“Dijkstra 求每个节点最短距离，最后取最大，若有无穷则不可达。”；LeetCode 778 水位上升的泳池中游泳：模型是“到达某格子的代价是路径上最大高度，而不是高度和。”，优化抓手是“用 Dijkstra 按当前路径最大高度扩展，松弛时取 max 当前代价和邻格高度。”。这些案例共同推出的规律是：先确认目标是步数最少还是代价最小，再检查边权是否非负。非负权用 Dijkstra，等权用 BFS，限制边数用分层松弛或 Bellman-Ford 思想。 方法名只是复盘后的标签，真正要掌握的是从题面约束、暴力失败、状态设计到边界反例的推导链。

\begin{principlebox}{图示：从 LeetCode 案例到 最短路与带权图推导链}
\begin{tabular}{@{}p{0.18\linewidth}p{0.74\linewidth}@{}}
暴力基线 & 枚举候选、完整模拟或逐个检查，先保证覆盖全部可能答案。\\
瓶颈定位 & 路径数量可能指数级。最短路算法的核心是用当前已知最小代价组织扩展顺序，而不是盲目枚举路径。\\
结构选择 & 非负权图用 Dijkstra，小顶堆保存当前距离最小的候选；边权为一用 BFS；有负边要考虑 Bellman-Ford 或重新建模。路径代价若是最大边权，也可用 Dijkstra 或二分答案。\\
不变量 & Dijkstra 的不变量是：每次弹出的未过期节点，其距离已经是最终最短距离。过期堆元素必须跳过，否则会重复扩展旧状态。\\
代码落地 & 由于 \code{priority_queue} 没有 decrease-key，更新距离时压入新状态，弹出时比较距离数组。距离用 \code{long long} 更稳，图用邻接表保存 \code{(to, weight)}。\\
\end{tabular}
\end{principlebox}

\subsection{LeetCode 案例分流}

\begin{principlebox}{从 LeetCode 案例分流：最短路与带权图}
\textbf{标准非负最短路。}
代表案例：LeetCode 743 网络延迟时间、LeetCode 1514 概率最大路径。先看这些题的题面条件：边权非负且总代价是边权和。由此推出的处理方式是：Dijkstra 用小顶堆按当前最短距离扩展。风险点：堆中旧状态要跳过，概率问题可换成最大堆或负对数。

\textbf{路径代价非求和。}
代表案例：LeetCode 1631 最小体力消耗路径、LeetCode 778 水位上升的泳池中游泳。先看这些题的题面条件：路径代价是最大边、最小边或其他单调合成值。由此推出的处理方式是：只要代价沿路径扩展不变小，就可改造 Dijkstra。风险点：要重新证明弹出状态已经最优。

\textbf{0-1 和分层边权。}
代表案例：LeetCode 1368 使网格图至少有一条有效路径的最小代价、LeetCode 1293 网格障碍消除最短路。先看这些题的题面条件：边权只有零和一，或步数和资源层数有限。由此推出的处理方式是：0-1 BFS 用双端队列，或按层松弛状态。风险点：普通队列不再保证代价递增。

\textbf{二分可达。}
代表案例：LeetCode 1631 最小体力消耗路径、LeetCode 1102 得分最高路径。先看这些题的题面条件：阈值越大可用边越多，存在可达性单调边界。由此推出的处理方式是：二分阈值，每次用 BFS/DFS 检查起终点是否连通。风险点：检查函数只回答可达，不保留具体最短路。

\end{principlebox}

\subsection{案例矩阵}

本节案例覆盖 LeetCode 505 迷宫 II、LeetCode 743 网络延迟时间、LeetCode 778 水位上升的泳池中游泳、LeetCode 787 K 站中转内最便宜的航班、LeetCode 1102 得分最高的路径、LeetCode 1334 阈值距离内邻居最少的城市、LeetCode 1368 使网格图至少有一条有效路径的最小代价、LeetCode 1514 概率最大的路径、LeetCode 1631 最小体力消耗路径、LeetCode 1928 规定时间内到达终点的最小花费、LeetCode 1976 到达目的地的方案数、LeetCode 2045 到达目的地的第二短时间。阅读时不要按题号背答案，而要先判断题面落在哪个分流场景：查询目标是什么，输入约束给了什么单调性或等价关系，暴力慢在哪里，优化结构保存了什么状态。

\begin{principlebox}{案例矩阵}
\textbf{LeetCode 505 迷宫 II。}
小球一旦选择方向就会滚到墙前停止，边权是滚动经过的格子数。单点风险：起点就是终点、目标无法停下、绕远路更短、同一停点多次被松弛。

\textbf{LeetCode 743 网络延迟时间。}
有向正权图从起点到所有节点的最短路最大值。单点风险：节点编号从一开始、不可达、重边、过期堆元素。

\textbf{LeetCode 778 水位上升的泳池中游泳。}
到达某格子的代价是路径上最大高度，而不是高度和。单点风险：单格、起点高度很高、需要绕路、多个相同高度。

\textbf{LeetCode 787 K 站中转内最便宜的航班。}
路径代价受中转次数限制，不能只按费用最小弹出节点后永久确定。单点风险：起点终点相同、不可达、K 为零、便宜路径中转过多。

\textbf{LeetCode 1102 得分最高的路径。}
路径得分是沿途最小格子值，要最大化这个最小值。单点风险：单格、起点或终点很小、必须绕路、多个相同得分路径。

\textbf{LeetCode 1334 阈值距离内邻居最少的城市。}
每个城市都需要知道到其他城市的最短距离是否不超过阈值。单点风险：距离等于阈值、并列时选编号最大、不可达、无向边。

\textbf{LeetCode 1368 使网格图至少有一条有效路径的最小代价。}
沿着格子箭头走代价为零，改方向走代价为一。单点风险：单格、箭头指向越界、多个零代价路径、过期状态。

\textbf{LeetCode 1514 概率最大的路径。}
路径概率是边概率乘积，目标是最大乘积路径。单点风险：不可达、概率为零、起点等于终点、浮点比较。

\textbf{LeetCode 1631 最小体力消耗路径。}
路径代价是沿途最大高度差，不是边权和。单点风险：单格、平坦网格、必须绕路、多个相同代价。

\textbf{LeetCode 1928 规定时间内到达终点的最小花费。}
路径同时有时间约束和费用目标，单一距离无法描述状态。单点风险：时间刚好到达、费用高但时间短、不可达、状态数量由 maxTime 决定。

\textbf{LeetCode 1976 到达目的地的方案数。}
需要在最短距离之外统计有多少条最短路径。单点风险：多条等长路径、取模、旧堆状态、不可达。

\textbf{LeetCode 2045 到达目的地的第二短时间。}
每个节点需要保留最短和次短到达步数，再考虑红绿灯等待时间。单点风险：第二短不能等于最短、红灯等待、无向图往返、到达终点后继续找次短。

\end{principlebox}

\subsection{共性落点}

\begin{principlebox}{本组共性落点}
\textbf{慢版本。} 暴力做法可能枚举所有路径，或者把带权边错误当成普通 BFS。只要边权不全相等，按步数最少就不等于总代价最小。
\textbf{瓶颈。} 路径数量可能指数级。最短路算法的核心是用当前已知最小代价组织扩展顺序，而不是盲目枚举路径。
\textbf{状态字段。} 距离数组、优先队列状态、松弛候选、旧状态判定、不可达哨兵；资源约束题还要把资源维度放进状态。
\textbf{不变量。} Dijkstra 的不变量是：每次弹出的未过期节点，其距离已经是最终最短距离。过期堆元素必须跳过，否则会重复扩展旧状态。
\textbf{实现检查。} 松弛时只在候选更优时更新；priority\_queue 弹出旧状态要跳过；距离加法用更宽整数。
\textbf{C++ 落地。} 由于 \code{priority_queue} 没有 decrease-key，更新距离时压入新状态，弹出时比较距离数组。距离用 \code{long long} 更稳，图用邻接表保存 \code{(to, weight)}。
\textbf{证明抓手。} 证明从扩展顺序开始：当前弹出的状态在代价规则下已经不会被未来路径改得更优。
\end{principlebox}

\subsection{案例推导}

\noindent\textbf{LeetCode 505 迷宫 II。}

\textbf{题目说明。} LeetCode 505 迷宫 II 的题面入口要先还原成具体任务：小球一旦选择方向就会滚到墙前停止，边权是滚动经过的格子数。

\textbf{样例入口。} 拿迷宫中球从 (0, 4) 滚向左手算，它不会每格停下，只会直到撞墙才停。普通 BFS 按格走会错，边权是滚动距离。

\textbf{暴力基线。} 慢版本可以枚举路径，或先用普通队列试跑一个小图。它会暴露步数最少和代价最小是否一致；一旦不一致，就必须围绕代价组织状态。放到本题就是：先按“小球一旦选择方向就会滚到墙前停止，边权是滚动经过的格子数”完整试慢版本；样例里已经能看到“规律是状态是停球位置，邻居是四个方向滚到的终点，距离要用 Dijkstra 更新”，所以优化前必须先能说清慢版本枚举了哪些候选。

\textbf{暴力过程。} 拿迷宫中球从 (0, 4) 滚向左手算，它不会每格停下，只会直到撞墙才停。普通 BFS 按格走会错，边权是滚动距离。

\textbf{重复工作。} 题面模型：小球一旦选择方向就会滚到墙前停止，边权是滚动经过的格子数。暴力版的慢点要落到本题：慢版本会围绕“小球一旦选择方向就会滚到墙前停止，边权是滚动经过的格子数”枚举路径或按错误顺序扩展状态。样例规律是“规律是状态是停球位置，邻居是四个方向滚到的终点，距离要用 Dijkstra 更新”，说明队列或堆的弹出顺序必须和代价定义一致，旧状态为什么能跳过也要讲清。后面的优化只消掉这类重复工作，不能改变答案集合。

\textbf{优化条件。} 本题的关键观察是：邻居生成不是走一步，而是沿四个方向模拟滚动到停点，再用 Dijkstra 按累计距离松弛。这句话必须能翻译成可维护状态；如果题目条件改变后观察不再成立，就不能继续套原来的写法。

\textbf{相邻状态。} 规律是状态是停球位置，邻居是四个方向滚到的终点，距离要用 Dijkstra 更新。把这个特例继续拆开看：队列或堆弹出的顺序必须和代价规则一致；旧状态被跳过，是因为已经有更小代价到达同一状态。只要边权或代价合成方式改变，就要重新证明扩展顺序。

\textbf{状态复用。} 把路径枚举压成距离数组和优先队列；每次只扩展当前最有希望的未过期状态。本题落点是：规律是状态是停球位置，邻居是四个方向滚到的终点，距离要用 Dijkstra 更新。手算时只改被新输入影响的那一小部分，而不是回头重算完整候选。

\textbf{指针和字段。} 状态字段要能说清：距离数组、优先队列状态、松弛候选、旧状态判定、不可达哨兵；资源约束题还要把资源维度放进状态。手算时按这个协议跟踪：拿迷宫中球从 (0, 4) 滚向左手算，它不会每格停下，只会直到撞墙才停。普通 BFS 按格走会错，边权是滚动距离。然后按协议复查：手算时记录每次弹出的代价、当前距离数组和松弛结果。旧状态被弹出时为什么跳过，也要写在表里。

\textbf{代码落点。} 松弛时只在候选更优时更新；priority\_queue 弹出旧状态要跳过；距离加法用更宽整数。关键代码行必须能逐句翻译成“旧状态怎样变成新状态”，否则就只是模板。

\textbf{正确性。} 证明从扩展顺序开始：当前弹出的状态在代价规则下已经不会被未来路径改得更优。

\textbf{边界实验。} 先用起点就是终点、目标无法停下、绕远路更短、同一停点多次被松弛做反例检查。实验时覆盖起点就是终点、目标无法停下、绕远路更短、同一停点多次被松弛，打印每次弹出的代价和是否过期。若普通 BFS 与 Dijkstra 结果不同，要能解释边权条件造成的差异。

\textbf{复杂度变化。} 暴力枚举路径可能指数级；Dijkstra 在邻接表和堆实现下按边数、点数和堆操作计算，0-1 BFS 可按节点数加边数计算。

\textbf{迁移追问。} 如果题目改成“若只问能否到达，可用 BFS 或 DFS；若问最少距离，就必须处理带权边”，先判断原状态是否仍然覆盖新答案集合，再决定是微调边界、改 value 语义，还是换算法。

\noindent\textbf{LeetCode 743 网络延迟时间。}

\textbf{题目说明。} LeetCode 743 网络延迟时间 的题面入口要先还原成具体任务：有向正权图从起点到所有节点的最短路最大值。

\textbf{样例入口。} 拿 times=[[2, 1, 1], [2, 3, 1], [3, 4, 1]]、起点 2 手算，2 到 1 和 3 距离都是 1，到 4 距离是 2。

\textbf{暴力基线。} 暴力枚举从起点到各点的所有路径，或错误地用普通 BFS 按边数扩展。

\textbf{暴力过程。} 以 times=[[2, 1, 1], [2, 3, 1], [3, 4, 1]]、起点 2 为例，2 到 1 和 3 的距离是 1，到 4 的距离是 2。

\textbf{重复工作。} 题面模型：有向正权图从起点到所有节点的最短路最大值。暴力版的慢点要落到本题：浪费在路径数量可能爆炸，而且普通队列不能保证带权路径的总代价最小。需要每次扩展当前距离最小的候选。后面的优化只消掉这类重复工作，不能改变答案集合。

\textbf{优化条件。} 本题的关键观察是：Dijkstra 求每个节点最短距离，最后取最大，若有无穷则不可达。这句话必须能翻译成可维护状态；如果题目条件改变后观察不再成立，就不能继续套原来的写法。

\textbf{相邻状态。} 规律是非负边最短路用 Dijkstra，最后取所有最短距离的最大值；不可达节点导致答案 -1。把这个特例继续拆开看：队列或堆弹出的顺序必须和代价规则一致；旧状态被跳过，是因为已经有更小代价到达同一状态。只要边权或代价合成方式改变，就要重新证明扩展顺序。

\textbf{状态复用。} Dijkstra：dist[node] 保存当前最短距离，小顶堆保存 (distance, node)。弹出过期距离时跳过，松弛所有出边。手算时只改被新输入影响的那一小部分，而不是回头重算完整候选。

\textbf{指针和字段。} 状态字段要能说清：dist[u] 是已知最短上界，heap top 是当前最小候选，edge weight 是从 u 到 v 的传播时间。手算时按这个协议跟踪：起点 2 距离 0 入堆；弹出 2 后更新 1 和 3 为 1；弹出 3 后更新 4 为 2。

\textbf{代码落点。} 关键是 priority\_queue 默认大顶堆，C++ 要用 greater 或存负距离；弹出时若 d!=dist[u] 就跳过旧状态。关键代码行必须能逐句翻译成“旧状态怎样变成新状态”，否则就只是模板。

\textbf{正确性。} 边权非负时，堆中最小距离节点一旦弹出，任何绕路都不会得到更小距离，因此它的 dist 已最终确定。

\textbf{边界实验。} 先用节点编号从一开始、不可达、重边、过期堆元素做反例检查。实验时覆盖节点编号从一开始、不可达、重边、过期堆元素，打印每次弹出的代价和是否过期。若普通 BFS 与 Dijkstra 结果不同，要能解释边权条件造成的差异。

\textbf{复杂度变化。} 枚举路径指数级；Dijkstra O((V+E)log V) 时间，O(V+E) 空间。

\textbf{迁移追问。} 如果题目改成“边权相同才可用普通 BFS”，先判断原状态是否仍然覆盖新答案集合，再决定是微调边界、改 value 语义，还是换算法。

\noindent\textbf{LeetCode 778 水位上升的泳池中游泳。}

\textbf{题目说明。} LeetCode 778 水位上升的泳池中游泳 的题面入口要先还原成具体任务：到达某格子的代价是路径上最大高度，而不是高度和。

\textbf{样例入口。} 拿 [[0, 2], [1, 3]] 手算，时间必须至少到 3 才能进入右下角；路径代价不是边权和，而是沿路最大高度。

\textbf{暴力基线。} 慢版本可以枚举路径，或先用普通队列试跑一个小图。它会暴露步数最少和代价最小是否一致；一旦不一致，就必须围绕代价组织状态。放到本题就是：先按“到达某格子的代价是路径上最大高度，而不是高度和”完整试慢版本；样例里已经能看到“规律是 Dijkstra 的距离定义为到达某格所需的最小最大高度，或二分时间做可达性”，所以优化前必须先能说清慢版本枚举了哪些候选。

\textbf{暴力过程。} 拿 [[0, 2], [1, 3]] 手算，时间必须至少到 3 才能进入右下角；路径代价不是边权和，而是沿路最大高度。

\textbf{重复工作。} 题面模型：到达某格子的代价是路径上最大高度，而不是高度和。暴力版的慢点要落到本题：慢版本会围绕“到达某格子的代价是路径上最大高度，而不是高度和”枚举路径或按错误顺序扩展状态。样例规律是“规律是 Dijkstra 的距离定义为到达某格所需的最小最大高度，或二分时间做可达性”，说明队列或堆的弹出顺序必须和代价定义一致，旧状态为什么能跳过也要讲清。后面的优化只消掉这类重复工作，不能改变答案集合。

\textbf{优化条件。} 本题的关键观察是：用 Dijkstra 按当前路径最大高度扩展，松弛时取 max 当前代价和邻格高度。这句话必须能翻译成可维护状态；如果题目条件改变后观察不再成立，就不能继续套原来的写法。

\textbf{相邻状态。} 规律是 Dijkstra 的距离定义为到达某格所需的最小最大高度，或二分时间做可达性。把这个特例继续拆开看：队列或堆弹出的顺序必须和代价规则一致；旧状态被跳过，是因为已经有更小代价到达同一状态。只要边权或代价合成方式改变，就要重新证明扩展顺序。

\textbf{状态复用。} 把路径枚举压成距离数组和优先队列；每次只扩展当前最有希望的未过期状态。本题落点是：规律是 Dijkstra 的距离定义为到达某格所需的最小最大高度，或二分时间做可达性。手算时只改被新输入影响的那一小部分，而不是回头重算完整候选。

\textbf{指针和字段。} 状态字段要能说清：距离数组、优先队列状态、松弛候选、旧状态判定、不可达哨兵；资源约束题还要把资源维度放进状态。手算时按这个协议跟踪：拿 [[0, 2], [1, 3]] 手算，时间必须至少到 3 才能进入右下角；路径代价不是边权和，而是沿路最大高度。然后按协议复查：手算时记录每次弹出的代价、当前距离数组和松弛结果。旧状态被弹出时为什么跳过，也要写在表里。

\textbf{代码落点。} 松弛时只在候选更优时更新；priority\_queue 弹出旧状态要跳过；距离加法用更宽整数。关键代码行必须能逐句翻译成“旧状态怎样变成新状态”，否则就只是模板。

\textbf{正确性。} 证明从扩展顺序开始：当前弹出的状态在代价规则下已经不会被未来路径改得更优。

\textbf{边界实验。} 先用单格、起点高度很高、需要绕路、多个相同高度做反例检查。实验时覆盖单格、起点高度很高、需要绕路、多个相同高度，打印每次弹出的代价和是否过期。若普通 BFS 与 Dijkstra 结果不同，要能解释边权条件造成的差异。

\textbf{复杂度变化。} 暴力枚举路径可能指数级；Dijkstra 在邻接表和堆实现下按边数、点数和堆操作计算，0-1 BFS 可按节点数加边数计算。

\textbf{迁移追问。} 如果题目改成“也可以按水位二分可达，或按高度排序并查集合并”，先判断原状态是否仍然覆盖新答案集合，再决定是微调边界、改 value 语义，还是换算法。

\noindent\textbf{LeetCode 787 K 站中转内最便宜的航班。}

\textbf{题目说明。} LeetCode 787 K 站中转内最便宜的航班 的题面入口要先还原成具体任务：路径代价受中转次数限制，不能只按费用最小弹出节点后永久确定。

\textbf{样例入口。} 拿航班 0->1 价格 100、1->2 价格 100、0->2 价格 500，K=1 手算，允许一次中转时 0->1->2 更便宜。

\textbf{暴力基线。} 慢版本可以枚举路径，或先用普通队列试跑一个小图。它会暴露步数最少和代价最小是否一致；一旦不一致，就必须围绕代价组织状态。放到本题就是：先按“路径代价受中转次数限制，不能只按费用最小弹出节点后永久确定”完整试慢版本；样例里已经能看到“规律是状态必须包含已用中转次数或边数，不能只保存每个城市一个最低价格”，所以优化前必须先能说清慢版本枚举了哪些候选。

\textbf{暴力过程。} 拿航班 0->1 价格 100、1->2 价格 100、0->2 价格 500，K=1 手算，允许一次中转时 0->1->2 更便宜。

\textbf{重复工作。} 题面模型：路径代价受中转次数限制，不能只按费用最小弹出节点后永久确定。暴力版的慢点要落到本题：慢版本会围绕“路径代价受中转次数限制，不能只按费用最小弹出节点后永久确定”枚举路径或按错误顺序扩展状态。样例规律是“规律是状态必须包含已用中转次数或边数，不能只保存每个城市一个最低价格”，说明队列或堆的弹出顺序必须和代价定义一致，旧状态为什么能跳过也要讲清。后面的优化只消掉这类重复工作，不能改变答案集合。

\textbf{优化条件。} 本题的关键观察是：用按边数分层的 Bellman-Ford 或带层数状态的优先队列，状态包含城市和已用边数。这句话必须能翻译成可维护状态；如果题目条件改变后观察不再成立，就不能继续套原来的写法。

\textbf{相邻状态。} 规律是状态必须包含已用中转次数或边数，不能只保存每个城市一个最低价格。把这个特例继续拆开看：队列或堆弹出的顺序必须和代价规则一致；旧状态被跳过，是因为已经有更小代价到达同一状态。只要边权或代价合成方式改变，就要重新证明扩展顺序。

\textbf{状态复用。} 把路径枚举压成距离数组和优先队列；每次只扩展当前最有希望的未过期状态。本题落点是：规律是状态必须包含已用中转次数或边数，不能只保存每个城市一个最低价格。手算时只改被新输入影响的那一小部分，而不是回头重算完整候选。

\textbf{指针和字段。} 状态字段要能说清：距离数组、优先队列状态、松弛候选、旧状态判定、不可达哨兵；资源约束题还要把资源维度放进状态。手算时按这个协议跟踪：拿航班 0->1 价格 100、1->2 价格 100、0->2 价格 500，K=1 手算，允许一次中转时 0->1->2 更便宜。然后按协议复查：手算时记录每次弹出的代价、当前距离数组和松弛结果。旧状态被弹出时为什么跳过，也要写在表里。

\textbf{代码落点。} 松弛时只在候选更优时更新；priority\_queue 弹出旧状态要跳过；距离加法用更宽整数。关键代码行必须能逐句翻译成“旧状态怎样变成新状态”，否则就只是模板。

\textbf{正确性。} 证明从扩展顺序开始：当前弹出的状态在代价规则下已经不会被未来路径改得更优。

\textbf{边界实验。} 先用起点终点相同、不可达、K 为零、便宜路径中转过多做反例检查。实验时覆盖起点终点相同、不可达、K 为零、便宜路径中转过多，打印每次弹出的代价和是否过期。若普通 BFS 与 Dijkstra 结果不同，要能解释边权条件造成的差异。

\textbf{复杂度变化。} 暴力枚举路径可能指数级；Dijkstra 在邻接表和堆实现下按边数、点数和堆操作计算，0-1 BFS 可按节点数加边数计算。

\textbf{迁移追问。} 如果题目改成“若去掉中转限制，普通 Dijkstra 才能直接按费用组织扩展”，先判断原状态是否仍然覆盖新答案集合，再决定是微调边界、改 value 语义，还是换算法。

\noindent\textbf{LeetCode 1102 得分最高的路径。}

\textbf{题目说明。} LeetCode 1102 得分最高的路径 的题面入口要先还原成具体任务：路径得分是沿途最小格子值，要最大化这个最小值。

\textbf{样例入口。} 拿 [[5, 4, 5], [1, 2, 6], [7, 4, 6]] 手算，路径分数是沿路最小值，要最大化这个最小值。

\textbf{暴力基线。} 慢版本可以枚举路径，或先用普通队列试跑一个小图。它会暴露步数最少和代价最小是否一致；一旦不一致，就必须围绕代价组织状态。放到本题就是：先按“路径得分是沿途最小格子值，要最大化这个最小值”完整试慢版本；样例里已经能看到“规律是优先扩展当前路径分数最高的格子，或二分阈值判断能否只走不低于阈值的格子”，所以优化前必须先能说清慢版本枚举了哪些候选。

\textbf{暴力过程。} 拿 [[5, 4, 5], [1, 2, 6], [7, 4, 6]] 手算，路径分数是沿路最小值，要最大化这个最小值。

\textbf{重复工作。} 题面模型：路径得分是沿途最小格子值，要最大化这个最小值。暴力版的慢点要落到本题：慢版本会围绕“路径得分是沿途最小格子值，要最大化这个最小值”枚举路径或按错误顺序扩展状态。样例规律是“规律是优先扩展当前路径分数最高的格子，或二分阈值判断能否只走不低于阈值的格子”，说明队列或堆的弹出顺序必须和代价定义一致，旧状态为什么能跳过也要讲清。后面的优化只消掉这类重复工作，不能改变答案集合。

\textbf{优化条件。} 本题的关键观察是：用最大堆按当前路径得分扩展，进入邻格时取当前得分和邻格值的较小者。这句话必须能翻译成可维护状态；如果题目条件改变后观察不再成立，就不能继续套原来的写法。

\textbf{相邻状态。} 规律是优先扩展当前路径分数最高的格子，或二分阈值判断能否只走不低于阈值的格子。把这个特例继续拆开看：队列或堆弹出的顺序必须和代价规则一致；旧状态被跳过，是因为已经有更小代价到达同一状态。只要边权或代价合成方式改变，就要重新证明扩展顺序。

\textbf{状态复用。} 把路径枚举压成距离数组和优先队列；每次只扩展当前最有希望的未过期状态。本题落点是：规律是优先扩展当前路径分数最高的格子，或二分阈值判断能否只走不低于阈值的格子。手算时只改被新输入影响的那一小部分，而不是回头重算完整候选。

\textbf{指针和字段。} 状态字段要能说清：距离数组、优先队列状态、松弛候选、旧状态判定、不可达哨兵；资源约束题还要把资源维度放进状态。手算时按这个协议跟踪：拿 [[5, 4, 5], [1, 2, 6], [7, 4, 6]] 手算，路径分数是沿路最小值，要最大化这个最小值。然后按协议复查：手算时记录每次弹出的代价、当前距离数组和松弛结果。旧状态被弹出时为什么跳过，也要写在表里。

\textbf{代码落点。} 松弛时只在候选更优时更新；priority\_queue 弹出旧状态要跳过；距离加法用更宽整数。关键代码行必须能逐句翻译成“旧状态怎样变成新状态”，否则就只是模板。

\textbf{正确性。} 证明从扩展顺序开始：当前弹出的状态在代价规则下已经不会被未来路径改得更优。

\textbf{边界实验。} 先用单格、起点或终点很小、必须绕路、多个相同得分路径做反例检查。实验时覆盖单格、起点或终点很小、必须绕路、多个相同得分路径，打印每次弹出的代价和是否过期。若普通 BFS 与 Dijkstra 结果不同，要能解释边权条件造成的差异。

\textbf{复杂度变化。} 暴力枚举路径可能指数级；Dijkstra 在邻接表和堆实现下按边数、点数和堆操作计算，0-1 BFS 可按节点数加边数计算。

\textbf{迁移追问。} 如果题目改成“也可把阈值二分后检查连通性，或按值从高到低并查集合并”，先判断原状态是否仍然覆盖新答案集合，再决定是微调边界、改 value 语义，还是换算法。

\noindent\textbf{LeetCode 1334 阈值距离内邻居最少的城市。}

\textbf{题目说明。} LeetCode 1334 阈值距离内邻居最少的城市 的题面入口要先还原成具体任务：每个城市都需要知道到其他城市的最短距离是否不超过阈值。

\textbf{样例入口。} 拿四个城市和距离阈值 4 手算，每个城市需要统计最短距离不超过 4 的其他城市数；并列时选编号更大的城市。

\textbf{暴力基线。} 慢版本可以枚举路径，或先用普通队列试跑一个小图。它会暴露步数最少和代价最小是否一致；一旦不一致，就必须围绕代价组织状态。放到本题就是：先按“每个城市都需要知道到其他城市的最短距离是否不超过阈值”完整试慢版本；样例里已经能看到“规律是先求全源最短路，再按邻居数量和编号规则选答案”，所以优化前必须先能说清慢版本枚举了哪些候选。

\textbf{暴力过程。} 拿四个城市和距离阈值 4 手算，每个城市需要统计最短距离不超过 4 的其他城市数；并列时选编号更大的城市。

\textbf{重复工作。} 题面模型：每个城市都需要知道到其他城市的最短距离是否不超过阈值。暴力版的慢点要落到本题：慢版本会围绕“每个城市都需要知道到其他城市的最短距离是否不超过阈值”枚举路径或按错误顺序扩展状态。样例规律是“规律是先求全源最短路，再按邻居数量和编号规则选答案”，说明队列或堆的弹出顺序必须和代价定义一致，旧状态为什么能跳过也要讲清。后面的优化只消掉这类重复工作，不能改变答案集合。

\textbf{优化条件。} 本题的关键观察是：节点数较小时可用 Floyd，边稀疏时也可从每个点跑 Dijkstra。这句话必须能翻译成可维护状态；如果题目条件改变后观察不再成立，就不能继续套原来的写法。

\textbf{相邻状态。} 规律是先求全源最短路，再按邻居数量和编号规则选答案。把这个特例继续拆开看：队列或堆弹出的顺序必须和代价规则一致；旧状态被跳过，是因为已经有更小代价到达同一状态。只要边权或代价合成方式改变，就要重新证明扩展顺序。

\textbf{状态复用。} 把路径枚举压成距离数组和优先队列；每次只扩展当前最有希望的未过期状态。本题落点是：规律是先求全源最短路，再按邻居数量和编号规则选答案。手算时只改被新输入影响的那一小部分，而不是回头重算完整候选。

\textbf{指针和字段。} 状态字段要能说清：距离数组、优先队列状态、松弛候选、旧状态判定、不可达哨兵；资源约束题还要把资源维度放进状态。手算时按这个协议跟踪：拿四个城市和距离阈值 4 手算，每个城市需要统计最短距离不超过 4 的其他城市数；并列时选编号更大的城市。然后按协议复查：手算时记录每次弹出的代价、当前距离数组和松弛结果。旧状态被弹出时为什么跳过，也要写在表里。

\textbf{代码落点。} 松弛时只在候选更优时更新；priority\_queue 弹出旧状态要跳过；距离加法用更宽整数。关键代码行必须能逐句翻译成“旧状态怎样变成新状态”，否则就只是模板。

\textbf{正确性。} 证明从扩展顺序开始：当前弹出的状态在代价规则下已经不会被未来路径改得更优。

\textbf{边界实验。} 先用距离等于阈值、并列时选编号最大、不可达、无向边做反例检查。实验时覆盖距离等于阈值、并列时选编号最大、不可达、无向边，打印每次弹出的代价和是否过期。若普通 BFS 与 Dijkstra 结果不同，要能解释边权条件造成的差异。

\textbf{复杂度变化。} 暴力枚举路径可能指数级；Dijkstra 在邻接表和堆实现下按边数、点数和堆操作计算，0-1 BFS 可按节点数加边数计算。

\textbf{迁移追问。} 如果题目改成“这是多源最短路选择问题，算法取决于 n、边数和查询密度”，先判断原状态是否仍然覆盖新答案集合，再决定是微调边界、改 value 语义，还是换算法。

\noindent\textbf{LeetCode 1368 使网格图至少有一条有效路径的最小代价。}

\textbf{题目说明。} LeetCode 1368 使网格图至少有一条有效路径的最小代价 的题面入口要先还原成具体任务：沿着格子箭头走代价为零，改方向走代价为一。

\textbf{样例入口。} 拿网格箭头指向右和下手算，沿箭头走的代价是 0，改方向才付 1。

\textbf{暴力基线。} 慢版本可以枚举路径，或先用普通队列试跑一个小图。它会暴露步数最少和代价最小是否一致；一旦不一致，就必须围绕代价组织状态。放到本题就是：先按“沿着格子箭头走代价为零，改方向走代价为一”完整试慢版本；样例里已经能看到“规律是边权只有 0 和 1，可用 0-1 BFS；状态是格子，松弛代价取决于当前箭头是否指向邻居”，所以优化前必须先能说清慢版本枚举了哪些候选。

\textbf{暴力过程。} 拿网格箭头指向右和下手算，沿箭头走的代价是 0，改方向才付 1。

\textbf{重复工作。} 题面模型：沿着格子箭头走代价为零，改方向走代价为一。暴力版的慢点要落到本题：慢版本会围绕“沿着格子箭头走代价为零，改方向走代价为一”枚举路径或按错误顺序扩展状态。样例规律是“规律是边权只有 0 和 1，可用 0-1 BFS；状态是格子，松弛代价取决于当前箭头是否指向邻居”，说明队列或堆的弹出顺序必须和代价定义一致，旧状态为什么能跳过也要讲清。后面的优化只消掉这类重复工作，不能改变答案集合。

\textbf{优化条件。} 本题的关键观察是：边权只有零和一，用 0-1 BFS：零代价邻居入队首，一代价邻居入队尾。这句话必须能翻译成可维护状态；如果题目条件改变后观察不再成立，就不能继续套原来的写法。

\textbf{相邻状态。} 规律是边权只有 0 和 1，可用 0-1 BFS；状态是格子，松弛代价取决于当前箭头是否指向邻居。把这个特例继续拆开看：队列或堆弹出的顺序必须和代价规则一致；旧状态被跳过，是因为已经有更小代价到达同一状态。只要边权或代价合成方式改变，就要重新证明扩展顺序。

\textbf{状态复用。} 把路径枚举压成距离数组和优先队列；每次只扩展当前最有希望的未过期状态。本题落点是：规律是边权只有 0 和 1，可用 0-1 BFS；状态是格子，松弛代价取决于当前箭头是否指向邻居。手算时只改被新输入影响的那一小部分，而不是回头重算完整候选。

\textbf{指针和字段。} 状态字段要能说清：距离数组、优先队列状态、松弛候选、旧状态判定、不可达哨兵；资源约束题还要把资源维度放进状态。手算时按这个协议跟踪：拿网格箭头指向右和下手算，沿箭头走的代价是 0，改方向才付 1。然后按协议复查：手算时记录每次弹出的代价、当前距离数组和松弛结果。旧状态被弹出时为什么跳过，也要写在表里。

\textbf{代码落点。} 松弛时只在候选更优时更新；priority\_queue 弹出旧状态要跳过；距离加法用更宽整数。关键代码行必须能逐句翻译成“旧状态怎样变成新状态”，否则就只是模板。

\textbf{正确性。} 证明从扩展顺序开始：当前弹出的状态在代价规则下已经不会被未来路径改得更优。

\textbf{边界实验。} 先用单格、箭头指向越界、多个零代价路径、过期状态做反例检查。实验时覆盖单格、箭头指向越界、多个零代价路径、过期状态，打印每次弹出的代价和是否过期。若普通 BFS 与 Dijkstra 结果不同，要能解释边权条件造成的差异。

\textbf{复杂度变化。} 暴力枚举路径可能指数级；Dijkstra 在邻接表和堆实现下按边数、点数和堆操作计算，0-1 BFS 可按节点数加边数计算。

\textbf{迁移追问。} 如果题目改成“普通 BFS 按步数扩展，不等于总代价最小”，先判断原状态是否仍然覆盖新答案集合，再决定是微调边界、改 value 语义，还是换算法。

\noindent\textbf{LeetCode 1514 概率最大的路径。}

\textbf{题目说明。} LeetCode 1514 概率最大的路径 的题面入口要先还原成具体任务：路径概率是边概率乘积，目标是最大乘积路径。

\textbf{样例入口。} 拿 0->1 概率 0.5、1->2 概率 0.5、0->2 概率 0.2 手算，经 1 到 2 的概率 0.25 更大。

\textbf{暴力基线。} 慢版本可以枚举路径，或先用普通队列试跑一个小图。它会暴露步数最少和代价最小是否一致；一旦不一致，就必须围绕代价组织状态。放到本题就是：先按“路径概率是边概率乘积，目标是最大乘积路径”完整试慢版本；样例里已经能看到“规律是把距离改成最大概率，优先队列每次弹出当前概率最高的节点并做乘法松弛”，所以优化前必须先能说清慢版本枚举了哪些候选。

\textbf{暴力过程。} 拿 0->1 概率 0.5、1->2 概率 0.5、0->2 概率 0.2 手算，经 1 到 2 的概率 0.25 更大。

\textbf{重复工作。} 题面模型：路径概率是边概率乘积，目标是最大乘积路径。暴力版的慢点要落到本题：慢版本会围绕“路径概率是边概率乘积，目标是最大乘积路径”枚举路径或按错误顺序扩展状态。样例规律是“规律是把距离改成最大概率，优先队列每次弹出当前概率最高的节点并做乘法松弛”，说明队列或堆的弹出顺序必须和代价定义一致，旧状态为什么能跳过也要讲清。后面的优化只消掉这类重复工作，不能改变答案集合。

\textbf{优化条件。} 本题的关键观察是：用大顶堆每次扩展当前概率最高节点，乘以边概率得到候选概率。这句话必须能翻译成可维护状态；如果题目条件改变后观察不再成立，就不能继续套原来的写法。

\textbf{相邻状态。} 规律是把距离改成最大概率，优先队列每次弹出当前概率最高的节点并做乘法松弛。把这个特例继续拆开看：队列或堆弹出的顺序必须和代价规则一致；旧状态被跳过，是因为已经有更小代价到达同一状态。只要边权或代价合成方式改变，就要重新证明扩展顺序。

\textbf{状态复用。} 把路径枚举压成距离数组和优先队列；每次只扩展当前最有希望的未过期状态。本题落点是：规律是把距离改成最大概率，优先队列每次弹出当前概率最高的节点并做乘法松弛。手算时只改被新输入影响的那一小部分，而不是回头重算完整候选。

\textbf{指针和字段。} 状态字段要能说清：距离数组、优先队列状态、松弛候选、旧状态判定、不可达哨兵；资源约束题还要把资源维度放进状态。手算时按这个协议跟踪：拿 0->1 概率 0.5、1->2 概率 0.5、0->2 概率 0.2 手算，经 1 到 2 的概率 0.25 更大。然后按协议复查：手算时记录每次弹出的代价、当前距离数组和松弛结果。旧状态被弹出时为什么跳过，也要写在表里。

\textbf{代码落点。} 松弛时只在候选更优时更新；priority\_queue 弹出旧状态要跳过；距离加法用更宽整数。关键代码行必须能逐句翻译成“旧状态怎样变成新状态”，否则就只是模板。

\textbf{正确性。} 证明从扩展顺序开始：当前弹出的状态在代价规则下已经不会被未来路径改得更优。

\textbf{边界实验。} 先用不可达、概率为零、起点等于终点、浮点比较做反例检查。实验时覆盖不可达、概率为零、起点等于终点、浮点比较，打印每次弹出的代价和是否过期。若普通 BFS 与 Dijkstra 结果不同，要能解释边权条件造成的差异。

\textbf{复杂度变化。} 暴力枚举路径可能指数级；Dijkstra 在邻接表和堆实现下按边数、点数和堆操作计算，0-1 BFS 可按节点数加边数计算。

\textbf{迁移追问。} 如果题目改成“取负对数后可转成最短路，但直接最大堆更直观”，先判断原状态是否仍然覆盖新答案集合，再决定是微调边界、改 value 语义，还是换算法。

\noindent\textbf{LeetCode 1631 最小体力消耗路径。}

\textbf{题目说明。} LeetCode 1631 最小体力消耗路径 的题面入口要先还原成具体任务：路径代价是沿途最大高度差，不是边权和。

\textbf{样例入口。} 拿 heights=[[1, 2, 2], [3, 8, 2], [5, 3, 5]] 手算，路径代价是相邻高度差的最大值，不是差值之和。

\textbf{暴力基线。} 慢版本可以枚举路径，或先用普通队列试跑一个小图。它会暴露步数最少和代价最小是否一致；一旦不一致，就必须围绕代价组织状态。放到本题就是：先按“路径代价是沿途最大高度差，不是边权和”完整试慢版本；样例里已经能看到“规律是 Dijkstra 的距离定义为到达某格的最小最大边权，或二分最大体力做可达性”，所以优化前必须先能说清慢版本枚举了哪些候选。

\textbf{暴力过程。} 拿 heights=[[1, 2, 2], [3, 8, 2], [5, 3, 5]] 手算，路径代价是相邻高度差的最大值，不是差值之和。

\textbf{重复工作。} 题面模型：路径代价是沿途最大高度差，不是边权和。暴力版的慢点要落到本题：慢版本会围绕“路径代价是沿途最大高度差，不是边权和”枚举路径或按错误顺序扩展状态。样例规律是“规律是 Dijkstra 的距离定义为到达某格的最小最大边权，或二分最大体力做可达性”，说明队列或堆的弹出顺序必须和代价定义一致，旧状态为什么能跳过也要讲清。后面的优化只消掉这类重复工作，不能改变答案集合。

\textbf{优化条件。} 本题的关键观察是：Dijkstra 的松弛改成取当前代价和新边代价的最大值。这句话必须能翻译成可维护状态；如果题目条件改变后观察不再成立，就不能继续套原来的写法。

\textbf{相邻状态。} 规律是 Dijkstra 的距离定义为到达某格的最小最大边权，或二分最大体力做可达性。把这个特例继续拆开看：队列或堆弹出的顺序必须和代价规则一致；旧状态被跳过，是因为已经有更小代价到达同一状态。只要边权或代价合成方式改变，就要重新证明扩展顺序。

\textbf{状态复用。} 把路径枚举压成距离数组和优先队列；每次只扩展当前最有希望的未过期状态。本题落点是：规律是 Dijkstra 的距离定义为到达某格的最小最大边权，或二分最大体力做可达性。手算时只改被新输入影响的那一小部分，而不是回头重算完整候选。

\textbf{指针和字段。} 状态字段要能说清：距离数组、优先队列状态、松弛候选、旧状态判定、不可达哨兵；资源约束题还要把资源维度放进状态。手算时按这个协议跟踪：拿 heights=[[1, 2, 2], [3, 8, 2], [5, 3, 5]] 手算，路径代价是相邻高度差的最大值，不是差值之和。然后按协议复查：手算时记录每次弹出的代价、当前距离数组和松弛结果。旧状态被弹出时为什么跳过，也要写在表里。

\textbf{代码落点。} 松弛时只在候选更优时更新；priority\_queue 弹出旧状态要跳过；距离加法用更宽整数。关键代码行必须能逐句翻译成“旧状态怎样变成新状态”，否则就只是模板。

\textbf{正确性。} 证明从扩展顺序开始：当前弹出的状态在代价规则下已经不会被未来路径改得更优。

\textbf{边界实验。} 先用单格、平坦网格、必须绕路、多个相同代价做反例检查。实验时覆盖单格、平坦网格、必须绕路、多个相同代价，打印每次弹出的代价和是否过期。若普通 BFS 与 Dijkstra 结果不同，要能解释边权条件造成的差异。

\textbf{复杂度变化。} 暴力枚举路径可能指数级；Dijkstra 在邻接表和堆实现下按边数、点数和堆操作计算，0-1 BFS 可按节点数加边数计算。

\textbf{迁移追问。} 如果题目改成“也可二分体力上限后 BFS 检查可达”，先判断原状态是否仍然覆盖新答案集合，再决定是微调边界、改 value 语义，还是换算法。

\noindent\textbf{LeetCode 1928 规定时间内到达终点的最小花费。}

\textbf{题目说明。} LeetCode 1928 规定时间内到达终点的最小花费 的题面入口要先还原成具体任务：路径同时有时间约束和费用目标，单一距离无法描述状态。

\textbf{样例入口。} 拿 0->1 用时 10 费用 5，1->2 用时 10 费用 5，0->2 用时 25 费用 100、maxTime=30 手算，费用最小的是经 1 到 2。

\textbf{暴力基线。} 慢版本可以枚举路径，或先用普通队列试跑一个小图。它会暴露步数最少和代价最小是否一致；一旦不一致，就必须围绕代价组织状态。放到本题就是：先按“路径同时有时间约束和费用目标，单一距离无法描述状态”完整试慢版本；样例里已经能看到“规律是每个节点不能只保存一个最小费用，状态要包含已用时间”，所以优化前必须先能说清慢版本枚举了哪些候选。

\textbf{暴力过程。} 拿 0->1 用时 10 费用 5，1->2 用时 10 费用 5，0->2 用时 25 费用 100、maxTime=30 手算，费用最小的是经 1 到 2。

\textbf{重复工作。} 题面模型：路径同时有时间约束和费用目标，单一距离无法描述状态。暴力版的慢点要落到本题：慢版本会围绕“路径同时有时间约束和费用目标，单一距离无法描述状态”枚举路径或按错误顺序扩展状态。样例规律是“规律是每个节点不能只保存一个最小费用，状态要包含已用时间”，说明队列或堆的弹出顺序必须和代价定义一致，旧状态为什么能跳过也要讲清。后面的优化只消掉这类重复工作，不能改变答案集合。

\textbf{优化条件。} 本题的关键观察是：状态包含节点和已用时间，DP 或 Dijkstra 在时间维度内松弛最小费用。这句话必须能翻译成可维护状态；如果题目条件改变后观察不再成立，就不能继续套原来的写法。

\textbf{相邻状态。} 规律是每个节点不能只保存一个最小费用，状态要包含已用时间。把这个特例继续拆开看：队列或堆弹出的顺序必须和代价规则一致；旧状态被跳过，是因为已经有更小代价到达同一状态。只要边权或代价合成方式改变，就要重新证明扩展顺序。

\textbf{状态复用。} 把路径枚举压成距离数组和优先队列；每次只扩展当前最有希望的未过期状态。本题落点是：规律是每个节点不能只保存一个最小费用，状态要包含已用时间。手算时只改被新输入影响的那一小部分，而不是回头重算完整候选。

\textbf{指针和字段。} 状态字段要能说清：距离数组、优先队列状态、松弛候选、旧状态判定、不可达哨兵；资源约束题还要把资源维度放进状态。手算时按这个协议跟踪：拿 0->1 用时 10 费用 5，1->2 用时 10 费用 5，0->2 用时 25 费用 100、maxTime=30 手算，费用最小的是经 1 到 2。然后按协议复查：手算时记录每次弹出的代价、当前距离数组和松弛结果。旧状态被弹出时为什么跳过，也要写在表里。

\textbf{代码落点。} 松弛时只在候选更优时更新；priority\_queue 弹出旧状态要跳过；距离加法用更宽整数。关键代码行必须能逐句翻译成“旧状态怎样变成新状态”，否则就只是模板。

\textbf{正确性。} 证明从扩展顺序开始：当前弹出的状态在代价规则下已经不会被未来路径改得更优。

\textbf{边界实验。} 先用时间刚好到达、费用高但时间短、不可达、状态数量由 maxTime 决定做反例检查。实验时覆盖时间刚好到达、费用高但时间短、不可达、状态数量由 maxTime 决定，打印每次弹出的代价和是否过期。若普通 BFS 与 Dijkstra 结果不同，要能解释边权条件造成的差异。

\textbf{复杂度变化。} 暴力枚举路径可能指数级；Dijkstra 在邻接表和堆实现下按边数、点数和堆操作计算，0-1 BFS 可按节点数加边数计算。

\textbf{迁移追问。} 如果题目改成“这是资源约束最短路，不能只保存每个节点一个最小费用”，先判断原状态是否仍然覆盖新答案集合，再决定是微调边界、改 value 语义，还是换算法。

\noindent\textbf{LeetCode 1976 到达目的地的方案数。}

\textbf{题目说明。} LeetCode 1976 到达目的地的方案数 的题面入口要先还原成具体任务：需要在最短距离之外统计有多少条最短路径。

\textbf{样例入口。} 拿 0->1->3 和 0->2->3 两条长度都为 2 的路手算，最短距离相同，所以到 3 的方案数应累加为 2。

\textbf{暴力基线。} 慢版本可以枚举路径，或先用普通队列试跑一个小图。它会暴露步数最少和代价最小是否一致；一旦不一致，就必须围绕代价组织状态。放到本题就是：先按“需要在最短距离之外统计有多少条最短路径”完整试慢版本；样例里已经能看到“规律是 Dijkstra 松弛到更短距离时重置方案数，等长时累加方案数”，所以优化前必须先能说清慢版本枚举了哪些候选。

\textbf{暴力过程。} 拿 0->1->3 和 0->2->3 两条长度都为 2 的路手算，最短距离相同，所以到 3 的方案数应累加为 2。

\textbf{重复工作。} 题面模型：需要在最短距离之外统计有多少条最短路径。暴力版的慢点要落到本题：慢版本会围绕“需要在最短距离之外统计有多少条最短路径”枚举路径或按错误顺序扩展状态。样例规律是“规律是 Dijkstra 松弛到更短距离时重置方案数，等长时累加方案数”，说明队列或堆的弹出顺序必须和代价定义一致，旧状态为什么能跳过也要讲清。后面的优化只消掉这类重复工作，不能改变答案集合。

\textbf{优化条件。} 本题的关键观察是：Dijkstra 松弛时若得到更短距离就重置方案数，若等于最短距离就累加方案数。这句话必须能翻译成可维护状态；如果题目条件改变后观察不再成立，就不能继续套原来的写法。

\textbf{相邻状态。} 规律是 Dijkstra 松弛到更短距离时重置方案数，等长时累加方案数。把这个特例继续拆开看：队列或堆弹出的顺序必须和代价规则一致；旧状态被跳过，是因为已经有更小代价到达同一状态。只要边权或代价合成方式改变，就要重新证明扩展顺序。

\textbf{状态复用。} 把路径枚举压成距离数组和优先队列；每次只扩展当前最有希望的未过期状态。本题落点是：规律是 Dijkstra 松弛到更短距离时重置方案数，等长时累加方案数。手算时只改被新输入影响的那一小部分，而不是回头重算完整候选。

\textbf{指针和字段。} 状态字段要能说清：距离数组、优先队列状态、松弛候选、旧状态判定、不可达哨兵；资源约束题还要把资源维度放进状态。手算时按这个协议跟踪：拿 0->1->3 和 0->2->3 两条长度都为 2 的路手算，最短距离相同，所以到 3 的方案数应累加为 2。然后按协议复查：手算时记录每次弹出的代价、当前距离数组和松弛结果。旧状态被弹出时为什么跳过，也要写在表里。

\textbf{代码落点。} 松弛时只在候选更优时更新；priority\_queue 弹出旧状态要跳过；距离加法用更宽整数。关键代码行必须能逐句翻译成“旧状态怎样变成新状态”，否则就只是模板。

\textbf{正确性。} 证明从扩展顺序开始：当前弹出的状态在代价规则下已经不会被未来路径改得更优。

\textbf{边界实验。} 先用多条等长路径、取模、旧堆状态、不可达做反例检查。实验时覆盖多条等长路径、取模、旧堆状态、不可达，打印每次弹出的代价和是否过期。若普通 BFS 与 Dijkstra 结果不同，要能解释边权条件造成的差异。

\textbf{复杂度变化。} 暴力枚举路径可能指数级；Dijkstra 在邻接表和堆实现下按边数、点数和堆操作计算，0-1 BFS 可按节点数加边数计算。

\textbf{迁移追问。} 如果题目改成“这是最短路和计数 DP 的结合，更新顺序必须和距离松弛一致”，先判断原状态是否仍然覆盖新答案集合，再决定是微调边界、改 value 语义，还是换算法。

\noindent\textbf{LeetCode 2045 到达目的地的第二短时间。}

\textbf{题目说明。} LeetCode 2045 到达目的地的第二短时间 的题面入口要先还原成具体任务：每个节点需要保留最短和次短到达步数，再考虑红绿灯等待时间。

\textbf{样例入口。} 拿 1--2、time=3、change=2 手算，第一次到 2 是最短；第二短必须从 2 回 1 再到 2，且出发时可能遇到红灯等待。

\textbf{暴力基线。} 慢版本可以枚举路径，或先用普通队列试跑一个小图。它会暴露步数最少和代价最小是否一致；一旦不一致，就必须围绕代价组织状态。放到本题就是：先按“每个节点需要保留最短和次短到达步数，再考虑红绿灯等待时间”完整试慢版本；样例里已经能看到“规律是每个节点保存两个不同到达时间，不能把次短等同于最短”，所以优化前必须先能说清慢版本枚举了哪些候选。

\textbf{暴力过程。} 拿 1--2、time=3、change=2 手算，第一次到 2 是最短；第二短必须从 2 回 1 再到 2，且出发时可能遇到红灯等待。

\textbf{重复工作。} 题面模型：每个节点需要保留最短和次短到达步数，再考虑红绿灯等待时间。暴力版的慢点要落到本题：慢版本会围绕“每个节点需要保留最短和次短到达步数，再考虑红绿灯等待时间”枚举路径或按错误顺序扩展状态。样例规律是“规律是每个节点保存两个不同到达时间，不能把次短等同于最短”，说明队列或堆的弹出顺序必须和代价定义一致，旧状态为什么能跳过也要讲清。后面的优化只消掉这类重复工作，不能改变答案集合。

\textbf{优化条件。} 本题的关键观察是：BFS 式扩展边数或时间，只有候选时间不同且能进入前两名时才入队。这句话必须能翻译成可维护状态；如果题目条件改变后观察不再成立，就不能继续套原来的写法。

\textbf{相邻状态。} 规律是每个节点保存两个不同到达时间，不能把次短等同于最短。把这个特例继续拆开看：队列或堆弹出的顺序必须和代价规则一致；旧状态被跳过，是因为已经有更小代价到达同一状态。只要边权或代价合成方式改变，就要重新证明扩展顺序。

\textbf{状态复用。} 把路径枚举压成距离数组和优先队列；每次只扩展当前最有希望的未过期状态。本题落点是：规律是每个节点保存两个不同到达时间，不能把次短等同于最短。手算时只改被新输入影响的那一小部分，而不是回头重算完整候选。

\textbf{指针和字段。} 状态字段要能说清：距离数组、优先队列状态、松弛候选、旧状态判定、不可达哨兵；资源约束题还要把资源维度放进状态。手算时按这个协议跟踪：拿 1--2、time=3、change=2 手算，第一次到 2 是最短；第二短必须从 2 回 1 再到 2，且出发时可能遇到红灯等待。然后按协议复查：手算时记录每次弹出的代价、当前距离数组和松弛结果。旧状态被弹出时为什么跳过，也要写在表里。

\textbf{代码落点。} 松弛时只在候选更优时更新；priority\_queue 弹出旧状态要跳过；距离加法用更宽整数。关键代码行必须能逐句翻译成“旧状态怎样变成新状态”，否则就只是模板。

\textbf{正确性。} 证明从扩展顺序开始：当前弹出的状态在代价规则下已经不会被未来路径改得更优。

\textbf{边界实验。} 先用第二短不能等于最短、红灯等待、无向图往返、到达终点后继续找次短做反例检查。实验时覆盖第二短不能等于最短、红灯等待、无向图往返、到达终点后继续找次短，打印每次弹出的代价和是否过期。若普通 BFS 与 Dijkstra 结果不同，要能解释边权条件造成的差异。

\textbf{复杂度变化。} 暴力枚举路径可能指数级；Dijkstra 在邻接表和堆实现下按边数、点数和堆操作计算，0-1 BFS 可按节点数加边数计算。

\textbf{迁移追问。} 如果题目改成“它说明最短路状态有时不是一个最优值，而是前两个不同值”，先判断原状态是否仍然覆盖新答案集合，再决定是微调边界、改 value 语义，还是换算法。

\subsection{实验复盘}

追问通常会问为什么普通 BFS 不行、Dijkstra 的非负边条件、堆中旧状态如何处理。请准备一个不同边权反例。

复盘本节时先从 LeetCode 505 迷宫 II、LeetCode 743 网络延迟时间、LeetCode 778 水位上升的泳池中游泳 开始：每道题都先手写一个能覆盖全部候选的暴力版，再指出优化结构具体保存了哪一段历史信息，最后用相邻案例比较为什么不能互换解法。

\topic{并查集与动态连通性}
这一节先看 LeetCode 305 岛屿数量 II、LeetCode 399 除法求值、LeetCode 547 省份数量 这几道 LeetCode 题，再整理同一题型的分流和推导。阅读顺序不是先背原理，而是先看案例的暴力瓶颈，再把重复出现的观察抽象成方法。

\subsection{原理和适用条件}

以 LeetCode 305 岛屿数量 II、LeetCode 399 除法求值、LeetCode 547 省份数量 为主线：LeetCode 305 岛屿数量 II：模型是“陆地逐步加入，每次加入只可能影响相邻陆地所在集合。”，优化抓手是“新增陆地先让岛屿数加一，再和四邻已存在陆地合并，合并成功则岛屿数减一。”；LeetCode 399 除法求值：模型是“变量之间的除法关系可以看成带权连通分量。”，优化抓手是“并查集维护节点到根的乘积权值，合并时根据等式调整根之间权重。”；LeetCode 547 省份数量：模型是“城市连接关系形成无向图，省份就是连通分量。”，优化抓手是“并查集合并所有直接连接的城市，最后统计根数量。”。这些案例共同推出的规律是：先用搜索回答连通性，再发现查询和合并交替出现。并查集只维护集合代表，如果题目还要路径、距离、比例或奇偶性，就必须扩展状态。 方法名只是复盘后的标签，真正要掌握的是从题面约束、暴力失败、状态设计到边界反例的推导链。

\begin{principlebox}{图示：从 LeetCode 案例到 并查集与动态连通性推导链}
\begin{tabular}{@{}p{0.18\linewidth}p{0.74\linewidth}@{}}
暴力基线 & 枚举候选、完整模拟或逐个检查，先保证覆盖全部可能答案。\\
瓶颈定位 & 瓶颈是连通性查询和合并交替出现。若只需要知道是否同组，不需要路径本身，并查集正好匹配。\\
结构选择 & 路径压缩把查找路径上的节点直接挂到根，按大小或按秩合并控制树高。邮箱、等式、边、网格都可以映射成元素集合。\\
不变量 & 每个集合有一个代表元，\code{parent[root] == root}。合并只改变一个根的父亲，非根节点的父亲可能不是最终根，但 \code{find} 会返回正确代表。\\
代码落地 & 类成员访问使用 \code{this->}。迭代版 \code{find} 可以避免递归。若需要维护集合大小、边数、权值差或奇偶关系，要扩展额外数组并在合并时同步更新。\\
\end{tabular}
\end{principlebox}

\subsection{LeetCode 案例分流}

\begin{principlebox}{从 LeetCode 案例分流：并查集与动态连通性}
\textbf{静态连通块。}
代表案例：LeetCode 547 省份数量、LeetCode 721 账户合并。先看这些题的题面条件：输入给出所有连接关系，只问最终分组数量或合并结果。由此推出的处理方式是：把元素编号后合并直接相连对象，最后按根统计。风险点：字符串编号、输出分组排序和孤立元素不能漏。

\textbf{成环检测。}
代表案例：LeetCode 684 冗余连接、LeetCode 685 冗余连接 II。先看这些题的题面条件：边按顺序加入，遇到同集合两端说明形成环。由此推出的处理方式是：合并失败就是冗余边或冲突边。风险点：有向图还要处理入度为二的情况。

\textbf{动态岛屿。}
代表案例：LeetCode 305 岛屿数量 II、LeetCode 990 等式方程可满足性。先看这些题的题面条件：操作逐步添加陆地或连接关系，查询每步连通块变化。由此推出的处理方式是：新增元素初始化为集合，和相邻已存在元素合并。风险点：重复添加和无效位置要先处理。

\textbf{带权和扩展并查集。}
代表案例：LeetCode 399 除法求值、LeetCode 990 等式方程可满足性。先看这些题的题面条件：集合关系还带比例、距离、奇偶或差值。由此推出的处理方式是：合并根时同步维护相对权值，find 时压缩并更新权值。风险点：只会普通 parent 不够，必须说明附加信息语义。

\end{principlebox}

\subsection{案例矩阵}

本节案例覆盖 LeetCode 305 岛屿数量 II、LeetCode 399 除法求值、LeetCode 547 省份数量、LeetCode 684 冗余连接、LeetCode 685 冗余连接 II、LeetCode 721 账户合并、LeetCode 947 移除最多的同行或同列石头、LeetCode 959 由斜杠划分区域、LeetCode 990 等式方程的可满足性、LeetCode 1061 按字典序排列最小的等效字符串、LeetCode 1202 交换字符串中的元素、LeetCode 1319 连通网络的操作次数、LeetCode 1579 保证图可完全遍历。阅读时不要按题号背答案，而要先判断题面落在哪个分流场景：查询目标是什么，输入约束给了什么单调性或等价关系，暴力慢在哪里，优化结构保存了什么状态。

\begin{principlebox}{案例矩阵}
\textbf{LeetCode 305 岛屿数量 II。}
陆地逐步加入，每次加入只可能影响相邻陆地所在集合。单点风险：重复加入同一格、边界位置、四邻属于同一岛、空操作列表。

\textbf{LeetCode 399 除法求值。}
变量之间的除法关系可以看成带权连通分量。单点风险：查询未知变量、自除、矛盾数据、路径压缩时权值同步更新。

\textbf{LeetCode 547 省份数量。}
城市连接关系形成无向图，省份就是连通分量。单点风险：邻接矩阵对称、自连接、孤立城市、只扫描上三角。

\textbf{LeetCode 684 冗余连接。}
树加一条边会形成唯一环。单点风险：节点编号、最后一条成环边、重复边、并查集初始化大小。

\textbf{LeetCode 685 冗余连接 II。}
有向图中额外边可能导致某节点入度为二，也可能导致环。单点风险：只有环、只有入度二、两者同时存在、返回最后出现的边。

\textbf{LeetCode 721 账户合并。}
共享邮箱说明账户属于同一人，邮箱之间形成连通分量。单点风险：同名不同人、一个账户多个邮箱、重复邮箱、输出排序。

\textbf{LeetCode 947 移除最多的同行或同列石头。}
同一行或同一列的石头属于同一连通分量，每个分量至少留一块。单点风险：单块石头、行列编号冲突、多个分量、重复坐标。

\textbf{LeetCode 959 由斜杠划分区域。}
每个网格可拆成四个小三角，斜杠决定内部哪些三角相连。单点风险：空格、正斜杠、反斜杠、边界合并、区域计数。

\textbf{LeetCode 990 等式方程的可满足性。}
相等关系形成集合，不等关系要求两端不在同一集合。单点风险：自等式、自不等式、传递相等、多组变量。

\textbf{LeetCode 1061 按字典序排列最小的等效字符串。}
等价字符形成集合，每个集合应选择字典序最小代表。单点风险：传递等价、重复等价关系、自等价、未出现字符。

\textbf{LeetCode 1202 交换字符串中的元素。}
可交换下标形成连通分量，同一分量内字符可以任意重排。单点风险：没有交换、多个分量、重复字符、分量下标排序。

\textbf{LeetCode 1319 连通网络的操作次数。}
多余网线可以挪用来连接不同连通分量。单点风险：边数不足、已经连通、重复边、多分量。

\textbf{LeetCode 1579 保证图可完全遍历。}
Alice 和 Bob 各自需要图连通，公共边应优先服务两人。单点风险：公共边冗余、某一方不连通、边数很多、节点从一开始编号。

\end{principlebox}

\subsection{共性落点}

\begin{principlebox}{本组共性落点}
\textbf{慢版本。} 暴力做法每加一条边就重新 DFS 判断连通，或每次查询都扫描整个分量。这在动态合并场景下会重复处理大量已经稳定的连接关系。
\textbf{瓶颈。} 瓶颈是连通性查询和合并交替出现。若只需要知道是否同组，不需要路径本身，并查集正好匹配。
\textbf{状态字段。} parent、size 或 rank、find 返回的代表元、合并时的附加信息；带权或奇偶并查集还要维护到根的关系。
\textbf{不变量。} 每个集合有一个代表元，\code{parent[root] == root}。合并只改变一个根的父亲，非根节点的父亲可能不是最终根，但 \code{find} 会返回正确代表。
\textbf{实现检查。} find 路径压缩不能破坏附加权值；union 只能连接两个根；合并失败和重复边要保留题意解释。
\textbf{C++ 落地。} 类成员访问使用 \code{this->}。迭代版 \code{find} 可以避免递归。若需要维护集合大小、边数、权值差或奇偶关系，要扩展额外数组并在合并时同步更新。
\textbf{证明抓手。} 证明从等价关系开始：合并维护连通性的传递性，find 返回的代表元就是当前集合身份。
\end{principlebox}

\subsection{案例推导}

\noindent\textbf{LeetCode 305 岛屿数量 II。}

\textbf{题目说明。} LeetCode 305 岛屿数量 II 的题面入口要先还原成具体任务：陆地逐步加入，每次加入只可能影响相邻陆地所在集合。

\textbf{样例入口。} 拿 3x3 网格依次加 (0, 0)、(0, 1)、(1, 2) 手算，前两个陆地相邻要合并，岛屿数从 2 降回 1；第三个不相邻，岛屿数加 1。

\textbf{暴力基线。} 慢版本每次合并后用 DFS 或 BFS 重新判断连通性。它正确但重复遍历已有连接；当关系只会合并不会拆开时，可以压成集合代表。放到本题就是：先按“陆地逐步加入，每次加入只可能影响相邻陆地所在集合”完整试慢版本；样例里已经能看到“规律是每次加陆地先新增一个集合，再和四邻陆地 union，重复加点不能重复计数”，所以优化前必须先能说清慢版本枚举了哪些候选。

\textbf{暴力过程。} 拿 3x3 网格依次加 (0, 0)、(0, 1)、(1, 2) 手算，前两个陆地相邻要合并，岛屿数从 2 降回 1；第三个不相邻，岛屿数加 1。

\textbf{重复工作。} 题面模型：陆地逐步加入，每次加入只可能影响相邻陆地所在集合。暴力版的慢点要落到本题：慢版本会围绕“陆地逐步加入，每次加入只可能影响相邻陆地所在集合”反复搜索连通性。样例规律是“规律是每次加陆地先新增一个集合，再和四邻陆地 union，重复加点不能重复计数”，说明关系只需要被压成集合代表；每次 union 失败或成功都要能翻译回题意。后面的优化只消掉这类重复工作，不能改变答案集合。

\textbf{优化条件。} 本题的关键观察是：新增陆地先让岛屿数加一，再和四邻已存在陆地合并，合并成功则岛屿数减一。这句话必须能翻译成可维护状态；如果题目条件改变后观察不再成立，就不能继续套原来的写法。

\textbf{相邻状态。} 规律是每次加陆地先新增一个集合，再和四邻陆地 union，重复加点不能重复计数。把这个特例继续拆开看：union 只是在维护等价类，不是在保存具体路径。两个节点代表元相同只能说明连通或关系一致；如果题目还要距离、比例或奇偶性，必须把附加信息挂到根关系上。

\textbf{状态复用。} 把连通分量压成代表元；查询只比较代表，合并只发生在两个根之间。本题落点是：规律是每次加陆地先新增一个集合，再和四邻陆地 union，重复加点不能重复计数。手算时只改被新输入影响的那一小部分，而不是回头重算完整候选。

\textbf{指针和字段。} 状态字段要能说清：parent、size 或 rank、find 返回的代表元、合并时的附加信息；带权或奇偶并查集还要维护到根的关系。手算时按这个协议跟踪：拿 3x3 网格依次加 (0, 0)、(0, 1)、(1, 2) 手算，前两个陆地相邻要合并，岛屿数从 2 降回 1；第三个不相邻，岛屿数加 1。然后按协议复查：手算时写 parent 表和集合代表。每次 union 只改变根之间的关系，find 后代表相同才说明连通。

\textbf{代码落点。} find 路径压缩不能破坏附加权值；union 只能连接两个根；合并失败和重复边要保留题意解释。关键代码行必须能逐句翻译成“旧状态怎样变成新状态”，否则就只是模板。

\textbf{正确性。} 证明从等价关系开始：合并维护连通性的传递性，find 返回的代表元就是当前集合身份。

\textbf{边界实验。} 先用重复加入同一格、边界位置、四邻属于同一岛、空操作列表做反例检查。实验时覆盖重复加入同一格、边界位置、四邻属于同一岛、空操作列表，打印每个元素的代表元和集合大小。重复合并、已经连通的边和字符串编号最容易暴露实现问题。

\textbf{复杂度变化。} 暴力每次查询都搜索连通性代价高；并查集把合并和查询摊还到近似常数，整体接近线性。

\textbf{迁移追问。} 如果题目改成“这是动态连通性，比每次重新 BFS 更适合并查集”，先判断原状态是否仍然覆盖新答案集合，再决定是微调边界、改 value 语义，还是换算法。

\noindent\textbf{LeetCode 399 除法求值。}

\textbf{题目说明。} LeetCode 399 除法求值 的题面入口要先还原成具体任务：变量之间的除法关系可以看成带权连通分量。

\textbf{样例入口。} 拿 a/b=2、b/c=3 手算，a 到 c 的比值是沿边相乘 6；若查询 c/a，则沿反向边乘 1/3 和 1/2。

\textbf{暴力基线。} 慢版本每次合并后用 DFS 或 BFS 重新判断连通性。它正确但重复遍历已有连接；当关系只会合并不会拆开时，可以压成集合代表。放到本题就是：先按“变量之间的除法关系可以看成带权连通分量”完整试慢版本；样例里已经能看到“规律是并查集或图搜索都要维护到代表元的倍率关系，合并时同步更新权值”，所以优化前必须先能说清慢版本枚举了哪些候选。

\textbf{暴力过程。} 拿 a/b=2、b/c=3 手算，a 到 c 的比值是沿边相乘 6；若查询 c/a，则沿反向边乘 1/3 和 1/2。

\textbf{重复工作。} 题面模型：变量之间的除法关系可以看成带权连通分量。暴力版的慢点要落到本题：慢版本会围绕“变量之间的除法关系可以看成带权连通分量”反复搜索连通性。样例规律是“规律是并查集或图搜索都要维护到代表元的倍率关系，合并时同步更新权值”，说明关系只需要被压成集合代表；每次 union 失败或成功都要能翻译回题意。后面的优化只消掉这类重复工作，不能改变答案集合。

\textbf{优化条件。} 本题的关键观察是：并查集维护节点到根的乘积权值，合并时根据等式调整根之间权重。这句话必须能翻译成可维护状态；如果题目条件改变后观察不再成立，就不能继续套原来的写法。

\textbf{相邻状态。} 规律是并查集或图搜索都要维护到代表元的倍率关系，合并时同步更新权值。把这个特例继续拆开看：union 只是在维护等价类，不是在保存具体路径。两个节点代表元相同只能说明连通或关系一致；如果题目还要距离、比例或奇偶性，必须把附加信息挂到根关系上。

\textbf{状态复用。} 把连通分量压成代表元；查询只比较代表，合并只发生在两个根之间。本题落点是：规律是并查集或图搜索都要维护到代表元的倍率关系，合并时同步更新权值。手算时只改被新输入影响的那一小部分，而不是回头重算完整候选。

\textbf{指针和字段。} 状态字段要能说清：parent、size 或 rank、find 返回的代表元、合并时的附加信息；带权或奇偶并查集还要维护到根的关系。手算时按这个协议跟踪：拿 a/b=2、b/c=3 手算，a 到 c 的比值是沿边相乘 6；若查询 c/a，则沿反向边乘 1/3 和 1/2。然后按协议复查：手算时写 parent 表和集合代表。每次 union 只改变根之间的关系，find 后代表相同才说明连通。

\textbf{代码落点。} find 路径压缩不能破坏附加权值；union 只能连接两个根；合并失败和重复边要保留题意解释。关键代码行必须能逐句翻译成“旧状态怎样变成新状态”，否则就只是模板。

\textbf{正确性。} 证明从等价关系开始：合并维护连通性的传递性，find 返回的代表元就是当前集合身份。

\textbf{边界实验。} 先用查询未知变量、自除、矛盾数据、路径压缩时权值同步更新做反例检查。实验时覆盖查询未知变量、自除、矛盾数据、路径压缩时权值同步更新，打印每个元素的代表元和集合大小。重复合并、已经连通的边和字符串编号最容易暴露实现问题。

\textbf{复杂度变化。} 暴力每次查询都搜索连通性代价高；并查集把合并和查询摊还到近似常数，整体接近线性。

\textbf{迁移追问。} 如果题目改成“也可建图 DFS，但多次查询时带权并查集更适合”，先判断原状态是否仍然覆盖新答案集合，再决定是微调边界、改 value 语义，还是换算法。

\noindent\textbf{LeetCode 547 省份数量。}

\textbf{题目说明。} LeetCode 547 省份数量 的题面入口要先还原成具体任务：城市连接关系形成无向图，省份就是连通分量。

\textbf{样例入口。} 拿 isConnected 中 0 和 1 相连、2 独立手算，0 和 1 union 后属于同一代表，2 保持自己，省份数为 2。

\textbf{暴力基线。} 慢版本每次合并后用 DFS 或 BFS 重新判断连通性。它正确但重复遍历已有连接；当关系只会合并不会拆开时，可以压成集合代表。放到本题就是：先按“城市连接关系形成无向图，省份就是连通分量”完整试慢版本；样例里已经能看到“规律是矩阵中的一条连接边对应集合合并，最后代表元数量就是连通分量数”，所以优化前必须先能说清慢版本枚举了哪些候选。

\textbf{暴力过程。} 拿 isConnected 中 0 和 1 相连、2 独立手算，0 和 1 union 后属于同一代表，2 保持自己，省份数为 2。

\textbf{重复工作。} 题面模型：城市连接关系形成无向图，省份就是连通分量。暴力版的慢点要落到本题：慢版本会围绕“城市连接关系形成无向图，省份就是连通分量”反复搜索连通性。样例规律是“规律是矩阵中的一条连接边对应集合合并，最后代表元数量就是连通分量数”，说明关系只需要被压成集合代表；每次 union 失败或成功都要能翻译回题意。后面的优化只消掉这类重复工作，不能改变答案集合。

\textbf{优化条件。} 本题的关键观察是：并查集合并所有直接连接的城市，最后统计根数量。这句话必须能翻译成可维护状态；如果题目条件改变后观察不再成立，就不能继续套原来的写法。

\textbf{相邻状态。} 规律是矩阵中的一条连接边对应集合合并，最后代表元数量就是连通分量数。把这个特例继续拆开看：union 只是在维护等价类，不是在保存具体路径。两个节点代表元相同只能说明连通或关系一致；如果题目还要距离、比例或奇偶性，必须把附加信息挂到根关系上。

\textbf{状态复用。} 把连通分量压成代表元；查询只比较代表，合并只发生在两个根之间。本题落点是：规律是矩阵中的一条连接边对应集合合并，最后代表元数量就是连通分量数。手算时只改被新输入影响的那一小部分，而不是回头重算完整候选。

\textbf{指针和字段。} 状态字段要能说清：parent、size 或 rank、find 返回的代表元、合并时的附加信息；带权或奇偶并查集还要维护到根的关系。手算时按这个协议跟踪：拿 isConnected 中 0 和 1 相连、2 独立手算，0 和 1 union 后属于同一代表，2 保持自己，省份数为 2。然后按协议复查：手算时写 parent 表和集合代表。每次 union 只改变根之间的关系，find 后代表相同才说明连通。

\textbf{代码落点。} find 路径压缩不能破坏附加权值；union 只能连接两个根；合并失败和重复边要保留题意解释。关键代码行必须能逐句翻译成“旧状态怎样变成新状态”，否则就只是模板。

\textbf{正确性。} 证明从等价关系开始：合并维护连通性的传递性，find 返回的代表元就是当前集合身份。

\textbf{边界实验。} 先用邻接矩阵对称、自连接、孤立城市、只扫描上三角做反例检查。实验时覆盖邻接矩阵对称、自连接、孤立城市、只扫描上三角，打印每个元素的代表元和集合大小。重复合并、已经连通的边和字符串编号最容易暴露实现问题。

\textbf{复杂度变化。} 暴力每次查询都搜索连通性代价高；并查集把合并和查询摊还到近似常数，整体接近线性。

\textbf{迁移追问。} 如果题目改成“BFS 也可做；并查集更突出动态合并”，先判断原状态是否仍然覆盖新答案集合，再决定是微调边界、改 value 语义，还是换算法。

\noindent\textbf{LeetCode 684 冗余连接。}

\textbf{题目说明。} LeetCode 684 冗余连接 的题面入口要先还原成具体任务：树加一条边会形成唯一环。

\textbf{样例入口。} 拿 edges=[[1, 2], [1, 3], [2, 3]] 手算，前两条边把 1、2、3 连成一棵树，第三条边连接的两个点已经同组，所以它就是冗余边。

\textbf{暴力基线。} 慢版本每次合并后用 DFS 或 BFS 重新判断连通性。它正确但重复遍历已有连接；当关系只会合并不会拆开时，可以压成集合代表。放到本题就是：先按“树加一条边会形成唯一环”完整试慢版本；样例里已经能看到“规律是无向图加边时，第一次 union 失败的边形成环”，所以优化前必须先能说清慢版本枚举了哪些候选。

\textbf{暴力过程。} 拿 edges=[[1, 2], [1, 3], [2, 3]] 手算，前两条边把 1、2、3 连成一棵树，第三条边连接的两个点已经同组，所以它就是冗余边。

\textbf{重复工作。} 题面模型：树加一条边会形成唯一环。暴力版的慢点要落到本题：慢版本会围绕“树加一条边会形成唯一环”反复搜索连通性。样例规律是“规律是无向图加边时，第一次 union 失败的边形成环”，说明关系只需要被压成集合代表；每次 union 失败或成功都要能翻译回题意。后面的优化只消掉这类重复工作，不能改变答案集合。

\textbf{优化条件。} 本题的关键观察是：按顺序加边，若一条边两端已连通，则它就是冗余边。这句话必须能翻译成可维护状态；如果题目条件改变后观察不再成立，就不能继续套原来的写法。

\textbf{相邻状态。} 规律是无向图加边时，第一次 union 失败的边形成环。把这个特例继续拆开看：union 只是在维护等价类，不是在保存具体路径。两个节点代表元相同只能说明连通或关系一致；如果题目还要距离、比例或奇偶性，必须把附加信息挂到根关系上。

\textbf{状态复用。} 把连通分量压成代表元；查询只比较代表，合并只发生在两个根之间。本题落点是：规律是无向图加边时，第一次 union 失败的边形成环。手算时只改被新输入影响的那一小部分，而不是回头重算完整候选。

\textbf{指针和字段。} 状态字段要能说清：parent、size 或 rank、find 返回的代表元、合并时的附加信息；带权或奇偶并查集还要维护到根的关系。手算时按这个协议跟踪：拿 edges=[[1, 2], [1, 3], [2, 3]] 手算，前两条边把 1、2、3 连成一棵树，第三条边连接的两个点已经同组，所以它就是冗余边。然后按协议复查：手算时写 parent 表和集合代表。每次 union 只改变根之间的关系，find 后代表相同才说明连通。

\textbf{代码落点。} find 路径压缩不能破坏附加权值；union 只能连接两个根；合并失败和重复边要保留题意解释。关键代码行必须能逐句翻译成“旧状态怎样变成新状态”，否则就只是模板。

\textbf{正确性。} 证明从等价关系开始：合并维护连通性的传递性，find 返回的代表元就是当前集合身份。

\textbf{边界实验。} 先用节点编号、最后一条成环边、重复边、并查集初始化大小做反例检查。实验时覆盖节点编号、最后一条成环边、重复边、并查集初始化大小，打印每个元素的代表元和集合大小。重复合并、已经连通的边和字符串编号最容易暴露实现问题。

\textbf{复杂度变化。} 暴力每次查询都搜索连通性代价高；并查集把合并和查询摊还到近似常数，整体接近线性。

\textbf{迁移追问。} 如果题目改成“有向冗余连接需要同时处理入度为二和环，复杂度更高”，先判断原状态是否仍然覆盖新答案集合，再决定是微调边界、改 value 语义，还是换算法。

\noindent\textbf{LeetCode 685 冗余连接 II。}

\textbf{题目说明。} LeetCode 685 冗余连接 II 的题面入口要先还原成具体任务：有向图中额外边可能导致某节点入度为二，也可能导致环。

\textbf{样例入口。} 拿有向边 [[1, 2], [1, 3], [2, 3]] 手算，节点 3 有两个父亲 1 和 2；需要判断去掉哪条边后剩余图是有根树。

\textbf{暴力基线。} 慢版本每次合并后用 DFS 或 BFS 重新判断连通性。它正确但重复遍历已有连接；当关系只会合并不会拆开时，可以压成集合代表。放到本题就是：先按“有向图中额外边可能导致某节点入度为二，也可能导致环”完整试慢版本；样例里已经能看到“规律是先处理双父节点候选，再用并查集检测环，双父和成环组合时要按题意返回正确候选边”，所以优化前必须先能说清慢版本枚举了哪些候选。

\textbf{暴力过程。} 拿有向边 [[1, 2], [1, 3], [2, 3]] 手算，节点 3 有两个父亲 1 和 2；需要判断去掉哪条边后剩余图是有根树。

\textbf{重复工作。} 题面模型：有向图中额外边可能导致某节点入度为二，也可能导致环。暴力版的慢点要落到本题：慢版本会围绕“有向图中额外边可能导致某节点入度为二，也可能导致环”反复搜索连通性。样例规律是“规律是先处理双父节点候选，再用并查集检测环，双父和成环组合时要按题意返回正确候选边”，说明关系只需要被压成集合代表；每次 union 失败或成功都要能翻译回题意。后面的优化只消掉这类重复工作，不能改变答案集合。

\textbf{优化条件。} 本题的关键观察是：先记录入度为二的两条候选边，再用并查集排除其中一条测试是否成树。这句话必须能翻译成可维护状态；如果题目条件改变后观察不再成立，就不能继续套原来的写法。

\textbf{相邻状态。} 规律是先处理双父节点候选，再用并查集检测环，双父和成环组合时要按题意返回正确候选边。把这个特例继续拆开看：union 只是在维护等价类，不是在保存具体路径。两个节点代表元相同只能说明连通或关系一致；如果题目还要距离、比例或奇偶性，必须把附加信息挂到根关系上。

\textbf{状态复用。} 把连通分量压成代表元；查询只比较代表，合并只发生在两个根之间。本题落点是：规律是先处理双父节点候选，再用并查集检测环，双父和成环组合时要按题意返回正确候选边。手算时只改被新输入影响的那一小部分，而不是回头重算完整候选。

\textbf{指针和字段。} 状态字段要能说清：parent、size 或 rank、find 返回的代表元、合并时的附加信息；带权或奇偶并查集还要维护到根的关系。手算时按这个协议跟踪：拿有向边 [[1, 2], [1, 3], [2, 3]] 手算，节点 3 有两个父亲 1 和 2；需要判断去掉哪条边后剩余图是有根树。然后按协议复查：手算时写 parent 表和集合代表。每次 union 只改变根之间的关系，find 后代表相同才说明连通。

\textbf{代码落点。} find 路径压缩不能破坏附加权值；union 只能连接两个根；合并失败和重复边要保留题意解释。关键代码行必须能逐句翻译成“旧状态怎样变成新状态”，否则就只是模板。

\textbf{正确性。} 证明从等价关系开始：合并维护连通性的传递性，find 返回的代表元就是当前集合身份。

\textbf{边界实验。} 先用只有环、只有入度二、两者同时存在、返回最后出现的边做反例检查。实验时覆盖只有环、只有入度二、两者同时存在、返回最后出现的边，打印每个元素的代表元和集合大小。重复合并、已经连通的边和字符串编号最容易暴露实现问题。

\textbf{复杂度变化。} 暴力每次查询都搜索连通性代价高；并查集把合并和查询摊还到近似常数，整体接近线性。

\textbf{迁移追问。} 如果题目改成“它比无向冗余连接多了有向树入度约束”，先判断原状态是否仍然覆盖新答案集合，再决定是微调边界、改 value 语义，还是换算法。

\noindent\textbf{LeetCode 721 账户合并。}

\textbf{题目说明。} LeetCode 721 账户合并 的题面入口要先还原成具体任务：共享邮箱说明账户属于同一人，邮箱之间形成连通分量。

\textbf{样例入口。} 拿 John 的账号 [a, b] 和另一个 [b, c] 手算，邮箱 b 重叠，所以 a、b、c 属于同一人；名字只是输出标签，合并依据是邮箱。

\textbf{暴力基线。} 慢版本每次合并后用 DFS 或 BFS 重新判断连通性。它正确但重复遍历已有连接；当关系只会合并不会拆开时，可以压成集合代表。放到本题就是：先按“共享邮箱说明账户属于同一人，邮箱之间形成连通分量”完整试慢版本；样例里已经能看到“规律是邮箱映射到编号后并查集合并，最后按代表收集并排序邮箱”，所以优化前必须先能说清慢版本枚举了哪些候选。

\textbf{暴力过程。} 拿 John 的账号 [a, b] 和另一个 [b, c] 手算，邮箱 b 重叠，所以 a、b、c 属于同一人；名字只是输出标签，合并依据是邮箱。

\textbf{重复工作。} 题面模型：共享邮箱说明账户属于同一人，邮箱之间形成连通分量。暴力版的慢点要落到本题：慢版本会围绕“共享邮箱说明账户属于同一人，邮箱之间形成连通分量”反复搜索连通性。样例规律是“规律是邮箱映射到编号后并查集合并，最后按代表收集并排序邮箱”，说明关系只需要被压成集合代表；每次 union 失败或成功都要能翻译回题意。后面的优化只消掉这类重复工作，不能改变答案集合。

\textbf{优化条件。} 本题的关键观察是：给邮箱编号并查集合并，同根邮箱收集后排序输出。这句话必须能翻译成可维护状态；如果题目条件改变后观察不再成立，就不能继续套原来的写法。

\textbf{相邻状态。} 规律是邮箱映射到编号后并查集合并，最后按代表收集并排序邮箱。把这个特例继续拆开看：union 只是在维护等价类，不是在保存具体路径。两个节点代表元相同只能说明连通或关系一致；如果题目还要距离、比例或奇偶性，必须把附加信息挂到根关系上。

\textbf{状态复用。} 把连通分量压成代表元；查询只比较代表，合并只发生在两个根之间。本题落点是：规律是邮箱映射到编号后并查集合并，最后按代表收集并排序邮箱。手算时只改被新输入影响的那一小部分，而不是回头重算完整候选。

\textbf{指针和字段。} 状态字段要能说清：parent、size 或 rank、find 返回的代表元、合并时的附加信息；带权或奇偶并查集还要维护到根的关系。手算时按这个协议跟踪：拿 John 的账号 [a, b] 和另一个 [b, c] 手算，邮箱 b 重叠，所以 a、b、c 属于同一人；名字只是输出标签，合并依据是邮箱。然后按协议复查：手算时写 parent 表和集合代表。每次 union 只改变根之间的关系，find 后代表相同才说明连通。

\textbf{代码落点。} find 路径压缩不能破坏附加权值；union 只能连接两个根；合并失败和重复边要保留题意解释。关键代码行必须能逐句翻译成“旧状态怎样变成新状态”，否则就只是模板。

\textbf{正确性。} 证明从等价关系开始：合并维护连通性的传递性，find 返回的代表元就是当前集合身份。

\textbf{边界实验。} 先用同名不同人、一个账户多个邮箱、重复邮箱、输出排序做反例检查。实验时覆盖同名不同人、一个账户多个邮箱、重复邮箱、输出排序，打印每个元素的代表元和集合大小。重复合并、已经连通的边和字符串编号最容易暴露实现问题。

\textbf{复杂度变化。} 暴力每次查询都搜索连通性代价高；并查集把合并和查询摊还到近似常数，整体接近线性。

\textbf{迁移追问。} 如果题目改成“姓名不能作为合并依据，只能作为输出字段”，先判断原状态是否仍然覆盖新答案集合，再决定是微调边界、改 value 语义，还是换算法。

\noindent\textbf{LeetCode 947 移除最多的同行或同列石头。}

\textbf{题目说明。} LeetCode 947 移除最多的同行或同列石头 的题面入口要先还原成具体任务：同一行或同一列的石头属于同一连通分量，每个分量至少留一块。

\textbf{样例入口。} 拿石头 (0, 0), (0, 1), (1, 0) 手算，它们通过同一行或同一列连成一个分量，三个石头最多移除 2 个。

\textbf{暴力基线。} 慢版本每次合并后用 DFS 或 BFS 重新判断连通性。它正确但重复遍历已有连接；当关系只会合并不会拆开时，可以压成集合代表。放到本题就是：先按“同一行或同一列的石头属于同一连通分量，每个分量至少留一块”完整试慢版本；样例里已经能看到“规律是每个连通分量至少留一个，答案是石头数减分量数”，所以优化前必须先能说清慢版本枚举了哪些候选。

\textbf{暴力过程。} 拿石头 (0, 0), (0, 1), (1, 0) 手算，它们通过同一行或同一列连成一个分量，三个石头最多移除 2 个。

\textbf{重复工作。} 题面模型：同一行或同一列的石头属于同一连通分量，每个分量至少留一块。暴力版的慢点要落到本题：慢版本会围绕“同一行或同一列的石头属于同一连通分量，每个分量至少留一块”反复搜索连通性。样例规律是“规律是每个连通分量至少留一个，答案是石头数减分量数”，说明关系只需要被压成集合代表；每次 union 失败或成功都要能翻译回题意。后面的优化只消掉这类重复工作，不能改变答案集合。

\textbf{优化条件。} 本题的关键观察是：把行和列编号后合并，答案是石头数减去连通分量数。这句话必须能翻译成可维护状态；如果题目条件改变后观察不再成立，就不能继续套原来的写法。

\textbf{相邻状态。} 规律是每个连通分量至少留一个，答案是石头数减分量数。把这个特例继续拆开看：union 只是在维护等价类，不是在保存具体路径。两个节点代表元相同只能说明连通或关系一致；如果题目还要距离、比例或奇偶性，必须把附加信息挂到根关系上。

\textbf{状态复用。} 把连通分量压成代表元；查询只比较代表，合并只发生在两个根之间。本题落点是：规律是每个连通分量至少留一个，答案是石头数减分量数。手算时只改被新输入影响的那一小部分，而不是回头重算完整候选。

\textbf{指针和字段。} 状态字段要能说清：parent、size 或 rank、find 返回的代表元、合并时的附加信息；带权或奇偶并查集还要维护到根的关系。手算时按这个协议跟踪：拿石头 (0, 0), (0, 1), (1, 0) 手算，它们通过同一行或同一列连成一个分量，三个石头最多移除 2 个。然后按协议复查：手算时写 parent 表和集合代表。每次 union 只改变根之间的关系，find 后代表相同才说明连通。

\textbf{代码落点。} find 路径压缩不能破坏附加权值；union 只能连接两个根；合并失败和重复边要保留题意解释。关键代码行必须能逐句翻译成“旧状态怎样变成新状态”，否则就只是模板。

\textbf{正确性。} 证明从等价关系开始：合并维护连通性的传递性，find 返回的代表元就是当前集合身份。

\textbf{边界实验。} 先用单块石头、行列编号冲突、多个分量、重复坐标做反例检查。实验时覆盖单块石头、行列编号冲突、多个分量、重复坐标，打印每个元素的代表元和集合大小。重复合并、已经连通的边和字符串编号最容易暴露实现问题。

\textbf{复杂度变化。} 暴力每次查询都搜索连通性代价高；并查集把合并和查询摊还到近似常数，整体接近线性。

\textbf{迁移追问。} 如果题目改成“也可把石头当节点建图，但行列并查集更省边”，先判断原状态是否仍然覆盖新答案集合，再决定是微调边界、改 value 语义，还是换算法。

\noindent\textbf{LeetCode 959 由斜杠划分区域。}

\textbf{题目说明。} LeetCode 959 由斜杠划分区域 的题面入口要先还原成具体任务：每个网格可拆成四个小三角，斜杠决定内部哪些三角相连。

\textbf{样例入口。} 拿 1x1 格子里一条斜杠手算，它把小方格分成两个区域；把每格拆成 4 个三角形后，斜杠决定哪些三角形不能合并。

\textbf{暴力基线。} 慢版本每次合并后用 DFS 或 BFS 重新判断连通性。它正确但重复遍历已有连接；当关系只会合并不会拆开时，可以压成集合代表。放到本题就是：先按“每个网格可拆成四个小三角，斜杠决定内部哪些三角相连”完整试慢版本；样例里已经能看到“规律是格内按字符合并，格间相邻边合并，最终代表元数量就是区域数”，所以优化前必须先能说清慢版本枚举了哪些候选。

\textbf{暴力过程。} 拿 1x1 格子里一条斜杠手算，它把小方格分成两个区域；把每格拆成 4 个三角形后，斜杠决定哪些三角形不能合并。

\textbf{重复工作。} 题面模型：每个网格可拆成四个小三角，斜杠决定内部哪些三角相连。暴力版的慢点要落到本题：慢版本会围绕“每个网格可拆成四个小三角，斜杠决定内部哪些三角相连”反复搜索连通性。样例规律是“规律是格内按字符合并，格间相邻边合并，最终代表元数量就是区域数”，说明关系只需要被压成集合代表；每次 union 失败或成功都要能翻译回题意。后面的优化只消掉这类重复工作，不能改变答案集合。

\textbf{优化条件。} 本题的关键观察是：给每格四个区域编号，先按字符合并内部，再合并相邻格边界区域。这句话必须能翻译成可维护状态；如果题目条件改变后观察不再成立，就不能继续套原来的写法。

\textbf{相邻状态。} 规律是格内按字符合并，格间相邻边合并，最终代表元数量就是区域数。把这个特例继续拆开看：union 只是在维护等价类，不是在保存具体路径。两个节点代表元相同只能说明连通或关系一致；如果题目还要距离、比例或奇偶性，必须把附加信息挂到根关系上。

\textbf{状态复用。} 把连通分量压成代表元；查询只比较代表，合并只发生在两个根之间。本题落点是：规律是格内按字符合并，格间相邻边合并，最终代表元数量就是区域数。手算时只改被新输入影响的那一小部分，而不是回头重算完整候选。

\textbf{指针和字段。} 状态字段要能说清：parent、size 或 rank、find 返回的代表元、合并时的附加信息；带权或奇偶并查集还要维护到根的关系。手算时按这个协议跟踪：拿 1x1 格子里一条斜杠手算，它把小方格分成两个区域；把每格拆成 4 个三角形后，斜杠决定哪些三角形不能合并。然后按协议复查：手算时写 parent 表和集合代表。每次 union 只改变根之间的关系，find 后代表相同才说明连通。

\textbf{代码落点。} find 路径压缩不能破坏附加权值；union 只能连接两个根；合并失败和重复边要保留题意解释。关键代码行必须能逐句翻译成“旧状态怎样变成新状态”，否则就只是模板。

\textbf{正确性。} 证明从等价关系开始：合并维护连通性的传递性，find 返回的代表元就是当前集合身份。

\textbf{边界实验。} 先用空格、正斜杠、反斜杠、边界合并、区域计数做反例检查。实验时覆盖空格、正斜杠、反斜杠、边界合并、区域计数，打印每个元素的代表元和集合大小。重复合并、已经连通的边和字符串编号最容易暴露实现问题。

\textbf{复杂度变化。} 暴力每次查询都搜索连通性代价高；并查集把合并和查询摊还到近似常数，整体接近线性。

\textbf{迁移追问。} 如果题目改成“这是把几何连通性离散成并查集元素的典型题”，先判断原状态是否仍然覆盖新答案集合，再决定是微调边界、改 value 语义，还是换算法。

\noindent\textbf{LeetCode 990 等式方程的可满足性。}

\textbf{题目说明。} LeetCode 990 等式方程的可满足性 的题面入口要先还原成具体任务：相等关系形成集合，不等关系要求两端不在同一集合。

\textbf{样例入口。} 拿 a==b、b==c、a!=c 手算，前两条等式把 a, b, c 合并，最后不等式发现二者同组，矛盾。

\textbf{暴力基线。} 慢版本每次合并后用 DFS 或 BFS 重新判断连通性。它正确但重复遍历已有连接；当关系只会合并不会拆开时，可以压成集合代表。放到本题就是：先按“相等关系形成集合，不等关系要求两端不在同一集合”完整试慢版本；样例里已经能看到“规律是先处理所有等式建立集合，再检查不等式；顺序反过来会漏掉后续合并造成的冲突”，所以优化前必须先能说清慢版本枚举了哪些候选。

\textbf{暴力过程。} 拿 a==b、b==c、a!=c 手算，前两条等式把 a, b, c 合并，最后不等式发现二者同组，矛盾。

\textbf{重复工作。} 题面模型：相等关系形成集合，不等关系要求两端不在同一集合。暴力版的慢点要落到本题：慢版本会围绕“相等关系形成集合，不等关系要求两端不在同一集合”反复搜索连通性。样例规律是“规律是先处理所有等式建立集合，再检查不等式；顺序反过来会漏掉后续合并造成的冲突”，说明关系只需要被压成集合代表；每次 union 失败或成功都要能翻译回题意。后面的优化只消掉这类重复工作，不能改变答案集合。

\textbf{优化条件。} 本题的关键观察是：先合并所有等式，再检查每个不等式是否冲突。这句话必须能翻译成可维护状态；如果题目条件改变后观察不再成立，就不能继续套原来的写法。

\textbf{相邻状态。} 规律是先处理所有等式建立集合，再检查不等式；顺序反过来会漏掉后续合并造成的冲突。把这个特例继续拆开看：union 只是在维护等价类，不是在保存具体路径。两个节点代表元相同只能说明连通或关系一致；如果题目还要距离、比例或奇偶性，必须把附加信息挂到根关系上。

\textbf{状态复用。} 把连通分量压成代表元；查询只比较代表，合并只发生在两个根之间。本题落点是：规律是先处理所有等式建立集合，再检查不等式；顺序反过来会漏掉后续合并造成的冲突。手算时只改被新输入影响的那一小部分，而不是回头重算完整候选。

\textbf{指针和字段。} 状态字段要能说清：parent、size 或 rank、find 返回的代表元、合并时的附加信息；带权或奇偶并查集还要维护到根的关系。手算时按这个协议跟踪：拿 a==b、b==c、a!=c 手算，前两条等式把 a, b, c 合并，最后不等式发现二者同组，矛盾。然后按协议复查：手算时写 parent 表和集合代表。每次 union 只改变根之间的关系，find 后代表相同才说明连通。

\textbf{代码落点。} find 路径压缩不能破坏附加权值；union 只能连接两个根；合并失败和重复边要保留题意解释。关键代码行必须能逐句翻译成“旧状态怎样变成新状态”，否则就只是模板。

\textbf{正确性。} 证明从等价关系开始：合并维护连通性的传递性，find 返回的代表元就是当前集合身份。

\textbf{边界实验。} 先用自等式、自不等式、传递相等、多组变量做反例检查。实验时覆盖自等式、自不等式、传递相等、多组变量，打印每个元素的代表元和集合大小。重复合并、已经连通的边和字符串编号最容易暴露实现问题。

\textbf{复杂度变化。} 暴力每次查询都搜索连通性代价高；并查集把合并和查询摊还到近似常数，整体接近线性。

\textbf{迁移追问。} 如果题目改成“先处理不等式会漏掉后续等式带来的传递冲突”，先判断原状态是否仍然覆盖新答案集合，再决定是微调边界、改 value 语义，还是换算法。

\noindent\textbf{LeetCode 1061 按字典序排列最小的等效字符串。}

\textbf{题目说明。} LeetCode 1061 按字典序排列最小的等效字符串 的题面入口要先还原成具体任务：等价字符形成集合，每个集合应选择字典序最小代表。

\textbf{样例入口。} 拿 s1=parker、s2=morris、base=parser 手算，p 和 m 等价，a 和 o 等价；每个等价类输出字典序最小字符。

\textbf{暴力基线。} 慢版本每次合并后用 DFS 或 BFS 重新判断连通性。它正确但重复遍历已有连接；当关系只会合并不会拆开时，可以压成集合代表。放到本题就是：先按“等价字符形成集合，每个集合应选择字典序最小代表”完整试慢版本；样例里已经能看到“规律是并查集合并字符时让较小根做代表，转换 base 时查代表”，所以优化前必须先能说清慢版本枚举了哪些候选。

\textbf{暴力过程。} 拿 s1=parker、s2=morris、base=parser 手算，p 和 m 等价，a 和 o 等价；每个等价类输出字典序最小字符。

\textbf{重复工作。} 题面模型：等价字符形成集合，每个集合应选择字典序最小代表。暴力版的慢点要落到本题：慢版本会围绕“等价字符形成集合，每个集合应选择字典序最小代表”反复搜索连通性。样例规律是“规律是并查集合并字符时让较小根做代表，转换 base 时查代表”，说明关系只需要被压成集合代表；每次 union 失败或成功都要能翻译回题意。后面的优化只消掉这类重复工作，不能改变答案集合。

\textbf{优化条件。} 本题的关键观察是：合并两个字符集合时让较小根成为代表，扫描目标串时替换为根。这句话必须能翻译成可维护状态；如果题目条件改变后观察不再成立，就不能继续套原来的写法。

\textbf{相邻状态。} 规律是并查集合并字符时让较小根做代表，转换 base 时查代表。把这个特例继续拆开看：union 只是在维护等价类，不是在保存具体路径。两个节点代表元相同只能说明连通或关系一致；如果题目还要距离、比例或奇偶性，必须把附加信息挂到根关系上。

\textbf{状态复用。} 把连通分量压成代表元；查询只比较代表，合并只发生在两个根之间。本题落点是：规律是并查集合并字符时让较小根做代表，转换 base 时查代表。手算时只改被新输入影响的那一小部分，而不是回头重算完整候选。

\textbf{指针和字段。} 状态字段要能说清：parent、size 或 rank、find 返回的代表元、合并时的附加信息；带权或奇偶并查集还要维护到根的关系。手算时按这个协议跟踪：拿 s1=parker、s2=morris、base=parser 手算，p 和 m 等价，a 和 o 等价；每个等价类输出字典序最小字符。然后按协议复查：手算时写 parent 表和集合代表。每次 union 只改变根之间的关系，find 后代表相同才说明连通。

\textbf{代码落点。} find 路径压缩不能破坏附加权值；union 只能连接两个根；合并失败和重复边要保留题意解释。关键代码行必须能逐句翻译成“旧状态怎样变成新状态”，否则就只是模板。

\textbf{正确性。} 证明从等价关系开始：合并维护连通性的传递性，find 返回的代表元就是当前集合身份。

\textbf{边界实验。} 先用传递等价、重复等价关系、自等价、未出现字符做反例检查。实验时覆盖传递等价、重复等价关系、自等价、未出现字符，打印每个元素的代表元和集合大小。重复合并、已经连通的边和字符串编号最容易暴露实现问题。

\textbf{复杂度变化。} 暴力每次查询都搜索连通性代价高；并查集把合并和查询摊还到近似常数，整体接近线性。

\textbf{迁移追问。} 如果题目改成“如果集合代表还要保存其他权值，合并规则需要同步维护”，先判断原状态是否仍然覆盖新答案集合，再决定是微调边界、改 value 语义，还是换算法。

\noindent\textbf{LeetCode 1202 交换字符串中的元素。}

\textbf{题目说明。} LeetCode 1202 交换字符串中的元素 的题面入口要先还原成具体任务：可交换下标形成连通分量，同一分量内字符可以任意重排。

\textbf{样例入口。} 拿 s=dcab、pairs=[[0, 3], [1, 2]] 手算，0 和 3 可交换成一组，1 和 2 可交换成一组；每组内部字符排序后放回最小位置。

\textbf{暴力基线。} 慢版本每次合并后用 DFS 或 BFS 重新判断连通性。它正确但重复遍历已有连接；当关系只会合并不会拆开时，可以压成集合代表。放到本题就是：先按“可交换下标形成连通分量，同一分量内字符可以任意重排”完整试慢版本；样例里已经能看到“规律是并查集找可交换连通分量，每个分量的字符可任意重排”，所以优化前必须先能说清慢版本枚举了哪些候选。

\textbf{暴力过程。} 拿 s=dcab、pairs=[[0, 3], [1, 2]] 手算，0 和 3 可交换成一组，1 和 2 可交换成一组；每组内部字符排序后放回最小位置。

\textbf{重复工作。} 题面模型：可交换下标形成连通分量，同一分量内字符可以任意重排。暴力版的慢点要落到本题：慢版本会围绕“可交换下标形成连通分量，同一分量内字符可以任意重排”反复搜索连通性。样例规律是“规律是并查集找可交换连通分量，每个分量的字符可任意重排”，说明关系只需要被压成集合代表；每次 union 失败或成功都要能翻译回题意。后面的优化只消掉这类重复工作，不能改变答案集合。

\textbf{优化条件。} 本题的关键观察是：并查集合并所有可交换下标，对每个分量收集字符并按字典序分配回最小下标。这句话必须能翻译成可维护状态；如果题目条件改变后观察不再成立，就不能继续套原来的写法。

\textbf{相邻状态。} 规律是并查集找可交换连通分量，每个分量的字符可任意重排。把这个特例继续拆开看：union 只是在维护等价类，不是在保存具体路径。两个节点代表元相同只能说明连通或关系一致；如果题目还要距离、比例或奇偶性，必须把附加信息挂到根关系上。

\textbf{状态复用。} 把连通分量压成代表元；查询只比较代表，合并只发生在两个根之间。本题落点是：规律是并查集找可交换连通分量，每个分量的字符可任意重排。手算时只改被新输入影响的那一小部分，而不是回头重算完整候选。

\textbf{指针和字段。} 状态字段要能说清：parent、size 或 rank、find 返回的代表元、合并时的附加信息；带权或奇偶并查集还要维护到根的关系。手算时按这个协议跟踪：拿 s=dcab、pairs=[[0, 3], [1, 2]] 手算，0 和 3 可交换成一组，1 和 2 可交换成一组；每组内部字符排序后放回最小位置。然后按协议复查：手算时写 parent 表和集合代表。每次 union 只改变根之间的关系，find 后代表相同才说明连通。

\textbf{代码落点。} find 路径压缩不能破坏附加权值；union 只能连接两个根；合并失败和重复边要保留题意解释。关键代码行必须能逐句翻译成“旧状态怎样变成新状态”，否则就只是模板。

\textbf{正确性。} 证明从等价关系开始：合并维护连通性的传递性，find 返回的代表元就是当前集合身份。

\textbf{边界实验。} 先用没有交换、多个分量、重复字符、分量下标排序做反例检查。实验时覆盖没有交换、多个分量、重复字符、分量下标排序，打印每个元素的代表元和集合大小。重复合并、已经连通的边和字符串编号最容易暴露实现问题。

\textbf{复杂度变化。} 暴力每次查询都搜索连通性代价高；并查集把合并和查询摊还到近似常数，整体接近线性。

\textbf{迁移追问。} 如果题目改成“它展示并查集不仅能回答连通性，还能把分量内资源重新组织”，先判断原状态是否仍然覆盖新答案集合，再决定是微调边界、改 value 语义，还是换算法。

\noindent\textbf{LeetCode 1319 连通网络的操作次数。}

\textbf{题目说明。} LeetCode 1319 连通网络的操作次数 的题面入口要先还原成具体任务：多余网线可以挪用来连接不同连通分量。

\textbf{样例入口。} 拿 n=4、连接 [[0, 1], [0, 2], [1, 2]] 手算，前三个点已有一条冗余边，但节点 3 孤立；总边数只有 3，刚好可以拿冗余边连上 3。

\textbf{暴力基线。} 慢版本每次合并后用 DFS 或 BFS 重新判断连通性。它正确但重复遍历已有连接；当关系只会合并不会拆开时，可以压成集合代表。放到本题就是：先按“多余网线可以挪用来连接不同连通分量”完整试慢版本；样例里已经能看到“规律是边数少于 n-1 直接失败，否则答案是连通分量数减一”，所以优化前必须先能说清慢版本枚举了哪些候选。

\textbf{暴力过程。} 拿 n=4、连接 [[0, 1], [0, 2], [1, 2]] 手算，前三个点已有一条冗余边，但节点 3 孤立；总边数只有 3，刚好可以拿冗余边连上 3。

\textbf{重复工作。} 题面模型：多余网线可以挪用来连接不同连通分量。暴力版的慢点要落到本题：慢版本会围绕“多余网线可以挪用来连接不同连通分量”反复搜索连通性。样例规律是“规律是边数少于 n-1 直接失败，否则答案是连通分量数减一”，说明关系只需要被压成集合代表；每次 union 失败或成功都要能翻译回题意。后面的优化只消掉这类重复工作，不能改变答案集合。

\textbf{优化条件。} 本题的关键观察是：先判断边数至少为 n 减一，再用并查集统计连通分量，答案是分量数减一。这句话必须能翻译成可维护状态；如果题目条件改变后观察不再成立，就不能继续套原来的写法。

\textbf{相邻状态。} 规律是边数少于 n-1 直接失败，否则答案是连通分量数减一。把这个特例继续拆开看：union 只是在维护等价类，不是在保存具体路径。两个节点代表元相同只能说明连通或关系一致；如果题目还要距离、比例或奇偶性，必须把附加信息挂到根关系上。

\textbf{状态复用。} 把连通分量压成代表元；查询只比较代表，合并只发生在两个根之间。本题落点是：规律是边数少于 n-1 直接失败，否则答案是连通分量数减一。手算时只改被新输入影响的那一小部分，而不是回头重算完整候选。

\textbf{指针和字段。} 状态字段要能说清：parent、size 或 rank、find 返回的代表元、合并时的附加信息；带权或奇偶并查集还要维护到根的关系。手算时按这个协议跟踪：拿 n=4、连接 [[0, 1], [0, 2], [1, 2]] 手算，前三个点已有一条冗余边，但节点 3 孤立；总边数只有 3，刚好可以拿冗余边连上 3。然后按协议复查：手算时写 parent 表和集合代表。每次 union 只改变根之间的关系，find 后代表相同才说明连通。

\textbf{代码落点。} find 路径压缩不能破坏附加权值；union 只能连接两个根；合并失败和重复边要保留题意解释。关键代码行必须能逐句翻译成“旧状态怎样变成新状态”，否则就只是模板。

\textbf{正确性。} 证明从等价关系开始：合并维护连通性的传递性，find 返回的代表元就是当前集合身份。

\textbf{边界实验。} 先用边数不足、已经连通、重复边、多分量做反例检查。实验时覆盖边数不足、已经连通、重复边、多分量，打印每个元素的代表元和集合大小。重复合并、已经连通的边和字符串编号最容易暴露实现问题。

\textbf{复杂度变化。} 暴力每次查询都搜索连通性代价高；并查集把合并和查询摊还到近似常数，整体接近线性。

\textbf{迁移追问。} 如果题目改成“合并失败的边代表冗余资源，但只要总边数足够即可”，先判断原状态是否仍然覆盖新答案集合，再决定是微调边界、改 value 语义，还是换算法。

\noindent\textbf{LeetCode 1579 保证图可完全遍历。}

\textbf{题目说明。} LeetCode 1579 保证图可完全遍历 的题面入口要先还原成具体任务：Alice 和 Bob 各自需要图连通，公共边应优先服务两人。

\textbf{样例入口。} 拿公共边能同时连接 Alice 和 Bob 的两个分量手算，先用类型 3 边可以让两个人都受益；若先处理私有边，公共边可能变成冗余。

\textbf{暴力基线。} 慢版本每次合并后用 DFS 或 BFS 重新判断连通性。它正确但重复遍历已有连接；当关系只会合并不会拆开时，可以压成集合代表。放到本题就是：先按“Alice 和 Bob 各自需要图连通，公共边应优先服务两人”完整试慢版本；样例里已经能看到“规律是先处理公共边，再分别处理 Alice 和 Bob 的边，最后两套并查集都必须连通”，所以优化前必须先能说清慢版本枚举了哪些候选。

\textbf{暴力过程。} 拿公共边能同时连接 Alice 和 Bob 的两个分量手算，先用类型 3 边可以让两个人都受益；若先处理私有边，公共边可能变成冗余。

\textbf{重复工作。} 题面模型：Alice 和 Bob 各自需要图连通，公共边应优先服务两人。暴力版的慢点要落到本题：慢版本会围绕“Alice 和 Bob 各自需要图连通，公共边应优先服务两人”反复搜索连通性。样例规律是“规律是先处理公共边，再分别处理 Alice 和 Bob 的边，最后两套并查集都必须连通”，说明关系只需要被压成集合代表；每次 union 失败或成功都要能翻译回题意。后面的优化只消掉这类重复工作，不能改变答案集合。

\textbf{优化条件。} 本题的关键观察是：先处理类型三公共边同步合并两个并查集，再分别处理各自边，最后检查两图连通。这句话必须能翻译成可维护状态；如果题目条件改变后观察不再成立，就不能继续套原来的写法。

\textbf{相邻状态。} 规律是先处理公共边，再分别处理 Alice 和 Bob 的边，最后两套并查集都必须连通。把这个特例继续拆开看：union 只是在维护等价类，不是在保存具体路径。两个节点代表元相同只能说明连通或关系一致；如果题目还要距离、比例或奇偶性，必须把附加信息挂到根关系上。

\textbf{状态复用。} 把连通分量压成代表元；查询只比较代表，合并只发生在两个根之间。本题落点是：规律是先处理公共边，再分别处理 Alice 和 Bob 的边，最后两套并查集都必须连通。手算时只改被新输入影响的那一小部分，而不是回头重算完整候选。

\textbf{指针和字段。} 状态字段要能说清：parent、size 或 rank、find 返回的代表元、合并时的附加信息；带权或奇偶并查集还要维护到根的关系。手算时按这个协议跟踪：拿公共边能同时连接 Alice 和 Bob 的两个分量手算，先用类型 3 边可以让两个人都受益；若先处理私有边，公共边可能变成冗余。然后按协议复查：手算时写 parent 表和集合代表。每次 union 只改变根之间的关系，find 后代表相同才说明连通。

\textbf{代码落点。} find 路径压缩不能破坏附加权值；union 只能连接两个根；合并失败和重复边要保留题意解释。关键代码行必须能逐句翻译成“旧状态怎样变成新状态”，否则就只是模板。

\textbf{正确性。} 证明从等价关系开始：合并维护连通性的传递性，find 返回的代表元就是当前集合身份。

\textbf{边界实验。} 先用公共边冗余、某一方不连通、边数很多、节点从一开始编号做反例检查。实验时覆盖公共边冗余、某一方不连通、边数很多、节点从一开始编号，打印每个元素的代表元和集合大小。重复合并、已经连通的边和字符串编号最容易暴露实现问题。

\textbf{复杂度变化。} 暴力每次查询都搜索连通性代价高；并查集把合并和查询摊还到近似常数，整体接近线性。

\textbf{迁移追问。} 如果题目改成“这是双并查集和资源优先级的组合”，先判断原状态是否仍然覆盖新答案集合，再决定是微调边界、改 value 语义，还是换算法。

\subsection{实验复盘}

追问通常会问并查集能否回答路径长度、路径压缩是否改变集合语义、字符串如何编号。请准备说明合并失败为什么意味着环。

复盘本节时先从 LeetCode 305 岛屿数量 II、LeetCode 399 除法求值、LeetCode 547 省份数量 开始：每道题都先手写一个能覆盖全部候选的暴力版，再指出优化结构具体保存了哪一段历史信息，最后用相邻案例比较为什么不能互换解法。
